\input{preamble}

% OK, start here.
%
\begin{document}

\title{Algèbre homologique}


\maketitle

\phantomsection
\label{section-phantom}

\tableofcontents

\section{Introduction}
\label{section-introduction}

\noindent
L'algèbre homologique de base sera exposée dans ce document.
Nous compléterons l'exposé selon les besoins des autres parties, puisqu'il existe
manifestement une quantité infinie de résultats sur le sujet.
On pourra consulter \cite{Maclane}.

\section{Notions de base}
\label{section-topology-basic}

\noindent
Les notions suivantes sont considérées comme élémentaires et ne seront
ni définies ni démontrées. Cela ne signifie pas qu'elles soient toutes
nécessairement faciles ou bien connues.

\begin{enumerate}
\item Rien pour l'instant.
\end{enumerate}




\section{Catégories préadditives et additives}
\label{section-additive-categories}

\noindent
Voici la définition d'une catégorie préadditive.

\begin{definition}
\label{definition-preadditive}
Une catégorie $\mathcal{A}$ est dite {\it préadditive} si chaque
ensemble de morphismes $\Mor_\mathcal{A}(x, y)$ est muni
d'une structure de groupe abélien telle que les
compositions
$$
\Mor(x, y) \times \Mor(y, z)
\longrightarrow
\Mor(x, z)
$$
soient bilinéaires. Un foncteur $F : \mathcal{A} \to \mathcal{B}$ entre
catégories préadditives est dit {\it additif} si et seulement si
$F : \Mor(x, y) \to \Mor(F(x), F(y))$
est un homomorphisme de groupes abéliens pour tous
$x, y \in \Ob(\mathcal{A})$.
\end{definition}

\noindent
En particulier, pour tous $x, y$, il existe au moins
un morphisme $x \to y$, à savoir le morphisme nul.

\begin{lemma}
\label{lemma-preadditive-zero}
Soit $\mathcal{A}$ une catégorie préadditive.
Soit $x$ un objet de $\mathcal{A}$.
Les assertions suivantes sont équivalentes :
\begin{enumerate}
\item $x$ est un objet initial,
\item $x$ est un objet final, et
\item $\text{id}_x = 0$ dans $\Mor_\mathcal{A}(x, x)$.
\end{enumerate}
De plus, si un tel objet $0$ existe, alors un morphisme
$\alpha : y \to z$ se factorise par $0$ si et seulement si $\alpha = 0$.
\end{lemma}

\begin{proof}
Supposons d'abord que $x$ soit soit (1) initial, soit (2) final.
Dans les deux cas, $\Mor(x,x)$ est un groupe abélien trivial
contenant $\text{id}_x$ ; ainsi $\text{id}_x = 0$ dans
$\Mor(x, x)$, ce qui montre que (1) et (2) impliquent chacun (3).

\medskip\noindent
Supposons maintenant que $\text{id}_x = 0$ dans $\Mor(x,x)$. Soit $y$
un objet arbitraire de $\mathcal{A}$ et soit $f \in \Mor(x ,y)$.
Notons $C : \Mor(x,x) \times \Mor(x,y) \to \Mor(x,y)$ l'application de composition.
Alors $f = C(0, f)$ et, puisque $C$ est bilinéaire, on a $C(0, f) = 0$.
Ainsi $f = 0$. Par conséquent, $x$ est initial dans $\mathcal{A}$.
Un argument semblable appliqué à $f \in \Mor(y, x)$ montre que
$x$ est aussi final. Ainsi (3) implique à la fois (1) et (2).
\end{proof}

\begin{definition}
\label{definition-zero-object}
Dans une catégorie préadditive $\mathcal{A}$, on appelle un objet à la fois
final et initial (comme dans le lemme \ref{lemma-preadditive-zero}) un
{\it objet nul}, et on le note $0$.
\end{definition}

\begin{lemma}
\label{lemma-preadditive-direct-sum}
Soit $\mathcal{A}$ une catégorie préadditive.
Soient $x, y \in \Ob(\mathcal{A})$.
Si le produit $x \times y$ existe, alors le
coproduit $x \amalg y$ existe aussi.
Si le coproduit $x \amalg y$ existe, alors le
produit $x \times y$ existe aussi. Dans ce cas,
on a également $x \amalg y \cong x \times y$.
\end{lemma}

\begin{proof}
Supposons que $z = x \times y$, avec projections
$p : z \to x$ et $q : z \to y$. Notons $i : x \to z$
le morphisme correspondant à $(1, 0)$. Notons $j : y \to z$
le morphisme correspondant à $(0, 1)$. On a donc le
diagramme commutatif
$$
\xymatrix{
x \ar[rr]^1 \ar[rd]^i & & x \\
& z \ar[ru]^p \ar[rd]^q & \\
y \ar[rr]^1 \ar[ru]^j & & y
}
$$
où les compositions diagonales sont nulles. Il s'ensuit que
$i \circ p + j \circ q : z \to z$ est l'identité, car
c'est un morphisme dont la composition avec $p$ donne $p$
et dont la composition avec $q$ donne $q$.
Supposons donnés des morphismes $a : x \to w$ et $b : y \to w$.
On peut alors former l'application $a \circ p + b \circ q : z \to w$.
On obtient ainsi une bijection $\Mor(z, w)
= \Mor(x, w) \times \Mor(y, w)$, ce qui
montre que $z = x \amalg y$.

\medskip\noindent
Nous laissons au lecteur le soin de construire les morphismes
$p, q$ lorsqu'on part d'un coproduit $x \amalg y$ plutôt que d'un
produit.
\end{proof}

\begin{definition}
\label{definition-direct-sum}
Étant donnés deux objets $x, y$ d'une catégorie préadditive $\mathcal{A}$,
la {\it somme directe} $x \oplus y$ de $x$ et $y$ est le
produit direct $x \times y$ muni des morphismes
$i, j, p, q$ du lemme \ref{lemma-preadditive-direct-sum} ci-dessus.
\end{definition}

\begin{remark}
\label{remark-direct-sum}
Remarquons que la preuve du lemme \ref{lemma-preadditive-direct-sum}
montre que, si $p$ et $q$ sont donnés, les morphismes $i$, $j$ sont uniquement
déterminés par les relations $p \circ i = \text{id}_x$,
$q \circ j = \text{id}_y$, $p \circ j = 0$, $q \circ i = 0$.
De plus, on a automatiquement
$i \circ p + j \circ q = \text{id}_{x \oplus y}$.
De même, si $i$, $j$ sont donnés, les morphismes $p$ et $q$ sont uniquement déterminés.
Enfin, étant donnés des objets $x, y, z$ et des morphismes
$i : x \to z$, $j : y \to z$, $p : z \to x$ et
$q : z \to y$ tels que $p \circ i = \text{id}_x$,
$q \circ j = \text{id}_y$, $p \circ j = 0$, $q \circ i = 0$
et $i \circ p + j \circ q = \text{id}_z$, alors $z$
est la somme directe de $x$ et $y$, les quatre morphismes
étant $i, j, p, q$.
\end{remark}

\begin{lemma}
\label{lemma-additive-additive}
Soient $\mathcal{A}$, $\mathcal{B}$ des catégories préadditives.
Soit $F : \mathcal{A} \to \mathcal{B}$ un foncteur additif.
Alors $F$ transforme les sommes directes en sommes directes et l'objet nul en objet nul.
\end{lemma}

\begin{proof}
Supposons $F$ additif. Une somme directe $z$
de $x$ et $y$ est caractérisée par l'existence de morphismes
$i : x \to z$, $j : y \to z$, $p : z \to x$ et
$q : z \to y$ tels que $p \circ i = \text{id}_x$,
$q \circ j = \text{id}_y$, $p \circ j = 0$, $q \circ i = 0$
et $i \circ p + j \circ q = \text{id}_z$, d'après la
remarque \ref{remark-direct-sum}. Il est clair que $F(x), F(y), F(z)$
et les morphismes $F(i), F(j), F(p), F(q)$ satisfont exactement les
mêmes relations (par additivité), et l'on voit que $F(z)$ est
une somme directe de $F(x)$ et $F(y)$.
Ainsi, $F$ transforme les sommes directes en sommes directes.

\medskip\noindent
Pour voir que $F$ transforme l'objet nul en objet nul, on utilise la
caractérisation (3) de l'objet nul donnée dans le
lemme \ref{lemma-preadditive-zero}.
\end{proof}

\begin{definition}
\label{definition-additive-category}
Une catégorie $\mathcal{A}$ est dite {\it additive}
si elle est préadditive et si les produits finis existent ; autrement
dit, elle possède un objet nul et des sommes directes.
\end{definition}

\noindent
En effet, le produit vide est un produit fini et,
s'il existe, c'est un objet final.

\begin{definition}
\label{definition-kernel}
Soit $\mathcal{A}$ une catégorie préadditive.
Soit $f : x \to y$ un morphisme.
\begin{enumerate}
\item Un {\it noyau} de $f$ est un morphisme
$i : z \to x$ tel que (a) $f \circ i = 0$ et (b),
pour tout $i' : z' \to x$ tel que $f \circ i' = 0$, il
existe un unique morphisme $g : z' \to z$ tel que
$i' = i \circ g$.
\item Si le noyau de $f$ existe, on le note
$\Ker(f) \to x$.
\item Un {\it conoyau} de $f$ est un morphisme
$p : y \to z$ tel que (a) $p \circ f = 0$ et (b),
pour tout $p' : y \to z'$ tel que $p' \circ f = 0$, il
existe un unique morphisme $g : z \to z'$ tel que
$p' = g \circ p$.
\item Si un conoyau de $f$ existe, on le note
$y \to \Coker(f)$.
\item Si un noyau de $f$ existe, alors une {\it coimage
de $f$} est un conoyau du morphisme $\Ker(f) \to x$.
\item Si un noyau et une coimage existent, on note cette dernière
$x \to \Coim(f)$.
\item Si un conoyau de $f$ existe, alors l'{\it image de
$f$} est un noyau du morphisme $y \to \Coker(f)$.
\item Si un conoyau et une image de $f$ existent, on note cette dernière
$\Im(f) \to y$.
\end{enumerate}
\end{definition}

\noindent
Dans la définition précédente, nous avons parlé du « noyau » et du
« conoyau », en utilisant tacitement leur unicité
à isomorphisme unique près. Celle-ci découle du lemme de Yoneda
(Catégories, section \ref{categories-section-opposite}),
car le noyau de $f : x \to y$ représente le foncteur qui
envoie un objet $z$ sur l'ensemble
$\Ker(\Mor_\mathcal{A}(z, x) \to \Mor_\mathcal{A}(z, y))$.
Le cas des conoyaux s'en déduit par dualité.

\medskip\noindent
Commençons par relier la somme directe aux noyaux comme suit.

\begin{lemma}
\label{lemma-additive-cat-biproduct-kernel}
Soit $\mathcal{C}$ une catégorie préadditive.
Soit $x \oplus y$, muni des morphismes $i, j, p, q$ du
lemme \ref{lemma-preadditive-direct-sum},
une somme directe dans $\mathcal{C}$. Alors $i : x \to x \oplus y$
est un noyau de $q : x \oplus y \rightarrow y$. Par dualité, $p$ est
un conoyau de $j$.
\end{lemma}

\begin{proof}
Soit $f : z' \to x \oplus y$ un morphisme tel que $q \circ f = 0$.
Nous devons montrer qu'il existe un unique morphisme $g : z' \to x$
tel que $f = i \circ g$. Puisque $i \circ p + j \circ q$ est l'identité de
$x \oplus y$, on voit que
$$
f = (i \circ p + j \circ q) \circ f = i \circ p \circ f
$$
et donc $g = p \circ f$ convient. L'unicité vient de ce que $p \circ i$
est l'identité de $x$. La preuve de la seconde assertion est duale.
\end{proof}

\begin{lemma}
\label{lemma-kernel-mono}
Soit $\mathcal{C}$ une catégorie préadditive.
Soit $f : x \to y$ un morphisme de $\mathcal{C}$.
\begin{enumerate}
\item Si un noyau de $f$ existe, alors
ce noyau est un monomorphisme.
\item Si un conoyau de $f$ existe, alors
ce conoyau est un épimorphisme.
\item Si un noyau et une coimage de $f$ existent, alors
la coimage est un épimorphisme.
\item Si un conoyau et une image de $f$ existent, alors
l'image est un monomorphisme.
\end{enumerate}
\end{lemma}

\begin{proof}
La partie (1) résulte aisément de l'unicité exigée dans la
définition d'un noyau. La preuve de (2) est duale.
La partie (3) découle de (2), puisque la coimage est un conoyau.
De même, (4) découle de (1).
\end{proof}

\begin{lemma}
\label{lemma-coim-im-map}
Soit $f : x \to y$ un morphisme dans une catégorie préadditive
tel que le noyau, le conoyau, l'image et la coimage existent tous.
Alors $f$ se factorise de manière unique sous la forme
$x \to \Coim(f) \to \Im(f) \to y$.
\end{lemma}

\begin{proof}
Il existe un morphisme canonique $\Coim(f) \to y$
car $\Ker(f) \to x \to y$ est nul.
La composition $\Coim(f) \to y \to \Coker(f)$
est nulle, car c'est l'unique morphisme qui donne
naissance au morphisme $x \to y \to \Coker(f)$, lequel
est nul
(l'unicité découle du
lemme \ref{lemma-kernel-mono} (3)).
Ainsi, $\Coim(f) \to y$ se factorise de manière unique par
$\Im(f) \to y$, ce qui fournit l'application voulue.
\end{proof}

\begin{example}
\label{example-not-abelian}
Soit $k$ un corps.
Considérons la catégorie
des espaces vectoriels filtrés sur $k$.
(Voir la définition \ref{definition-filtered}.)
Considérons les espaces vectoriels filtrés $(V, F)$ et $(W, F)$ avec
$V = W = k$ et
$$
F^iV
=
\left\{
\begin{matrix}
V & \text{si} & i < 0 \\
0 & \text{si} & i \geq 0
\end{matrix}
\right.
\text{ et }
F^iW
=
\left\{
\begin{matrix}
W & \text{si} & i \leq 0 \\
0 & \text{si} & i > 0
\end{matrix}
\right.
$$
L'application $f : V \to W$ correspondant à $\text{id}_k$ sur les espaces
vectoriels sous-jacents a un noyau et un conoyau triviaux, mais n'est pas
un isomorphisme. Remarquons aussi que $\Coim(f) = V$ et $\Im(f) = W$.
Cela signifie que la catégorie des espaces vectoriels filtrés sur $k$
n'est pas abélienne.
\end{example}

\noindent
Nous terminons cette section par quelques lemmes sur le scindement des images
d'endomorphismes idempotents. Rappelons qu'un endomorphisme $f$
est idempotent si $f \circ f = f$.

\begin{remark}
\label{remark-idempotent-symmetry}
Soient $\mathcal{A}$ une catégorie préadditive,
$x$ un objet de $\mathcal{A}$,
et $f : x \to x$ un idempotent.
En posant $g = \text{id} _ x - f$,
on a $f \circ g = 0$, $g \circ f = 0$ et $g \circ g = g$.
\end{remark}

\begin{lemma}
\label{lemma-idempotent-kernel-cokernel}
Soient $\mathcal{A}$ une catégorie préadditive,
$x$ un objet de $\mathcal{A}$,
et $f : x \to x$ un idempotent.
Posons $g = \text{id} _ x - f$
et rappelons, d'après la remarque \ref{remark-idempotent-symmetry},
que $f \circ g = g \circ f = 0$.
\begin{enumerate}
\item Si $f$ a pour noyau $i : \Ker (f) \to x$
et si $p$ est l'unique morphisme tel que $g = i \circ p$,
alors $p \circ i = \text{id} _ {\Ker (f)}$
et $p$ est le conoyau de $f$.
\item Si $f$ a pour conoyau $p : x \to \Coker (f)$
et si $i$ est l'unique morphisme tel que $g = i \circ p$,
alors $p \circ i = \text{id} _ {\Coker (f)}$
et $i$ est le noyau de $f$.
\item Si $g$ a pour noyau $j : \Ker (g) \to x$
et si $q$ est l'unique morphisme tel que $f = j \circ q$,
alors $q \circ j = \text{id} _ {\Ker (g)}$
et $q$ est le conoyau de $g$.
\item Si $g$ a pour conoyau $q : x \to \Coker (g)$
et si $j$ est l'unique morphisme tel que $f = j \circ q$,
alors $q \circ j = \text{id} _ {\Coker (g)}$
et $j$ est le noyau de $g$.
\end{enumerate}
\end{lemma}

\begin{proof}
Pour (1), calculons d'abord que $i = (f + g) \circ i = g \circ i$.
Par construction de $p$, cela entraîne
$i \circ p \circ i = g \circ i = i \circ \text{id} _ {\Ker (f)}$.
Le lemme \ref{lemma-kernel-mono} permet d'en conclure que
$p \circ i = \text{id} _ {\Ker (f)}$.

\medskip\noindent
Il reste à montrer que $p : x \to \Ker (f)$
satisfait la propriété universelle du conoyau de $f$,
c'est-à-dire que, pour tout morphisme $a : x \to y$ tel que $a \circ f = 0$,
il existe un unique morphisme $b : \Ker (f) \to y$
tel que $a = b \circ p$.
Pour l'existence, on calcule d'abord que $a = a \circ (f + g) = a \circ g$,
puis on déduit de la construction de $p$
que $a = a \circ g = a \circ i \circ p$
(on peut donc prendre $b = a \circ i$).
L'unicité résulte de $b = b \circ p \circ i = a \circ i$
(ainsi $a$ détermine $b$).

\medskip\noindent
Par dualité,
par symétrie (c'est-à-dire parce que $g$ est lui-même idempotent ;
cf. remarque \ref{remark-idempotent-symmetry}),
et respectivement par dualité et symétrie, (2), (3) et (4) se déduisent comme (1).
\end{proof}

\begin{lemma}
\label{lemma-idempotent-splitting}
Soient $\mathcal{A}$ une catégorie préadditive,
$x$ un objet de $\mathcal{A}$,
et $f : x \to x$ un idempotent.
Posons $g = \text{id} _ x - f$.
Si chacun de $f$ et $g$ possède au moins un noyau ou un conoyau,
et si l'on construit successivement les morphismes $i , j , p , q$
comme dans l'énoncé du lemme \ref{lemma-idempotent-kernel-cokernel},
alors les quatre (co)noyaux existent et
\begin{align*}
x
& \simeq \Ker (f) \oplus \Ker (g) \\
& \simeq \Coker (f) \oplus \Ker (g) \\
& \simeq \Ker (f) \oplus \Coker (g) \\
& \simeq \Coker (f) \oplus \Coker (g)
\end{align*}
la structure de chaque produit direct étant donnée par $i , j , p , q$.
\end{lemma}

\begin{proof}
D'après le lemme \ref{lemma-idempotent-kernel-cokernel},
$p \circ i = \text{id} _ {\Ker (f)} = \text{id} _ {\Coker (f)}$
et $q \circ j = \text{id} _ {\Ker (g)} = \text{id} _ {\Coker (g)}$.
Par construction, $i \circ p + j \circ q = f + g = \text{id} _ x$.
L'assertion résulte de la remarque \ref{remark-direct-sum}.
\end{proof}

\begin{lemma}
\label{lemma-split-morphism-kernel-cokernel}
Soient $\mathcal{A}$ une catégorie préadditive,
$x$ et $y$ des objets de $\mathcal{A}$,
et $j : y \to x$ et $q : x \to y$ des morphismes
vérifiant $q \circ j = \text{id} _ y$.
Alors :
\begin{enumerate}
\item $\text{id} _ x - j \circ q$ a pour noyau $j$.
\item $\text{id} _ x - j \circ q$ a pour conoyau $q$.
\end{enumerate}
\end{lemma}

\begin{proof}
Pour (1), calculons d'abord que $(\text{id} _ x - j \circ q) \circ j = 0$.
Nous devons montrer que $j$ satisfait la propriété universelle
du noyau de $\text{id} _ {x} - j \circ q$,
c'est-à-dire que, pour tout morphisme $a : z \to x$
tel que $(\text{id} _ x - j \circ q) \circ a = 0$,
il existe un unique morphisme $b : z \to y$
tel que $a = j \circ b$.
L'existence résulte du calcul $a = j \circ q \circ a$
(on peut donc prendre $b = q \circ a$).
L'unicité résulte de $b = q \circ j \circ b = q \circ a$
(ainsi $a$ détermine $b$).

\medskip\noindent
Par dualité, (2) se déduit comme (1).
\end{proof}

\begin{lemma}
\label{lemma-split-morphism-splitting}
Soient $\mathcal{A}$ une catégorie préadditive,
$x$ et $y$ des objets de $\mathcal{A}$,
et $j : y \to x$ et $q : x \to y$ des morphismes
vérifiant $q \circ j = \text{id} _ y$.
Si $j \circ q$ possède au moins un noyau ou un conoyau,
alors les deux (co)noyaux existent et
\begin{align*}
x
& \simeq \Ker (j \circ q) \oplus y \\
& \simeq \Coker (j \circ q) \oplus y
\end{align*}
la structure de chacun des produits directs comprenant $j , q$
et le morphisme canonique de (co)noyau,
la dernière application étant déterminée de manière unique.
\end{lemma}

\begin{proof}
Posons $f = j \circ q$,
calculons que $f \circ f = f$,
et appliquons le lemme \ref{lemma-idempotent-splitting}
et le lemme \ref{lemma-split-morphism-kernel-cokernel}.
\end{proof}




\section{Catégories karoubiennes}
\label{section-karoubian}

\noindent
On pourra omettre cette section lors d'une première lecture.

\begin{definition}
\label{definition-karoubian}
Soit $\mathcal{C}$ une catégorie préadditive. On dit que $\mathcal{C}$
est {\it karoubienne} si tout endomorphisme idempotent d'un objet
de $\mathcal{C}$ possède un noyau.
\end{definition}

\noindent
La notion duale demanderait que tout endomorphisme idempotent d'un
objet possède un conoyau. Toutefois, au vu du (dual du)
lemme suivant, cette notion serait équivalente.

\begin{lemma}
\label{lemma-karoubian}
Soit $\mathcal{C}$ une catégorie préadditive. Les assertions
suivantes sont équivalentes :
\begin{enumerate}
\item $\mathcal{C}$ est karoubienne,
\item tout endomorphisme idempotent d'un objet de $\mathcal{C}$ possède un
conoyau, et
\item pour tout endomorphisme idempotent $p : z \to z$ de $\mathcal{C}$,
il existe une décomposition en somme directe $z = x \oplus y$ telle
que $p$ corresponde à la projection sur $y$.
\end{enumerate}
\end{lemma}

\begin{proof}
Cela découle immédiatement des lemmes \ref{lemma-idempotent-kernel-cokernel},
\ref{lemma-idempotent-splitting} et \ref{lemma-additive-cat-biproduct-kernel}.
\end{proof}

\begin{lemma}
\label{lemma-projectors-have-images}
Soit $\mathcal{D}$ une catégorie préadditive.
\begin{enumerate}
\item Si $\mathcal{D}$ possède des produits dénombrables et les noyaux des applications
qui admettent un inverse à droite, alors $\mathcal{D}$ est karoubienne.
\item Si $\mathcal{D}$ possède des coproduits dénombrables et les conoyaux des
applications qui admettent un inverse à gauche, alors $\mathcal{D}$ est karoubienne.
\end{enumerate}
\end{lemma}

\begin{proof}
Soit $X$ un objet de $\mathcal{D}$ et soit $e : X \to X$ un idempotent.
Le foncteur
$$
W \longmapsto \Ker(
\Mor_\mathcal{D}(W, X)
\xrightarrow{e}
\Mor_\mathcal{D}(W, X)
)
$$
est représentable si et seulement si $e$ possède un noyau. Remarquons que, pour tout
groupe abélien $A$ et tout endomorphisme idempotent $e : A \to A$, on a
$$
\Ker(e : A \to A)
= \Ker(\Phi :
\prod\nolimits_{n \in \mathbf{N}} A
\to
\prod\nolimits_{n \in \mathbf{N}} A
)
$$
où
$$
\Phi(a_1, a_2, a_3, \ldots) = (ea_1 + (1 - e)a_2, ea_2 + (1 - e)a_3, \ldots)
$$
De plus, $\Phi$ possède l'inverse à droite
$$
\Psi(a_1, a_2, a_3, \ldots) =
(a_1, (1 - e)a_1 + ea_2, (1 - e)a_2 + ea_3, \ldots).
$$
Ainsi, (1) est vrai. La preuve de (2) est duale (en utilisant la définition duale
d'une catégorie karoubienne, à savoir la condition (2) du
lemme \ref{lemma-karoubian}).
\end{proof}








\section{Catégories abéliennes}
\label{section-abelian-categories}

\noindent
Une catégorie abélienne est une catégorie qui satisfait juste assez d'axiomes pour que le
lemme du serpent soit valable. L'un de ces axiomes (parfois oublié)
est que l'application canonique $\Coim(f) \to \Im(f)$
du lemme \ref{lemma-coim-im-map} soit toujours un isomorphisme.
L'exemple \ref{example-not-abelian} montre que cet axiome est nécessaire.

\begin{definition}
\label{definition-abelian-category}
Une catégorie $\mathcal{A}$ est {\it abélienne} si
elle est additive, si tous les noyaux et conoyaux existent,
et si l'application naturelle $\Coim(f) \to \Im(f)$
est un isomorphisme pour tout morphisme $f$ de
$\mathcal{A}$.
\end{definition}

\begin{lemma}
\label{lemma-abelian-opposite}
Soit $\mathcal{A}$ une catégorie préadditive.
Les additions sur les ensembles de morphismes munissent
$\mathcal{A}^{opp}$ d'une structure de catégorie préadditive.
De plus, $\mathcal{A}$ est additive si et seulement si $\mathcal{A}^{opp}$
est additive, et
$\mathcal{A}$ est abélienne si et seulement si $\mathcal{A}^{opp}$ est abélienne.
\end{lemma}

\begin{proof}
La première assertion est immédiate.
Pour voir que $\mathcal{A}$ est additive si et seulement si $\mathcal{A}^{opp}$
est additive, rappelons que l'additivité peut se caractériser par
l'existence d'un objet nul et de sommes directes, qui sont tous deux
préservés par passage à la catégorie opposée.
Enfin, pour voir que
$\mathcal{A}$ est abélienne si et seulement si $\mathcal{A}^{opp}$ est abélienne,
observons que les noyaux, conoyaux, images et coimages dans
$\mathcal{A}^{opp}$ correspondent respectivement aux
conoyaux, noyaux, coimages et images dans $\mathcal{A}$.
\end{proof}

\begin{definition}
\label{definition-injective-surjective}
Soit $f : x \to y$ un morphisme dans une catégorie abélienne.
\begin{enumerate}
\item On dit que $f$ est {\it injectif} si $\Ker(f) = 0$.
\item On dit que $f$ est {\it surjectif} si $\Coker(f) = 0$.
\end{enumerate}
Si $x \to y$ est injectif, on dit que $x$ est un {\it sous-objet}
de $y$ et l'on emploie la notation $x \subset y$. Si $x \to y$ est
surjectif, on dit que $y$ est un {\it quotient} de $x$.
\end{definition}

\begin{lemma}
\label{lemma-characterize-injective}
Soit $f : x \to y$ un morphisme dans une catégorie abélienne $\mathcal{A}$. Alors :
\begin{enumerate}
\item $f$ est injectif si et seulement si $f$ est un monomorphisme,
\item $f$ est surjectif si et seulement si $f$ est un épimorphisme, et
\item $f$ est un isomorphisme si et seulement si $f$ est injectif et surjectif.
\end{enumerate}
\end{lemma}

\begin{proof}
Prouvons (1). Rappelons que $\Ker(f)$ est un objet représentant le
foncteur qui envoie $z$ sur
$\Ker(\Mor_\mathcal{A}(z, x) \to \Mor_\mathcal{A}(z, y))$, voir la
définition \ref{definition-kernel}.
Ainsi, $\Ker(f)$ vaut $0$ si et seulement si
$\Mor_\mathcal{A}(z, x) \to \Mor_\mathcal{A}(z, y)$
est injective pour tout $z$, si et seulement si $f$ est un monomorphisme.
La preuve de (2) est analogue.
Pour prouver (3), remarquons qu'un isomorphisme est à la fois un monomorphisme
et un épimorphisme, ce qui, par (1) et (2), montre que $f$ est injectif
et surjectif. Si $f$ est à la fois injectif et surjectif, alors
$x = \Coim(f)$ et $y = \Im(f)$, d'où $f$ est un isomorphisme.
\end{proof}

\noindent
Dans une catégorie abélienne, si $x \subset y$ est un sous-objet,
on note
$$
y/x = \Coker(x \to y).
$$

\begin{lemma}
\label{lemma-colimit-abelian-category}
Soit $\mathcal{A}$ une catégorie abélienne.
Toutes les limites finies et toutes les colimites finies existent dans $\mathcal{A}$.
\end{lemma}

\begin{proof}
Pour montrer que les limites finies existent, il suffit de montrer
que les produits finis et les égalisateurs existent, voir
Catégories, lemme \ref{categories-lemma-finite-limits-exist}.
Les produits finis existent
par définition, et l'égalisateur de $a, b : x \to y$ est
le noyau de $a - b$. L'argument pour les colimites finies
est analogue, par dualité.
\end{proof}

\begin{example}
\label{example-fibre-product-pushouts}
Soit $\mathcal{A}$ une catégorie abélienne.
Les sommes amalgamées et les produits fibrés dans $\mathcal{A}$ admettent les
descriptions simples suivantes :
\begin{enumerate}
\item Si $a : x \to y$, $b : z \to y$ sont des morphismes de $\mathcal{A}$, alors
on a le produit fibré :
$x \times_y z = \Ker((a, -b) : x \oplus z \to y)$.
\item Si $a : y \to x$, $b : y \to z$ sont des morphismes de $\mathcal{A}$, alors
on a la somme amalgamée :
$x \amalg_y z = \Coker((a, -b) : y \to x \oplus z)$.
\end{enumerate}
\end{example}

\begin{definition}
\label{definition-exact}
Soit $\mathcal{A}$ une catégorie additive. Considérons une suite de morphismes
$$
\ldots \to x \to y \to z \to \ldots
\quad\text{ou}\quad
x_1 \to x_2 \to \ldots \to x_n
$$
dans $\mathcal{A}$. On dit qu'une telle suite est un {\it complexe} si la
composition de deux flèches consécutives quelconques (parmi celles représentées) est nulle.
Si $\mathcal{A}$ est abélienne, on dit qu'un complexe du premier
type ci-dessus est {\it exact en $y$} si $\Im(x \to y) = \Ker(y \to z)$,
et qu'un complexe du second type est {\it exact en $x_i$},
où $1 < i < n$, si
$\Im(x_{i - 1} \to x_i) = \Ker(x_i \to x_{i + 1})$. On dit qu'une
suite comme ci-dessus est {\it exacte}, ou est une {\it suite exacte}, ou est un
{\it complexe exact}, si c'est un complexe exact en tout objet (dans
le premier cas), ou exact en $x_i$ pour tout $1 < i < n$ (dans le second cas).
Il existe des variantes de ces notions pour les suites de la forme
$$
\ldots \to x_{-3} \to x_{-2} \to x_{-1}
\quad\text{et}\quad
x_1 \to x_2 \to x_3 \to \ldots
$$
Une {\it suite exacte courte} est un complexe exact de la forme
$$
0 \to A  \to B \to C \to 0.
$$
\end{definition}

\noindent
Dans le lemme suivant, nous supposons que le lecteur sait ce que signifie
l'exactitude d'une suite de groupes abéliens.

\begin{lemma}
\label{lemma-check-exactness}
Soit $\mathcal{A}$ une catégorie abélienne.
Soit $0 \to M_1 \to M_2 \to M_3 \to 0$ un complexe de $\mathcal{A}$.
\begin{enumerate}
\item $M_1 \to M_2 \to M_3 \to 0$ est exact si et seulement si
$$
0 \to \Hom_\mathcal{A}(M_3, N) \to
\Hom_\mathcal{A}(M_2, N) \to \Hom_\mathcal{A}(M_1, N)
$$
est une suite exacte de groupes abéliens pour tout objet $N$ de
$\mathcal{A}$, et
\item $0 \to M_1 \to M_2 \to M_3$ est exact si et seulement si
$$
0 \to \Hom_\mathcal{A}(N, M_1) \to \Hom_\mathcal{A}(N, M_2) \to
\Hom_\mathcal{A}(N, M_3)
$$
est une suite exacte de groupes abéliens pour tout objet $N$ de $\mathcal{A}$.
\end{enumerate}
\end{lemma}

\begin{proof}
Omis. Indication : voir
Algèbre, lemme \ref{algebra-lemma-hom-exact}.
\end{proof}

\begin{definition}
\label{definition-ses-split}
Soit $\mathcal{A}$ une catégorie abélienne.
Soient $i : A \to B$ et $q : B \to C$ des morphismes
de $\mathcal{A}$ tels que
$0 \to A \to B \to C \to 0$ soit une suite
exacte courte. On dit que la suite exacte courte
est {\it scindée} s'il existe
des morphismes $j : C \to B$ et $p : B \to A$ tels
que $(B, i, j, p, q)$ soit la somme directe de $A$ et $C$.
\end{definition}

\begin{lemma}
\label{lemma-ses-split}
Soit $\mathcal{A}$ une catégorie abélienne.
Soit $0 \to A \to B \to C \to 0$
une suite exacte courte.
\begin{enumerate}
\item Étant donné un morphisme $s : C \to B$ inverse à droite de
$B \to C$, il existe un unique $\pi : B \to A$
tel que $(s, \pi)$ scinde la suite exacte courte
au sens de la définition \ref{definition-ses-split}.
\item Étant donné un morphisme $\pi : B \to A$ inverse à gauche de
$A \to B$, il existe un unique $s : C \to B$
tel que $(s, \pi)$ scinde la suite exacte courte
au sens de la définition \ref{definition-ses-split}.
\end{enumerate}
\end{lemma}

\begin{proof}
C'est un corollaire immédiat du lemme \ref{lemma-split-morphism-splitting}.
\end{proof}

\begin{lemma}
\label{lemma-characterize-cartesian}
Soit $\mathcal{A}$ une catégorie abélienne. Soit
$$
\xymatrix{
w\ar[r]^f\ar[d]_g
& y\ar[d]^h\\
x\ar[r]^k
& z
}
$$
un diagramme commutatif. 
\begin{enumerate}
\item Le diagramme est cartésien si et seulement si 
$$
0 \to w \xrightarrow{(g, f)} x \oplus y \xrightarrow{(k, -h)} z
$$
est exact.
\item Le diagramme est cocartésien si et seulement si 
$$
w \xrightarrow{(g, -f)} x \oplus y \xrightarrow{(k, h)} z \to 0
$$
est exact.
\end{enumerate}
\end{lemma}

\begin{proof}
Posons $u = (g, f) : w \to x \oplus y$ et $v = (k, -h) : x \oplus y \to z$. 
Soient $p : x \oplus y \to x$ et $q : x \oplus y \to y$ les projections 
canoniques. Soit $i : \Ker(v) \to x \oplus y$ l'injection 
canonique. D'après l'exemple \ref{example-fibre-product-pushouts}, le diagramme est 
cartésien si et seulement s'il existe un isomorphisme 
$r : \Ker(v) \to w$ tel que $f \circ r = q \circ i$ et 
$g \circ r = p \circ i$. La suite 
$0 \to w \overset{u} \to x \oplus y \overset{v} \to z$ est exacte si et 
seulement s'il existe un isomorphisme $r : \Ker(v) \to w$ tel que 
$u \circ r = i$. Or, étant donné $r : \Ker(v) \to w$, on a 
$f \circ r = q \circ i$ et $g \circ r = p \circ i$ si et 
seulement si $q \circ u \circ r= f \circ r = q \circ i$ et 
$p \circ u \circ r = g \circ r = p \circ i$, donc si et seulement si
$u \circ r = i$. Cela prouve (1), puis (2) en découle par dualité. 
\end{proof}

\begin{lemma}
\label{lemma-cartesian-kernel}
Soit $\mathcal{A}$ une catégorie abélienne. Soit
$$
\xymatrix{
w\ar[r]^f\ar[d]_g
& y\ar[d]^h\\
x\ar[r]^k
& z
}
$$
un diagramme commutatif.
\begin{enumerate}
\item Si le diagramme est cartésien, alors le morphisme 
$\Ker(f)\to\Ker(k)$ induit par $g$ est un isomorphisme.
\item Si le diagramme est cocartésien, alors le morphisme 
$\Coker(f)\to\Coker(k)$ induit par $h$ est un isomorphisme.
\end{enumerate}
\end{lemma}

\begin{proof}
Supposons le diagramme cartésien. Soit 
$e:\Ker(f)\to\Ker(k)$ induit par $g$. Soient 
$i:\Ker(f)\to w$ et $j:\Ker(k)\to x$ les injections 
canoniques. Il existe $t:\Ker(k)\to w$ tel que $f\circ t=0$ 
et $g\circ t=j$. Il existe donc $u:\Ker(k)\to\Ker(f)$ 
tel que $i\circ u=t$. Il s'ensuit que 
$g\circ i\circ u\circ e=g\circ t\circ e=j\circ e=g\circ i$ et 
$f\circ i\circ u\circ e=0=f\circ i$, donc $i\circ u\circ e=i$. Comme 
$i$ est un monomorphisme, ceci implique $u\circ e=\text{id}_{\Ker(f)}$.
De plus, on a $j\circ e\circ u=g\circ i\circ u=g\circ t=j$. 
Comme $j$ est un monomorphisme, ceci implique $e\circ u=\text{id}_{\Ker(k)}$.
Cela prouve (1). L'assertion (2) en découle par dualité.
\end{proof}

\begin{lemma}
\label{lemma-cartesian-cocartesian}
Soit $\mathcal{A}$ une catégorie abélienne. Soit
$$
\xymatrix{
w\ar[r]^f\ar[d]_g
& y\ar[d]^h\\
x\ar[r]^k
& z
}
$$
un diagramme commutatif.
\begin{enumerate}
\item Si le diagramme est cartésien et si $k$ est un épimorphisme, 
alors le diagramme est cocartésien et $f$ est un épimorphisme.
\item Si le diagramme est cocartésien et si $g$ est un monomorphisme, 
alors le diagramme est cartésien et $h$ est un monomorphisme.
\end{enumerate}
\end{lemma}

\begin{proof}
Supposons le diagramme cartésien et $k$ épimorphique. 
Soit $u = (g, f) : w \to x \oplus y$ et soit $v = (k, -h) : x \oplus y \to z$. 
Comme $k$ est un épimorphisme, $v$ en est un également. Par conséquent, 
et d'après le lemme \ref{lemma-characterize-cartesian}, la suite 
$0\to w\overset{u}\to x\oplus y\overset{v}\to z\to 0$ est exacte. Le 
diagramme est donc cocartésien par le lemme \ref{lemma-characterize-cartesian}. Enfin, 
$f$ est un épimorphisme d'après le lemme \ref{lemma-cartesian-kernel} et le 
lemme \ref{lemma-characterize-injective}. Cela prouve (1), et (2) 
en découle par dualité.
\end{proof}

\begin{lemma}
\label{lemma-epimorphism-universal-abelian-category}
Soit $\mathcal{A}$ une catégorie abélienne.
\begin{enumerate}
\item Si $x \to y$ est surjectif, alors, pour tout $z \to y$, la
projection $x \times_y z \to z$ est surjective.
\item Si $x \to y$ est injectif, alors, pour tout $x \to z$, le
morphisme $z \to z \amalg_x y$ est injectif.
\end{enumerate}
\end{lemma}

\begin{proof}
Cela découle immédiatement du lemme \ref{lemma-characterize-injective} et du
lemme \ref{lemma-cartesian-cocartesian}.
\end{proof}

\begin{lemma}
\label{lemma-check-exactness-fibre-product}
Soit $\mathcal{A}$ une catégorie abélienne. Soient $f:x\to y$ et $g:y\to z$ 
des morphismes tels que $g\circ f=0$. Alors les assertions suivantes sont équivalentes :
\begin{enumerate}
\item La suite $x\overset{f}\to y\overset{g}\to z$ est exacte.
\item Pour tout $h:w\to y$ tel que $g\circ h=0$, il existe un objet $v$, 
un épimorphisme $k:v\to w$ et un morphisme $l:v\to x$ tels que $h\circ k=f\circ l$.
\end{enumerate}
\end{lemma}

\begin{proof}
Soit $i:\Ker(g)\to y$ l'injection canonique. Soit 
$p:x\to\Coim(f)$ la projection canonique. Soit 
$j:\Im(f)\to\Ker(g)$ l'injection canonique.

\medskip\noindent
Supposons (1). Soit $h:w\to y$ tel que $g\circ h=0$. Il existe 
$c:w\to\Ker(g)$ tel que $i\circ c=h$. 
Soit $v=x\times_{\Ker(g)}w$, avec projections canoniques 
$k:v\to w$ et $l:v\to x$, de sorte que $c\circ k=j\circ p\circ l$. 
Alors $h\circ k=i\circ c\circ k=i\circ j\circ p\circ l=f\circ l$. 
Comme $j\circ p$ est un épimorphisme par hypothèse, $k$ est un 
épimorphisme d'après le lemme \ref{lemma-cartesian-cocartesian}. Cela implique (2).

\medskip\noindent
Supposons (2). Alors $g\circ i=0$. Il existe donc un objet 
$w$, un épimorphisme $k:w\to\Ker(g)$ et un morphisme 
$l:w\to x$ tels que $f\circ l=i\circ k$. Il s'ensuit que 
$i\circ j\circ p\circ l=f\circ l=i\circ k$. Comme $i$ est un 
monomorphisme, on voit que $j\circ p\circ l=k$ est un épimorphisme. 
Ainsi, $j$ est un épimorphisme, donc un isomorphisme. Cela implique (1).
\end{proof}

\begin{lemma}
\label{lemma-exact-kernel-sequence}
Soit $\mathcal{A}$ une catégorie abélienne. Soit
$$
\xymatrix{
x \ar[r]^f \ar[d]^\alpha &
y \ar[r]^g \ar[d]^\beta &
z \ar[d]^\gamma\\
u \ar[r]^k & v \ar[r]^l & w
}
$$
un diagramme commutatif.
\begin{enumerate}
\item Si la première ligne est exacte et si $k$ est un monomorphisme, alors la suite 
induite $\Ker(\alpha) \to \Ker(\beta) \to \Ker(\gamma)$ 
est exacte.
\item Si la seconde ligne est exacte et si $g$ est un épimorphisme, alors la suite 
induite
$\Coker(\alpha) \to \Coker(\beta) \to \Coker(\gamma)$ 
est exacte.
\end{enumerate}
\end{lemma}

\begin{proof}
Supposons la première ligne exacte et $k$ monomorphique. Soient 
$a:\Ker(\alpha)\to\Ker(\beta)$ et 
$b:\Ker(\beta)\to\Ker(\gamma)$ les morphismes induits. 
Soient $h:\Ker(\alpha)\to x$, $i:\Ker(\beta)\to y$ et 
$j:\Ker(\gamma)\to z$ les injections canoniques. Comme $j$ est 
un monomorphisme, on a $b\circ a=0$. Soit $c:s\to\Ker(\beta)$ 
tel que $b\circ c=0$. Alors $g\circ i\circ c=j\circ b\circ c=0$. D'après le 
lemme \ref{lemma-check-exactness-fibre-product}, il existe un objet $t$, un 
épimorphisme $d:t\to s$ et un morphisme $e:t\to x$ tels que 
$i\circ c\circ d=f\circ e$. Alors 
$k\circ \alpha\circ e=\beta\circ f\circ e=\beta\circ i\circ c\circ d=0$. 
Comme $k$ est un monomorphisme, on obtient $\alpha\circ e=0$. Il existe donc 
$m:t\to\Ker(\alpha)$ tel que $h\circ m=e$. Il s'ensuit que 
$i\circ a\circ m=f\circ h\circ m=f\circ e=i\circ c\circ d$. 
Comme $i$ est un monomorphisme, on obtient $a\circ m=c\circ d$. Ainsi, le 
lemme \ref{lemma-check-exactness-fibre-product} implique (1), puis 
(2) en découle par dualité.
\end{proof}

\begin{lemma}
\label{lemma-snake}
Soit $\mathcal{A}$ une catégorie abélienne. Soit
$$
\xymatrix{
& x \ar[r]^f \ar[d]^\alpha &
y \ar[r]^g \ar[d]^\beta &
z \ar[r] \ar[d]^\gamma &
0 \\
0 \ar[r] & u \ar[r]^k & v \ar[r]^l & w
}
$$
un diagramme commutatif dont les lignes sont exactes.
\begin{enumerate}
\item Il existe un unique morphisme
$\delta : \Ker(\gamma) \to \Coker(\alpha)$
tel que le diagramme
$$
\xymatrix{
y \ar[d]_\beta &
y \times_z \Ker(\gamma) \ar[l]_{\pi'} \ar[r]^{\pi} &
\Ker(\gamma) \ar[d]^\delta \\
v \ar[r]^{\iota'} & \Coker(\alpha) \amalg_u v &
\Coker(\alpha) \ar[l]_\iota
}
$$
soit commutatif, où $\pi$ et $\pi'$ sont les projections canoniques,
et $\iota$ et $\iota'$ les coprojections canoniques.
\item La suite induite
$$
\Ker(\alpha) \xrightarrow{f'} \Ker(\beta) \xrightarrow{g'}
\Ker(\gamma) \xrightarrow{\delta} \Coker(\alpha) \xrightarrow{k'}
\Coker(\beta) \xrightarrow{l'} \Coker(\gamma)
$$
est exacte. Si $f$ est injectif, alors $f'$ l'est aussi, et si $l$ est
surjectif, alors $l'$ l'est aussi.
\end{enumerate}
\end{lemma}

\begin{proof}
Comme $\pi$ est un épimorphisme et $\iota$ un monomorphisme d'après le
lemme \ref{lemma-cartesian-cocartesian}, l'unicité de $\delta$ est claire.
Posons $p = y \times_z \Ker(\gamma)$ et $q = \Coker(\alpha) \amalg_u v$.
Soient $h : \Ker(\beta) \to y$, $i : \Ker(\gamma) \to z$ et
$j : \Ker(\pi) \to p$ les injections canoniques.
Soit $\pi'' : u \to \Coker(\alpha)$ la projection canonique.
En gardant à l'esprit le lemme \ref{lemma-cartesian-cocartesian}, on obtient le diagramme
commutatif à lignes exactes
$$
\xymatrix{
0 \ar[r] &
\Ker(\pi) \ar[r]^j &
p \ar[r]^{\pi} \ar[d]_{\pi'} &
\Ker(\gamma) \ar[d]_i \ar[r] & 0 \\
& x \ar[r]^f \ar[d]_\alpha & y \ar[r]^g \ar[d]_\beta &
z \ar[d]_\gamma \ar[r] & 0 \\
0 \ar[r] & u \ar[r]^k \ar[d]_{\pi''} &
v \ar[r]^l \ar[d]_{\iota'} & w & \\
0 \ar[r] & \Coker(\alpha) \ar[r]^\iota & q & &
}
$$
Comme $l \circ \beta \circ \pi' = \gamma \circ i \circ \pi = 0$ et que la troisième
ligne du diagramme ci-dessus est exacte, il existe $a : p \to u$
tel que $k \circ a = \beta \circ \pi'$. Comme le quadrilatère supérieur droit du
diagramme ci-dessus est cartésien, le lemme \ref{lemma-cartesian-kernel} fournit un
épimorphisme $b : x \to \Ker(\pi)$ tel que $\pi' \circ j \circ b = f$.
Il s'ensuit que
$k \circ a \circ j \circ b = \beta \circ \pi' \circ j \circ b =
\beta \circ f = k \circ \alpha$.
Comme $k$ est un monomorphisme, ceci implique $a \circ j \circ b = \alpha$. Il s'ensuit que
$\pi'' \circ a \circ j \circ b = \pi'' \circ \alpha = 0$. Comme $b$ est un
épimorphisme, ceci
implique $\pi'' \circ a \circ j = 0$. Par conséquent, puisque la ligne supérieure du diagramme
ci-dessus est exacte, il existe
$\delta : \Ker(\gamma) \to \Coker(\alpha)$ tel que
$\delta \circ \pi = \pi'' \circ a$. Il s'ensuit que
$\iota \circ \delta \circ \pi = \iota \circ \pi'' \circ a =
\iota' \circ k \circ a = \iota' \circ \beta \circ \pi'$,
comme souhaité.

\medskip\noindent
Comme le quadrilatère supérieur droit du diagramme ci-dessus est cartésien, il
existe $c : \Ker(\beta) \to p$ tel que $\pi' \circ c = h$ et $\pi \circ c = g'$.
Il s'ensuit que
$\iota \circ \delta \circ g' = \iota \circ \delta \circ \pi \circ c =
\iota' \circ \beta \circ \pi' \circ c = \iota' \circ \beta \circ h = 0$.
Comme $\iota$ est un monomorphisme, ceci implique $\delta \circ g' = 0$.

\medskip\noindent
Soit ensuite $d : r \to \Ker(\gamma)$ tel que $\delta \circ d = 0$. L'application du
lemme \ref{lemma-check-exactness-fibre-product} à la suite exacte
$p \xrightarrow{\pi} \Ker(\gamma) \to 0$ et à $d$ fournit un objet $s$,
un épimorphisme $m : s \to r$ et un morphisme $n : s \to p$ tels que
$\pi \circ n = d \circ m$. Comme
$\pi'' \circ a \circ n = \delta \circ d \circ m = 0$,
l'application du lemme \ref{lemma-check-exactness-fibre-product} à la suite exacte
$x \xrightarrow{\alpha} u \xrightarrow{\pi''} \Coker(\alpha)$ et à
$a \circ n$ fournit un objet $t$, un épimorphisme $\varepsilon : t \to s$ et
un morphisme $\zeta : t \to x$ tels que
$a \circ n \circ \varepsilon = \alpha \circ \zeta$.
On a
$\beta \circ \pi' \circ n \circ \varepsilon =
k \circ \alpha \circ \zeta = \beta \circ f \circ \zeta$.
Posons $\eta = \pi' \circ n \circ \varepsilon - f \circ \zeta : t \to y$. Alors,
$\beta \circ \eta = 0$. Il existe donc
$\vartheta : t \to \Ker(\beta)$ tel que $\eta = h \circ \vartheta$. On a
$i \circ g' \circ \vartheta = g \circ h \circ \vartheta =
g \circ \pi' \circ n \circ \varepsilon - g \circ f \circ \zeta =
i \circ \pi \circ n \circ \varepsilon = i \circ d \circ m \circ \varepsilon$.
Comme $i$ est un monomorphisme, on obtient
$g' \circ \vartheta = d \circ m \circ \varepsilon$.
Ainsi, puisque $m \circ \varepsilon$ est un épimorphisme, le
lemme \ref{lemma-check-exactness-fibre-product} implique que
$\Ker(\beta) \xrightarrow{g'} \Ker(\gamma) \xrightarrow{\delta} \Coker(\alpha)$
est exacte. L'assertion découle alors du lemme \ref{lemma-exact-kernel-sequence}
et de la dualité.
\end{proof}

\begin{lemma}
\label{lemma-snake-natural}
Soit $\mathcal{A}$ une catégorie abélienne. Soit 
$$
\xymatrix{
& & & x\ar[ld]\ar[rr]\ar[dd]^(.4)\alpha
& & y\ar[ld]\ar[rr]\ar[dd]^(.4)\beta
& & z\ar[ld]\ar[rr]\ar[dd]^(.4)\gamma
& & 0\\
& & x'\ar[rr]\ar[dd]^(.4){\alpha'}
& & y'\ar[rr]\ar[dd]^(.4){\beta'}
& & z'\ar[rr]\ar[dd]^(.4){\gamma'}
& & 0
& \\
& 0\ar[rr]
& & u\ar[ld]\ar[rr]
& & v\ar[ld]\ar[rr]
& & w\ar[ld]
& & \\
0\ar[rr]
& & u'\ar[rr]
& & v'\ar[rr]
& & w'
& & &
}
$$
un diagramme commutatif à lignes exactes. Alors le diagramme induit
$$
\xymatrix@C=15pt{
\Ker(\alpha) \ar[r] \ar[d] &
\Ker(\beta) \ar[r] \ar[d] &
\Ker(\gamma) \ar[r]^(.45){\delta} \ar[d] &
\Coker(\alpha) \ar[r] \ar[d] &
\Coker(\beta) \ar[r] \ar[d] &
\Coker(\gamma) \ar[d] \\
\Ker(\alpha') \ar[r] &
\Ker(\beta') \ar[r] &
\Ker(\gamma') \ar[r]^(.45){\delta'} &
\Coker(\alpha') \ar[r] &
\Coker(\beta') \ar[r] &
\Coker(\gamma')
}
$$
est commutatif.
\end{lemma}

\begin{proof}
Omis.
\end{proof}

\begin{lemma}
\label{lemma-four-lemma}
Soit $\mathcal{A}$ une catégorie abélienne. Soit
$$
\xymatrix{
w \ar[r] \ar[d]^\alpha & x \ar[r] \ar[d]^\beta & y \ar[r] \ar[d]^\gamma &
z \ar[d]^\delta \\
w' \ar[r] & x' \ar[r] & y' \ar[r] & z'
}
$$
un diagramme commutatif à lignes exactes.
\begin{enumerate}
\item Si $\alpha, \gamma$ sont surjectifs et $\delta$ injectif, alors
$\beta$ est surjectif.
\item Si $\beta, \delta$ sont injectifs et $\alpha$ surjectif, alors
$\gamma$ est injectif.
\end{enumerate}
\end{lemma}

\begin{proof}
Supposons $\alpha, \gamma$ surjectifs et $\delta$ injectif.
On peut remplacer $w'$ par $\Im(w' \to x')$, c'est-à-dire supposer
que $w' \to x'$ est injectif.
On peut remplacer $z$ par $\Im(y \to z)$, c'est-à-dire supposer que
$y \to z$ est surjectif. On peut alors appliquer le
lemme \ref{lemma-snake}
à
$$
\xymatrix{
& \Ker(y \to z) \ar[r] \ar[d] & y \ar[r] \ar[d] & z \ar[r] \ar[d] & 0 \\
0 \ar[r] & \Ker(y' \to z') \ar[r] & y' \ar[r] & z'
}
$$
pour conclure que $\Ker(y \to z) \to \Ker(y' \to z')$ est
surjectif. Enfin, on applique le
lemme \ref{lemma-snake}
à
$$
\xymatrix{
& w \ar[r] \ar[d] & x \ar[r] \ar[d] & \Ker(y \to z) \ar[r] \ar[d] & 0 \\
0 \ar[r] & w' \ar[r] & x' \ar[r] & \Ker(y' \to z')
}
$$
pour conclure que $x \to x'$ est surjectif. Cela prouve (1). La preuve
de (2) est duale.
\end{proof}

\begin{lemma}
\label{lemma-five-lemma}
\begin{reference}
\cite[Lemma 4.5 page 16]{Eilenberg-Steenrod}
\end{reference}
Soit $\mathcal{A}$ une catégorie abélienne. Soit
$$
\xymatrix{
v \ar[r] \ar[d]^\alpha &
w \ar[r] \ar[d]^\beta &
x \ar[r] \ar[d]^\gamma &
y \ar[r] \ar[d]^\delta &
z \ar[d]^\epsilon \\
v' \ar[r] & w' \ar[r] & x' \ar[r] & y' \ar[r] & z'
}
$$
un diagramme commutatif à lignes exactes. Si $\beta, \delta$
sont des isomorphismes, si $\epsilon$ est injectif et si $\alpha$ est surjectif,
alors $\gamma$ est un isomorphisme.
\end{lemma}

\begin{proof}
C'est une conséquence immédiate du
lemme \ref{lemma-four-lemma}.
\end{proof}



\section{Extensions}
\label{section-extensions}

\begin{definition}
\label{definition-extension}
Soit $\mathcal{A}$ une catégorie abélienne.
Soient $A, B \in \Ob(\mathcal{A})$.
Une {\it extension $E$ de $B$ par $A$} est une suite
exacte courte
$$
0 \to A \to E \to B \to 0.
$$
Un {\it morphisme d'extensions} entre deux
extensions $0 \to A \to E \to B \to 0$ et
$0 \to A \to F \to B \to 0$ désigne un morphisme
$f : E \to F$ dans $\mathcal{A}$ qui rend le diagramme
$$
\xymatrix{
0 \ar[r] &
A \ar[r] \ar[d]^{\text{id}} &
E \ar[r] \ar[d]^f &
B \ar[r] \ar[d]^{\text{id}} &
0 \\
0 \ar[r] &
A \ar[r] &
F \ar[r] &
B \ar[r] &
0
}
$$
commutatif.
Ainsi, les extensions de $B$ par $A$ forment une catégorie.
\end{definition}

\noindent
Par abus de langage, nous omettons souvent de mentionner les
morphismes $A \to E$ et $E \to B$, bien qu'ils fassent
assurément partie de la structure d'une extension.

\begin{definition}
\label{definition-ext-group}
Soit $\mathcal{A}$ une catégorie abélienne.
Soient $A, B \in \Ob(\mathcal{A})$.
L'ensemble des classes d'isomorphisme d'extensions
de $B$ par $A$ se note
$$
\Ext_\mathcal{A}(B, A).
$$
On l'appelle le {\it groupe $\Ext$}.
\end{definition}

\noindent
Cette définition est valide car, selon nos conventions,
$\Ob(\mathcal{A})$ est un ensemble, et donc
$\Ext_\mathcal{A}(B, A)$ est un ensemble.
Dans chacun des cas de « grandes » catégories abéliennes
énumérés dans Catégories, remarque \ref{categories-remark-big-categories},
on peut aussi vérifier directement que $\Ext_\mathcal{A}(B, A)$
est un ensemble. Nous verrons en outre plus loin que c'est
toujours le cas lorsque $\mathcal{A}$ a suffisamment de projectifs
ou suffisamment d'injectifs. Insérer ici une référence ultérieure.

\medskip\noindent
En fait, on peut munir $\Ext_\mathcal{A}(-, -)$ d'une structure de
foncteur
$$
\mathcal{A} \times \mathcal{A}^{opp} \longrightarrow \textit{Ensembles}, \quad
(A, B) \longmapsto \Ext_\mathcal{A}(B, A)
$$
comme suit :
\begin{enumerate}
\item Étant donnés un morphisme $B' \to B$ et une extension
$E$ de $B$ par $A$, on définit $E' = E \times_B B'$.
D'après le lemme \ref{lemma-cartesian-kernel}, on a le diagramme
commutatif suivant de suites exactes courtes :
$$
\xymatrix{
0 \ar[r] & A \ar[r] \ar[d] & E' \ar[r] \ar[d] & B' \ar[r] \ar[d] & 0 \\
0 \ar[r] & A \ar[r] & E \ar[r] & B \ar[r] & 0
}
$$
L'extension $E'$ est appelée l'{\it image réciproque de $E$ par
$B' \to B$}.
\item Étant donnés un morphisme $A \to A'$ et une extension
$E$ de $B$ par $A$, on définit $E' = A' \amalg_A E$.
D'après le lemme \ref{lemma-cartesian-cocartesian},
on a le diagramme commutatif suivant
de suites exactes courtes :
$$
\xymatrix{
0 \ar[r] & A \ar[r] \ar[d] & E \ar[r] \ar[d] & B \ar[r] \ar[d] & 0 \\
0 \ar[r] & A' \ar[r] & E' \ar[r] & B \ar[r] & 0
}
$$
L'extension $E'$ est appelée l'{\it image directe de $E$ par
$A \to A'$}.
\end{enumerate}
Pour voir que cela définit bien le foncteur indiqué ci-dessus,
il faut vérifier plusieurs choses. Tout d'abord, la
fonctorialité en la variable $B$ exige que
$(E \times_B B') \times_{B'} B'' = E \times_B B''$,
ce qui est une propriété générale des produits fibrés.
Par dualité, on traite la fonctorialité en la
variable $A$. Enfin, étant donnés $A \to A'$ et
$B' \to B$, il faut montrer que
$$
A' \amalg_A (E \times_B B')
\cong
(A' \amalg_A E)\times_B B'
$$
comme extensions de $B'$ par $A'$. Rappelons que $A' \amalg_A E$
est un quotient de $A' \oplus E$.
Le membre de droite est donc un quotient de
$A' \oplus E \times_B B'$, et l'on vérifie immédiatement que
le noyau est précisément celui qu'il faut pour
obtenir le membre de gauche.

\medskip\noindent
Remarquons que si $E_1$ et $E_2$ sont des extensions de
$B$ par $A$, alors $E_1\oplus E_2$ est une extension
de $B \oplus B$ par $A\oplus A$. On prend
l'image directe par l'application somme $A \oplus A \to A$, puis
l'image réciproque par l'application diagonale $B \to B \oplus B$, pour obtenir
une extension $E_1 + E_2$ de $B$ par $A$.
$$
\xymatrix{
0 \ar[r] &
A \oplus A \ar[r] \ar[d]_\Sigma &
E_1 \oplus E_2 \ar[r] \ar[d] &
B \oplus B \ar[r] \ar[d] &
0 \\
0 \ar[r] &
A \ar[r] &
E' \ar[r] &
B \oplus B \ar[r] &
0\\
0 \ar[r] &
A \ar[r] \ar[u] &
E_1 + E_2 \ar[r] \ar[u] &
B \ar[r] \ar[u]^\Delta &
0
}
$$
L'extension $E_1 + E_2$ est appelée la {\it somme de Baer} des
extensions données.

\begin{lemma}
\label{lemma-baer-sum}
La construction $(E_1, E_2) \mapsto E_1 + E_2$
ci-dessus définit une loi de groupe commutative
sur $\Ext_\mathcal{A}(B, A)$, fonctorielle
en les deux variables.
\end{lemma}

\begin{proof}
Omis.
\end{proof}

\begin{lemma}
\label{lemma-six-term-sequence-ext}
Soit $\mathcal{A}$ une catégorie abélienne.
Soit $0 \to M_1 \to M_2 \to M_3 \to 0$ une suite exacte courte
dans $\mathcal{A}$.
\begin{enumerate}
\item Il existe une suite exacte canonique à six termes de groupes abéliens
$$
\xymatrix{
0 \ar[r] &
\Hom_\mathcal{A}(M_3, N) \ar[r] &
\Hom_\mathcal{A}(M_2, N) \ar[r] &
\Hom_\mathcal{A}(M_1, N) \ar[lld] \\
& \Ext_\mathcal{A}(M_3, N) \ar[r] &
\Ext_\mathcal{A}(M_2, N) \ar[r] &
\Ext_\mathcal{A}(M_1, N)
}
$$
pour tout objet $N$ de $\mathcal{A}$, et
\item il existe une suite exacte canonique à six termes de groupes abéliens
$$
\xymatrix{
0 \ar[r] &
\Hom_\mathcal{A}(N, M_1) \ar[r] &
\Hom_\mathcal{A}(N, M_2) \ar[r] &
\Hom_\mathcal{A}(N, M_3) \ar[lld] \\
& \Ext_\mathcal{A}(N, M_1) \ar[r] &
\Ext_\mathcal{A}(N, M_2) \ar[r] &
\Ext_\mathcal{A}(N, M_3)
}
$$
pour tout objet $N$ de $\mathcal{A}$.
\end{enumerate}
\end{lemma}

\begin{proof}
Omis. Indication : les morphismes de bord sont définis en utilisant soit l'image directe,
soit l'image réciproque de la suite exacte courte donnée.
\end{proof}







\section{Foncteurs additifs}
\label{section-functors}

\noindent
Commençons par un lemme tout à fait élémentaire qui caractérise les foncteurs additifs
entre catégories additives.

\begin{lemma}
\label{lemma-additive-functor}
Soient $\mathcal{A}$ et $\mathcal{B}$ des catégories additives.
Soit $F : \mathcal{A} \to \mathcal{B}$ un foncteur.
Les assertions suivantes sont équivalentes :
\begin{enumerate}
\item $F$ est additif,
\item $F(A) \oplus F(B) \to F(A \oplus B)$ est un isomorphisme pour
tous $A, B \in \mathcal{A}$, et
\item $F(A \oplus B) \to F(A) \oplus F(B)$ est un isomorphisme pour
tous $A, B \in \mathcal{A}$.
\end{enumerate}
\end{lemma}

\begin{proof}
Les foncteurs additifs commutent aux sommes directes d'après le
lemme \ref{lemma-additive-additive} ; ainsi (1)
implique (2) et (3). Si (2) est vrai et si $0$ désigne
l'objet nul de $\mathcal{A}$, alors le fait que l'application
$F(0) \oplus F(0) \to F(0 \oplus 0) \cong F(0)$
soit un isomorphisme implique que l'application diagonale
$\Hom(F(0), C) \to \Hom(F(0), C) \oplus \Hom(F(0), C)$
est bijective pour tout $C$ dans $\mathcal{B}$. Ainsi, $F(0)$
est l'objet nul de $\mathcal{B}$ et $F$ envoie le
morphisme nul entre deux objets quelconques sur le morphisme nul.
On obtient de même cette conclusion si (3) est vrai. Il s'ensuit que (2) et (3)
sont équivalents, car la composition
$F(A) \oplus F(B) \to F(A \oplus B) \to F(A) \oplus F(B)$
est donnée par la matrice
$$
\left(
\begin{matrix}
F(1) & F(0) \\
F(0) & F(1)
\end{matrix}
\right)
$$
et est donc l'application identité. Supposons (2) et (3) vraies.
Soient $f, g : A \to B$ des applications. Alors $f + g$ est égal à
la composition
$$
A \to A \oplus A \xrightarrow{\text{diag}(f, g)} B \oplus B \to B
$$
Appliquons le foncteur $F$ et considérons le diagramme suivant :
$$
\xymatrix{
F(A) \ar[r] \ar[rd] &
F(A \oplus A) \ar[rr]_{F(\text{diag}(f, g))} & &
F(B \oplus B) \ar[r] \ar[d] &
F(B) \\
&
F(A) \oplus F(A) \ar[u] \ar[rr]^{\text{diag}(F(f), F(g))} & &
F(B) \oplus F(B) \ar[ru]
}
$$
Affirmons qu'il est commutatif. Pour le carré central, on peut le vérifier
séparément pour chacune des quatre applications induites $F(A) \to F(B)$ ;
cela résulte du fait que $F$ est un foncteur et qu'il
transforme les morphismes nuls en morphismes nuls.
Pour le triangle de gauche, on utilise que
$F(A \oplus A) \to F(A) \oplus F(A)$ est un isomorphisme :
il suffit donc d'effectuer la vérification après composition avec
cette application, et cette vérification est immédiate. Par dualité, il en va de même de l'autre triangle.
Ainsi, le chemin inférieur est égal à $F(f + g)$, ce qui permet de conclure.
\end{proof}

\noindent
Rappelons que nous avons défini, dans
Catégories, définition \ref{categories-definition-exact},
les notions de foncteur « exact à droite », « exact à gauche » et
« exact » dans le cadre d'un foncteur entre
catégories possédant les (co)limites finies. Cela
s'applique donc en particulier aux foncteurs entre catégories
abéliennes.

\begin{lemma}
\label{lemma-exact-functor}
Soient $\mathcal{A}$ et $\mathcal{B}$ des catégories abéliennes.
Soit $F : \mathcal{A} \to \mathcal{B}$ un foncteur.
\begin{enumerate}
\item Si $F$ est exact à gauche ou exact à droite, alors il est additif.
\item $F$ est exact à gauche si et seulement si,
pour toute suite exacte courte
$0 \to A \to B \to C \to 0$,
la suite $0 \to F(A) \to F(B) \to F(C)$
est exacte.
\item $F$ est exact à droite si et seulement si, pour toute suite exacte courte
$0 \to A \to B \to C \to 0$,
la suite $F(A) \to F(B) \to F(C) \to 0$
est exacte.
\item $F$ est exact si et seulement si, pour toute suite exacte courte
$0 \to A \to B \to C \to 0$,
la suite $0 \to F(A) \to F(B) \to F(C) \to 0$
est exacte.
\end{enumerate}
\end{lemma}

\begin{proof}
Si $F$ est exact à gauche, c'est-à-dire si $F$ commute aux limites finies, alors
$F$ envoie les produits sur les produits, donc $F$ préserve les sommes directes,
et $F$ est donc additif d'après le lemme \ref{lemma-additive-functor}.
D'autre part, supposons que, pour toute suite exacte courte
$0 \to A \to B \to C \to 0$, la suite $0 \to F(A) \to F(B) \to F(C)$
soit exacte. Soient $A, B$ deux objets. On a alors une suite
exacte courte
$$
0 \to A \to A \oplus B \to B \to 0
$$
voir par exemple le lemme \ref{lemma-additive-cat-biproduct-kernel}.
Par hypothèse, la ligne inférieure du diagramme commutatif
$$
\xymatrix{
0 \ar[r] &
F(A) \ar[d] \ar[r] &
F(A) \oplus F(B) \ar[r] \ar[d] &
F(B) \ar[d] \ar[r] &
0 \\
0 \ar[r] &
F(A) \ar[r] &
F(A \oplus B) \ar[r] &
F(B)
}
$$
est exacte. Le lemme du serpent (lemme \ref{lemma-snake}) permet donc
de conclure que $F(A) \oplus F(B) \to F(A \oplus B)$ est un
isomorphisme. Ainsi, $F$ est également additif dans ce cas.
Dans la suite de la preuve, on peut donc supposer $F$ additif.

\medskip\noindent
Notons $f : B \to C$ une application de $B$ vers $C$.
L'exactitude de $0 \to A \to B \to C$ signifie précisément que
$A = \Ker(f)$. Il est clair que le noyau de $f$ est
l'égalisateur des deux applications $f$ et $0$ de $B$ vers $C$.
Ainsi, si $F$ commute aux limites, alors $F(\Ker(f))
= \Ker(F(f))$, ce qui signifie exactement que
$0 \to F(A) \to F(B) \to F(C)$ est exacte.

\medskip\noindent
Réciproquement, supposons que $F$ soit additif et
transforme toute suite exacte courte $0 \to A \to B \to C \to 0$ en
une suite exacte $0 \to F(A) \to F(B) \to F(C)$.
Puisqu'il est additif, il commute aux sommes directes,
et donc aux produits finis dans $\mathcal{A}$. Pour montrer
qu'il commute aux limites finies, il suffit donc
de montrer qu'il commute aux
égalisateurs. Or, dans une catégorie abélienne, les égalisateurs
sont les noyaux des applications différences ;
il suffit donc de montrer que $F$ commute à la
prise de noyaux. Soit $f : A \to B$ un morphisme.
Factorisons $f$ sous la forme $A \to I \to B$, avec $f' : A \to I$ surjectif
et $i : I \to B$ injectif. (Cela est possible par la
définition d'une catégorie abélienne.) Il est alors
clair que $\Ker(f) = \Ker(f')$. De plus,
$0 \to \Ker(f') \to A \to I \to 0$
et
$0 \to I \to B \to B/I \to 0$
sont des suites exactes courtes. La condition imposée à $F$
montre que
$0 \to F(\Ker(f')) \to F(A) \to F(I)$
et
$0 \to F(I) \to F(B) \to F(B/I)$
sont exactes. Par conséquent,
$F(\Ker(f'))$ est aussi le noyau de l'application
$F(A) \to F(B)$, ce qui conclut la preuve.

\medskip\noindent
La preuve de (3) est analogue à celle de (2).
L'assertion (4) combine (2) et (3).
\end{proof}

\begin{lemma}
\label{lemma-exact-functor-ext}
Soient $\mathcal{A}$ et $\mathcal{B}$ des catégories abéliennes.
Soit $F : \mathcal{A} \to \mathcal{B}$ un foncteur exact.
Pour toute paire d'objets $A, B$ de $\mathcal{A}$, le
foncteur $F$ induit un homomorphisme de groupes abéliens
$$
\Ext_\mathcal{A}(B, A)
\longrightarrow
\Ext_\mathcal{B}(F(B), F(A))
$$
qui envoie l'extension $E$ sur $F(E)$.
\end{lemma}

\begin{proof}
Omis.
\end{proof}

\noindent
Le lemme suivant est utilisé pour montrer que la catégorie des faisceaux
abéliens sur un site est abélienne, le foncteur $b$ étant la faisceautisation.

\begin{lemma}
\label{lemma-adjoint-get-abelian}
Soient $a : \mathcal{A} \to \mathcal{B}$ et $b : \mathcal{B} \to \mathcal{A}$
des foncteurs. Supposons que :
\begin{enumerate}
\item $\mathcal{A}$, $\mathcal{B}$ soient des catégories additives,
$a$, $b$ des foncteurs additifs, et $a$ un adjoint à droite de $b$,
\item $\mathcal{B}$ soit abélienne et $b$ exact à gauche, et
\item $ba \cong \text{id}_\mathcal{A}$.
\end{enumerate}
Alors $\mathcal{A}$ est abélienne.
\end{lemma}

\begin{proof}
Puisque $\mathcal{B}$ est abélienne, toutes les limites et colimites finies
existent dans $\mathcal{B}$ d'après le lemme \ref{lemma-colimit-abelian-category}.
Comme $b$ est un adjoint à gauche, $b$ est aussi exact à droite,
et donc exact, voir
Catégories, lemme \ref{categories-lemma-exact-adjoint}.
Soit $\varphi : B_1 \to B_2$ un morphisme de $\mathcal{B}$.
En particulier, si $K = \Ker(B_1 \to B_2)$, alors $K$ est
l'égalisateur de $0$ et $\varphi$ ; par conséquent,
$bK$ est l'égalisateur de $0$ et $b\varphi$, et donc
$bK$ est le noyau de $b\varphi$. De même, si
$Q = \Coker(B_1 \to B_2)$, alors $Q$ est
le coégalisateur de $0$ et $\varphi$ ; par conséquent,
$bQ$ est le coégalisateur de $0$ et $b\varphi$, et donc
$bQ$ est le conoyau de $b\varphi$. On voit ainsi que tout morphisme
de la forme $b\varphi$ dans $\mathcal{A}$ possède un noyau et un conoyau.
Or, puisque $ba \cong \text{id}$, tout morphisme de
$\mathcal{A}$ est de cette forme ; on conclut donc que les noyaux et
les conoyaux existent dans $\mathcal{A}$. En fait, l'argument montre que,
si $\psi : A_1 \to A_2$ est un morphisme, alors
$$
\Ker(\psi) = b\Ker(a\psi),
\quad\text{et}\quad
\Coker(\psi) = b\Coker(a\psi).
$$
Il reste à montrer que $\Coim(\psi)= \Im(\psi)$.
Procédons comme suit.
Remarquons d'abord que, puisque $\mathcal{A}$ possède des noyaux et des conoyaux, elle
possède toutes les limites et colimites finies (voir la preuve du
lemme \ref{lemma-colimit-abelian-category}).
On voit donc, par Catégories, lemme \ref{categories-lemma-exact-adjoint},
que $a$ est exact à gauche et
transforme par conséquent les noyaux (= égalisateurs) en noyaux.
\begin{align*}
\Coim(\psi)
& =
\Coker(\Ker(\psi) \to A_1)
& \text{par définition} \\
& =
b\Coker(a(\Ker(\psi) \to A_1))
& \text{par la formule ci-dessus} \\
& =
b\Coker(\Ker(a\psi) \to aA_1)
& a\text{ préserve les noyaux} \\
& =
b\Coim(a\psi)
& \text{par définition} \\
& =
b\Im(a\psi)
& \mathcal{B}\text{ est abélienne} \\
& =
b\Ker(aA_2 \to \Coker(a\psi))
& \text{par définition} \\
& =
\Ker(baA_2 \to b\Coker(a\psi))
& b\text{ préserve les noyaux} \\
& =
\Ker(A_2 \to b\Coker(a\psi))
& ba = \text{id}_\mathcal{A} \\
& =
\Ker(A_2 \to \Coker(\psi))
& \text{par la formule ci-dessus} \\
& =
\Im(\psi)
& \text{par définition}
\end{align*}
Le lemme est donc démontré.
\end{proof}





\section{Localisation}
\label{section-localization}

\noindent
Dans cette section, nous relevons comment la localisation de Gabriel-Zisman interagit avec
la structure additive d'une catégorie.

\begin{lemma}
\label{lemma-localization-preadditive}
Soit $\mathcal{C}$ une catégorie préadditive.
Soit $S$ un système multiplicatif à gauche ou à droite.
Il existe une unique structure préadditive sur
$S^{-1}\mathcal{C}$ telle que le foncteur de localisation
$Q : \mathcal{C} \to S^{-1}\mathcal{C}$ soit additif.
\end{lemma}

\begin{proof}
Nous allons le démontrer lorsque $S$ est un système multiplicatif à gauche.
Le cas où $S$ est un système multiplicatif à droite est dual.
Supposons que $X, Y$ soient des objets de $\mathcal{C}$ et que
$\alpha, \beta : X \to Y$ soient des morphismes de $S^{-1}\mathcal{C}$. D'après
Catégories, lemme \ref{categories-lemma-morphisms-left-localization},
on peut les représenter par des paires $s^{-1}f, s^{-1}g$ de dénominateur commun $s$.
Supposons que $S^{-1}\mathcal{C}$ soit muni d'une structure préadditive et
que $Q$ soit additif. Alors
\begin{align*}
s^{-1}f+s^{-1}g
&=Q(s)^{-1}\circ Q(f)+Q(s)^{-1}\circ Q(g)\\
&=Q(s)^{-1}\circ(Q(f)+Q(g))\\
&=Q(s)^{-1}\circ Q(f+g)\\
&=s^{-1}(f+g).
\end{align*}
Ainsi, si la structure préadditive sur $S^{-1}\mathcal{C}$ existe, elle est unique.
En conséquence, on définit $\alpha + \beta$ comme la
classe d'équivalence de $s^{-1}(f + g)$ ; dans la suite de la preuve, on montre
que cette définition ne dépend pas des choix et que la composition est bilinéaire.
Une fois cela établi, il est clair que $Q$ est un foncteur additif.

\medskip\noindent
Montrons que la construction ci-dessus ne dépend pas des choix.
Une manière abstraite de l'exprimer est de dire que les colimites filtrantes de
groupes abéliens coïncident avec les colimites filtrantes d'ensembles, puis d'utiliser
Catégories,
équation (\ref{categories-equation-left-localization-morphisms-colimit}).
On peut détailler quelque peu cet argument comme suit.
Soient $s : Y \to Y_1$ et $f, g : X \to Y_1$. Supposons disposer d'une seconde
représentation de $\alpha, \beta$ sous la forme $(s')^{-1}f', (s')^{-1}g'$, avec
$s' : Y \to Y_2$ et $f', g' : X \to Y_2$. D'après
Catégories, remarque \ref{categories-remark-left-localization-morphisms-colimit},
on peut trouver un morphisme $s_3 : Y \to Y_3$ et des morphismes
$a_1 : Y_1 \to Y_3$, $a_2 : Y_2 \to Y_3$ tels que
$a_1 \circ s = s_3 = a_2 \circ s'$ et aussi
$a_1 \circ f = a_2 \circ f'$ et $a_1 \circ g = a_2 \circ g'$.
On voit ainsi que $s^{-1}(f + g)$ est équivalent à
\begin{align*}
s_3^{-1}(a_1 \circ (f + g)) & =
s_3^{-1}(a_1 \circ f + a_1 \circ g) \\
& = s_3^{-1}(a_2 \circ f' + a_2 \circ g') \\
& = s_3^{-1}(a_2 \circ (f' + g'))
\end{align*}
qui est équivalent à $(s')^{-1}(f' + g')$.

\medskip\noindent
Fixons $s : Y \to Y'$ et $f, g : X \to Y'$ tels que
$\alpha = s^{-1}f$ et $\beta = s^{-1}g$ comme morphismes $X \to Y$
dans $S^{-1}\mathcal{C}$.
Pour montrer que la composition est bilinéaire, considérons d'abord le cas d'un
morphisme $\gamma : Y \to Z$ dans $S^{-1}\mathcal{C}$. Écrivons $\gamma = t^{-1}h$
pour certains $h : Y \to Z'$ et $t : Z \to Z'$ dans $S$. En utilisant LMS2, on
choisit des morphismes $a : Y' \to Z''$ et $t' : Z' \to Z''$ dans $S$ tels
que $a \circ s = t' \circ h$. Représentons la situation :
$$
\xymatrix{
& & Z \ar[d]^t \\
& Y \ar[r]^h \ar[d]^s & Z' \ar[d]^{t'} \\
X \ar[r]^{f, g} & Y' \ar[r]^a & Z''
}
$$
Alors
$\gamma \circ \alpha = (t' \circ t)^{-1}(a \circ f)$ et
$\gamma \circ \beta = (t' \circ t)^{-1}(a \circ g)$.
On voit donc que $\gamma \circ (\alpha + \beta)$ est représenté
par $(t' \circ t)^{-1}(a \circ (f + g)) =
(t' \circ t)^{-1}(a \circ f + a \circ g)$, qui représente
$\gamma \circ \alpha + \gamma \circ \beta$.

\medskip\noindent
Supposons enfin que $\delta : W \to X$ soit un autre morphisme de
$S^{-1}\mathcal{C}$. Écrivons $\delta = r^{-1}i$ pour certains
$i : W \to X'$ et $r : X \to X'$ dans $S$. Nous affirmons que l'on peut trouver
un morphisme $s' : Y' \to Y''$ dans $S$ et des morphismes $a'', b'' : X' \to Y''$
tels que le diagramme suivant soit commutatif :
$$
\xymatrix{
& & & Y \ar[d]^s \\
& X \ar[rr]^{f, g, f + g} \ar[d]^r & & Y' \ar[d]^{s'} \\
W \ar[r]^i & X' \ar[rr]^{a'', b'', a'' + b''} & & Y''
}
$$
En effet, en utilisant LMS2, on peut d'abord choisir
$s_1 : Y' \to Y_1$, $s_2 : Y' \to Y_2$ dans $S$ et
$a : X' \to Y_1$, $b : X' \to Y_2$ tels que
$a \circ r = s_1 \circ f$ et $b \circ r = s_2 \circ g$.
Puis, puisque la catégorie $Y'/S$ est filtrante (voir
Catégories, remarque \ref{categories-remark-left-localization-morphisms-colimit}),
on peut
trouver $s' : Y' \to Y''$ et des morphismes $a' : Y_1 \to Y''$, $b' : Y_2 \to Y''$
tels que $s' = a' \circ s_1$ et $s' = b' \circ s_2$. Il suffit alors de poser
$a'' = a' \circ a$ et $b'' = b' \circ b$.
À ce stade, on voit que les compositions
$\alpha \circ \delta$ et $\beta \circ \delta$ sont représentées par
$(s' \circ s)^{-1}(a'' \circ i)$ et $(s' \circ s)^{-1}(b'' \circ i)$.
Ainsi, $\alpha \circ \delta + \beta \circ \delta$ est représenté
par $(s' \circ s)^{-1}(a'' \circ i + b'' \circ i) =
(s' \circ s)^{-1}((a'' + b'') \circ i)$,
qui, toujours d'après le diagramme, est un représentant
de $(\alpha + \beta) \circ \delta$.
\end{proof}

\begin{lemma}
\label{lemma-localization-additive}
Soit $\mathcal{C}$ une catégorie additive.
Soit $S$ un système multiplicatif à gauche ou à droite.
Alors $S^{-1}\mathcal{C}$ est une catégorie additive et le foncteur de localisation
$Q : \mathcal{C} \to S^{-1}\mathcal{C}$ est additif.
\end{lemma}

\begin{proof}
D'après le lemme \ref{lemma-localization-preadditive},
$S^{-1}\mathcal{C}$ est préadditive et $Q$ est additif.
Le lemme \ref{lemma-additive-additive} permet d'en conclure que $S^{-1}\mathcal{C}$
possède un objet nul et des sommes directes.
\end{proof}

\noindent
Le lemme suivant décrit le « noyau »
du foncteur de localisation lorsque l'on inverse un système multiplicatif.

\begin{lemma}
\label{lemma-kernel-localization}
Soit $\mathcal{C}$ une catégorie additive. Soit $S$ un système
multiplicatif. Soit $X$ un objet
de $\mathcal{C}$. Les assertions suivantes sont équivalentes :
\begin{enumerate}
\item $Q(X) = 0$ dans $S^{-1}\mathcal{C}$,
\item il existe $Y \in \Ob(\mathcal{C})$ tel que
$0 : X \to Y$ appartienne à $S$, et
\item il existe $Z \in \Ob(\mathcal{C})$ tel que
$0 : Z \to X$ appartienne à $S$.
\end{enumerate}
\end{lemma}

\begin{proof}
Si (2) est vrai, alors $0 = Q(0) : Q(X) \to Q(Y)$ est un isomorphisme.
Dans la catégorie additive $S^{-1}\mathcal{C}$, ceci implique que $Q(X) = 0$.
Ainsi, (2) $\Rightarrow$ (1). De même, (3) $\Rightarrow$ (1).
Supposons que $Q(X) = 0$. Cela implique que le morphisme
$f : 0 \to X$ est transformé en isomorphisme dans $S^{-1}\mathcal{C}$.
Par conséquent, d'après
Catégories, lemme \ref{categories-lemma-what-gets-inverted},
il existe un morphisme $g : Z \to 0$ tel que $fg \in S$. Cela prouve
(1) $\Rightarrow$ (3). De même, (1) $\Rightarrow$ (2).
\end{proof}

\begin{lemma}
\label{lemma-localization-abelian}
Soit $\mathcal{A}$ une catégorie abélienne.
\begin{enumerate}
\item Si $S$ est un système multiplicatif à gauche, alors
la catégorie $S^{-1}\mathcal{A}$ possède des conoyaux et le foncteur
$Q : \mathcal{A} \to S^{-1}\mathcal{A}$ commute avec eux.
\item Si $S$ est un système multiplicatif à droite, alors
la catégorie $S^{-1}\mathcal{A}$ possède des noyaux et le foncteur
$Q : \mathcal{A} \to S^{-1}\mathcal{A}$ commute avec eux.
\item Si $S$ est un système multiplicatif, alors la catégorie
$S^{-1}\mathcal{A}$ est abélienne et le foncteur
$Q : \mathcal{A} \to S^{-1}\mathcal{A}$ est exact.
\end{enumerate}
\end{lemma}

\begin{proof}
Supposons que $S$ soit un système multiplicatif à gauche. Soit $a : X \to Y$ un morphisme
de $S^{-1}\mathcal{A}$. Alors $a = s^{-1}f$ pour certains $s : Y \to Y'$
dans $S$ et $f : X \to Y'$. Comme $Q(s)$ est un isomorphisme, l'existence
de $\Coker(a : X \to Y)$ équivaut à celle
de $\Coker(Q(f) : X \to Y')$. Puisque $\Coker(Q(f))$ est le
coégalisateur de $0$ et $Q(f)$, on voit que $\Coker(Q(f))$ est
représenté par $Q(\Coker(f))$, d'après
Catégories, lemme \ref{categories-lemma-left-localization-limits}.
Cela prouve (1).

\medskip\noindent
La partie (2) est duale de la partie (1).

\medskip\noindent
Si $S$ est un système multiplicatif, alors $S$ est à la fois un système multiplicatif
à gauche et à droite. Ainsi, $S^{-1}\mathcal{A}$ possède des
noyaux et des conoyaux, et $Q$ commute avec les noyaux et les conoyaux.
Pour achever la preuve de (3), il faut montrer que $\Coim = \Im$ dans
$S^{-1}\mathcal{A}$. En utilisant de nouveau le fait que toute flèche de $S^{-1}\mathcal{A}$
est isomorphe à une flèche $Q(f)$, on voit que le résultat découle
du résultat correspondant pour $\mathcal{A}$.
\end{proof}




\section{Jordan-H\"older}
\label{section-jordan-holder}

\noindent
Le lemme de Jordan-H\"older est le lemme \ref{lemma-jordan-holder}.
Commençons par quelques définitions.

\begin{definition}
\label{definition-simple}
Soit $\mathcal{A}$ une catégorie abélienne. Un objet $A$ de $\mathcal{A}$
est dit {\it simple} s'il est non nul et si les seuls sous-objets
de $A$ sont $0$ et $A$.
\end{definition}

\begin{definition}
\label{definition-Artinian}
Soit $\mathcal{A}$ une catégorie abélienne.
\begin{enumerate}
\item On dit qu'un objet $A$ de $\mathcal{A}$ est {\it artinien} si et seulement s'il
satisfait la condition de chaîne descendante pour les sous-objets.
\item On dit que $\mathcal{A}$ est {\it artinienne} si tout objet de
$\mathcal{A}$ est artinien.
\end{enumerate}
\end{definition}

\begin{definition}
\label{definition-Noetherian}
Soit $\mathcal{A}$ une catégorie abélienne.
\begin{enumerate}
\item On dit qu'un objet $A$ de $\mathcal{A}$ est {\it noethérien} si et seulement s'il
satisfait la condition de chaîne ascendante pour les sous-objets.
\item On dit que $\mathcal{A}$ est {\it noethérienne} si tout objet de
$\mathcal{A}$ est noethérien.
\end{enumerate}
\end{definition}

\begin{lemma}
\label{lemma-ses-artinian}
Soit $\mathcal{A}$ une catégorie abélienne. Soit $0 \to A_1 \to A_2 \to A_3 \to 0$
une suite exacte courte de $\mathcal{A}$. Alors $A_2$ est artinien
si et seulement si $A_1$ et $A_3$ sont artiniens.
\end{lemma}

\begin{proof}
Omis.
\end{proof}

\begin{lemma}
\label{lemma-ses-noetherian}
Soit $\mathcal{A}$ une catégorie abélienne. Soit $0 \to A_1 \to A_2 \to A_3 \to 0$
une suite exacte courte de $\mathcal{A}$. Alors $A_2$ est noethérien
si et seulement si $A_1$ et $A_3$ sont noethériens.
\end{lemma}

\begin{proof}
Omis.
\end{proof}

\begin{lemma}
\label{lemma-finite-length}
Soit $\mathcal{A}$ une catégorie abélienne. Soit $A$ un objet
de $\mathcal{A}$. Les assertions suivantes sont équivalentes :
\begin{enumerate}
\item $A$ est artinien et noethérien, et
\item il existe une filtration
$0 \subset A_1 \subset A_2 \subset \ldots \subset A_n = A$
par des sous-objets telle que $A_i/A_{i - 1}$ soit simple pour $i = 1, \ldots, n$.
\end{enumerate}
\end{lemma}

\begin{proof}
Supposons (1). Si $A$ est nul, alors (2) est vrai. Si $A$ n'est pas nul, la propriété
artinienne fournit un plus petit objet non nul $A_1 \subset A$.
Bien entendu, $A_1$ est simple. Si $A_1 = A$, la preuve est terminée.
Sinon, on peut trouver $A_1 \subset A_2 \subset A$, minimal
sous la condition $A_2 \not = A_1$. Alors $A_2/A_1$ est simple. En poursuivant ainsi, on
peut trouver une suite $0 \subset A_1 \subset A_2 \subset \ldots$
de sous-objets de $A$ telle que $A_i/A_{i - 1}$ soit simple. Puisque $A$
est noethérien, ce processus s'arrête. On en déduit (2).

\medskip\noindent
Supposons (2). Nous allons prouver (1) par récurrence sur $n$. Si $n = 1$, alors
$A$ est simple, et donc clairement noethérien et artinien. Si le résultat est vrai
pour $n - 1$, on utilise la suite exacte courte
$0 \to A_{n - 1} \to A_n \to A_n/A_{n - 1} \to 0$
et les lemmes \ref{lemma-ses-artinian} et \ref{lemma-ses-noetherian}
pour conclure au rang $n$.
\end{proof}

\begin{lemma}[Jordan-H\"older]
\label{lemma-jordan-holder}
Soit $\mathcal{A}$ une catégorie abélienne. Soit $A$ un objet
de $\mathcal{A}$ satisfaisant les conditions équivalentes du
lemme \ref{lemma-finite-length}. Étant données deux filtrations
$$
0 \subset A_1 \subset A_2 \subset \ldots \subset A_n = A
\quad\text{et}\quad
0 \subset B_1 \subset B_2 \subset \ldots \subset B_m = A
$$
dont les $S_i = A_i/A_{i - 1}$ et $T_j = B_j/B_{j - 1}$ sont des objets simples, on a
$n = m$, et il existe une permutation $\sigma$ de $\{1, \ldots, n\}$
telle que $S_i \cong T_{\sigma(i)}$ pour tout $i \in \{1, \ldots, n\}$.
\end{lemma}

\begin{proof}
Soit $j$ le plus petit indice tel que $A_1 \subset B_j$.
Alors l'application $S_1 = A_1 \to B_j/B_{j - 1} = T_j$ est un isomorphisme.
De plus, l'objet $A/A_1 = A_n/A_1 = B_m/A_1$
possède les deux filtrations
$$
0 \subset A_2/A_1 \subset A_3/A_1 \subset \ldots \subset A_n/A_1
$$
et
$$
0 \subset (B_1 + A_1)/A_1 \subset \ldots \subset
(B_{j - 1} + A_1)/A_1 = B_j/A_1 \subset B_{j + 1}/A_1
\subset \ldots \subset B_m/A_1
$$
On conclut par récurrence.
\end{proof}






\section{Sous-catégories de Serre}
\label{section-serre-subcategories}

\noindent
Dans \cite[chapitre I, section 1]{Serre_homotopie_classes},
une notion de « classe » de groupes abéliens est définie.
Cette notion a été étendue aux catégories abéliennes par de nombreux auteurs
(de manières légèrement différentes). Nous utiliserons la variante suivante,
qui est pratiquement identique à la définition originale de Serre.

\begin{definition}
\label{definition-serre-subcategory}
\begin{reference}
\cite[Condition (I) on page 259]{Serre_homotopie_classes}
\end{reference}
Soit $\mathcal{A}$ une catégorie abélienne.
\begin{enumerate}
\item Une {\it sous-catégorie de Serre} de $\mathcal{A}$ est une
sous-catégorie pleine non vide $\mathcal{C}$ de $\mathcal{A}$
telle que, pour toute suite exacte\footnote{D'après la
définition \ref{definition-exact}, cela signifie que $\Im(A \to B) = \Ker(B \to C)$.}
$$
A \to B \to C
$$
avec $A, C \in \Ob(\mathcal{C})$, on ait aussi
$B \in \Ob(\mathcal{C})$.
\item Une {\it sous-catégorie de Serre faible} de $\mathcal{A}$ est une
sous-catégorie pleine non vide $\mathcal{C}$ de $\mathcal{A}$ telle que, pour toute
suite exacte
$$
A_0 \to A_1 \to A_2 \to A_3 \to A_4
$$
avec $A_0, A_1, A_3, A_4$ dans $\mathcal{C}$, on ait aussi $A_2$ dans $\mathcal{C}$.
\end{enumerate}
\end{definition}

\noindent
Dans certaines références, la seconde notion est appelée sous-catégorie « épaisse »,
et dans d'autres, c'est la première notion qui porte ce nom.
Il semble toutefois universellement admis que la notion de sous-catégorie de Serre
est celle définie ci-dessus. Remarquons que, dans les deux cas, la catégorie
$\mathcal{C}$ est abélienne et que le foncteur d'inclusion
$\mathcal{C} \to \mathcal{A}$ est un foncteur exact pleinement fidèle.
Caractérisons plus précisément ces types de sous-catégories.

\begin{lemma}
\label{lemma-characterize-serre-subcategory}
Soit $\mathcal{A}$ une catégorie abélienne.
Soit $\mathcal{C}$ une sous-catégorie de $\mathcal{A}$.
Alors $\mathcal{C}$ est une sous-catégorie de Serre si et seulement si
les conditions suivantes sont satisfaites :
\begin{enumerate}
\item $0 \in \Ob(\mathcal{C})$,
\item $\mathcal{C}$ est une sous-catégorie strictement pleine de $\mathcal{A}$,
\item tout sous-objet ou quotient d'un objet de $\mathcal{C}$ est un objet
de $\mathcal{C}$,
\item si $A \in \Ob(\mathcal{A})$ est une extension d'objets de $\mathcal{C}$,
alors $A \in \Ob(\mathcal{C})$ également.
\end{enumerate}
De plus, une sous-catégorie de Serre est une catégorie abélienne et
le foncteur d'inclusion est exact.
\end{lemma}

\begin{proof}
Omis.
\end{proof}

\begin{lemma}
\label{lemma-characterize-weak-serre-subcategory}
Soit $\mathcal{A}$ une catégorie abélienne.
Soit $\mathcal{C}$ une sous-catégorie de $\mathcal{A}$.
Alors $\mathcal{C}$ est une sous-catégorie de Serre faible si et seulement si
les conditions suivantes sont satisfaites :
\begin{enumerate}
\item $0 \in \Ob(\mathcal{C})$,
\item $\mathcal{C}$ est une sous-catégorie strictement pleine de $\mathcal{A}$,
\item les noyaux et conoyaux dans $\mathcal{A}$ des morphismes
entre objets de $\mathcal{C}$ appartiennent à $\mathcal{C}$,
\item si $A \in \Ob(\mathcal{A})$ est une extension d'objets de $\mathcal{C}$,
alors $A \in \Ob(\mathcal{C})$ également.
\end{enumerate}
De plus, une sous-catégorie de Serre faible est une catégorie abélienne et
le foncteur d'inclusion est exact.
\end{lemma}

\begin{proof}
Omis.
\end{proof}

\begin{lemma}
\label{lemma-kernel-exact-functor}
Soient $\mathcal{A}$, $\mathcal{B}$ des catégories abéliennes.
Soit $F : \mathcal{A} \to \mathcal{B}$ un foncteur exact.
Alors la sous-catégorie pleine des objets $C$ de $\mathcal{A}$
tels que $F(C) = 0$ forme une sous-catégorie de Serre de $\mathcal{A}$.
\end{lemma}

\begin{proof}
Omis.
\end{proof}

\begin{definition}
\label{definition-kernel-category}
Soient $\mathcal{A}$, $\mathcal{B}$ des catégories abéliennes.
Soit $F : \mathcal{A} \to \mathcal{B}$ un foncteur exact.
La sous-catégorie pleine des objets $C$ de $\mathcal{A}$
tels que $F(C) = 0$ est alors appelée le {\it noyau du foncteur $F$},
et est parfois notée $\Ker(F)$.
\end{definition}

\noindent
Toute sous-catégorie de Serre d'une catégorie abélienne est le noyau
d'un foncteur exact. Dans
Exemples, section \ref{examples-section-serre-quotient-modulo-torsion-modules},
nous étudions ce fait pour l'exemple original de Serre des groupes de torsion.

\begin{lemma}
\label{lemma-serre-subcategory-is-kernel}
Soit $\mathcal{A}$ une catégorie abélienne.
Soit $\mathcal{C} \subset \mathcal{A}$ une sous-catégorie de Serre.
Il existe une catégorie abélienne $\mathcal{A}/\mathcal{C}$
et un foncteur exact
$$
F : \mathcal{A} \longrightarrow \mathcal{A}/\mathcal{C}
$$
qui est essentiellement surjectif et dont le noyau est $\mathcal{C}$.
La catégorie $\mathcal{A}/\mathcal{C}$ et le foncteur $F$ sont
caractérisés par la propriété universelle suivante : pour tout foncteur exact
$G : \mathcal{A} \to \mathcal{B}$ tel que
$\mathcal{C} \subset \Ker(G)$, il existe une factorisation
$G = H \circ F$ par un unique foncteur exact
$H : \mathcal{A}/\mathcal{C} \to \mathcal{B}$.
\end{lemma}

\begin{proof}
Considérons l'ensemble de flèches de $\mathcal{A}$ défini par
la formule suivante :
$$
S = \{f \in \text{Flèches}(\mathcal{A}) \mid
\Ker(f), \Coker(f) \in \Ob(\mathcal{C}) \}.
$$
Nous affirmons que $S$ est un système multiplicatif. Pour le prouver, il faut
vérifier MS1, MS2 et MS3, voir
Catégories, définition \ref{categories-definition-multiplicative-system}.

\medskip\noindent
Il est clair que les identités appartiennent à $S$. Supposons que
$f : A \to B$ et $g : B \to C$ appartiennent à $S$.
Il existe des suites exactes
$$
\begin{matrix}
0 \to \Ker(f) \to \Ker(gf) \to \Ker(g) \\
\Coker(f) \to \Coker(gf) \to \Coker(g) \to 0
\end{matrix}
$$
Il s'ensuit que $gf \in S$. Cela prouve MS1. (En fait, un argument
semblable montrera que $S$ est un système multiplicatif saturé, voir
Catégories, définition
\ref{categories-definition-saturated-multiplicative-system}.)

\medskip\noindent
Considérons un diagramme plein
$$
\xymatrix{
A \ar[d]_t \ar[r]_g & B \ar@{..>}[d]^s \\
C \ar@{..>}[r]^f & C \amalg_A B
}
$$
avec $t \in S$. Posons
$W = C \amalg_A B =  \Coker((t, -g) : A \to C \oplus B)$.
Alors $\Ker(t) \to \Ker(s)$ est surjectif et
$\Coker(t) \to \Coker(s)$ est un isomorphisme. Par conséquent,
$s$ appartient à $S$. Cela prouve LMS2 ; la preuve de RMS2 est duale.

\medskip\noindent
Enfin, considérons des morphismes $f, g : B \to C$ et un morphisme $s : A \to B$
dans $S$ tels que $f \circ s = g \circ s$. Cela signifie que
$(f - g) \circ s = 0$. À son tour, cela signifie que
$I = \Im(f - g) \subset C$ est un quotient de $\Coker(s)$,
donc un objet de $\mathcal{C}$. Ainsi, $t : C \to C' = C/I$ est un
élément de $S$ tel que $t \circ (f - g) = 0$, c'est-à-dire tel que
$t \circ f = t \circ g$. Cela prouve LMS3 ; la preuve de
RMS3 est duale.

\medskip\noindent
Ayant montré que $S$ est un système multiplicatif, on pose
$\mathcal{A}/\mathcal{C} = S^{-1}\mathcal{A}$ et l'on prend pour
$F$ le foncteur de localisation $Q$. D'après le
lemme \ref{lemma-localization-abelian},
la catégorie $\mathcal{A}/\mathcal{C}$ est abélienne et $F$ est exact.
Si $X$ appartient au noyau de $F = Q$, alors, d'après le
lemme \ref{lemma-kernel-localization},
on voit que $0 : X \to Z$ appartient à $S$, et par conséquent
$X$ est un objet de $\mathcal{C}$, c'est-à-dire que le noyau de
$F$ est $\mathcal{C}$.
Enfin, si $G$ est comme dans l'énoncé du lemme, alors $G$ transforme
chaque élément de $S$ en un isomorphisme. La propriété universelle de la localisation
fournit donc le foncteur $H : \mathcal{A}/\mathcal{C} \to \mathcal{B}$, voir
Catégories, lemme \ref{categories-lemma-properties-left-localization}.
Il reste à montrer que le foncteur $H$ est exact.
Il suffit pour cela de montrer que $H$ commute
à la prise de noyaux et de conoyaux, voir le lemme \ref{lemma-exact-functor}.
Soit $A \to B$ un morphisme de $\mathcal{A}/\mathcal{C}$.
On peut représenter $A \to B$ sous la forme $fs^{-1}$, où $s : A' \to A$
appartient à $S$ et $f : A' \to B$ est un morphisme arbitraire de $\mathcal{A}$.
Puisque $F = Q$ envoie $s$ sur un isomorphisme dans la catégorie quotient
$\mathcal{A}/\mathcal{C}$, il suffit de montrer que $H$ commute à la prise des
noyaux et conoyaux des morphismes $f : A \to B$ de $\mathcal{A}$.
Mais ici, on a $H(f) = G(f)$, et le résultat découle
du fait que $G$ est exact.
\end{proof}

\begin{lemma}
\label{lemma-quotient-by-kernel-exact-functor}
Soient $\mathcal{A}$, $\mathcal{B}$ des catégories abéliennes.
Soit $F : \mathcal{A} \to \mathcal{B}$ un foncteur exact.
Soit $\mathcal{C} \subset \mathcal{A}$ une sous-catégorie de Serre
contenue dans le noyau de $F$.
Alors $\mathcal{C} = \Ker(F)$ si et seulement si le foncteur induit
$\overline{F} : \mathcal{A}/\mathcal{C} \to \mathcal{B}$
(lemme \ref{lemma-serre-subcategory-is-kernel}) est fidèle.
\end{lemma}

\begin{proof}
Nous utiliserons sans plus de mention les résultats du lemme
\ref{lemma-serre-subcategory-is-kernel}.
L'implication directe est vraie car le noyau de $\overline{F}$ est nul
par construction. En effet, si $f : X \to Y$ est un morphisme de
$\mathcal{A}/\mathcal{C}$ tel que $\overline{F}(f) = 0$, alors
$\overline{F}(\Im(f)) = \Im(\overline{F}(f)) = 0$ ; ainsi, $\Im(f) = 0$ par
l'hypothèse sur le noyau de $F$. Donc $f = 0$.

\medskip\noindent
Pour l'implication réciproque, soit $X$ un objet de $\mathcal{A}$ tel que $F(X)
= 0$. Alors $\overline{F}(\text{id}_X) = \text{id}_{\overline{F}(X)} = 0$, donc
$\text{id}_X = 0$ dans $\mathcal{A}/\mathcal{C}$ par fidélité de
$\overline{F}$. Par conséquent, $X = 0$ dans $\mathcal{A}/\mathcal{C}$, c'est-à-dire $X \in
\Ob(\mathcal{C})$.
\end{proof}






\section{Groupes K}
\label{section-K-groups}

\noindent
Quelques mots sur le groupe $K_0$ d'une catégorie abélienne.

\begin{definition}
\label{definition-K-zero}
Soit $\mathcal{A}$ une catégorie abélienne.
On note $K_0(\mathcal{A})$ le
{\it groupe $K$ d'indice zéro de $\mathcal{A}$}.
C'est le groupe abélien construit comme suit.
On prend le groupe abélien libre
sur les objets de $\mathcal{A}$
et, pour toute suite exacte courte
$0 \to A \to B \to C \to 0$,
on impose la relation $[B] - [A] - [C] = 0$.
\end{definition}

\noindent
Autrement dit, il existe une présentation
$$
\bigoplus_{A \to B \to C\text{ sec}}
\mathbf{Z}[A \to B \to C]
\longrightarrow
\bigoplus_{A \in \Ob(\mathcal{A})}
\mathbf{Z}[A]
\longrightarrow
K_0(\mathcal{A})
\longrightarrow
0
$$
où la première flèche envoie $[A \to B \to C]$ sur $[B] - [A] - [C]$.
La suite exacte courte $0 \to 0 \to 0 \to 0 \to 0$
donne la relation $[0] = 0$ dans $K_0(\mathcal{A})$.
Il n'y a pas de difficulté ensembliste, car toutes nos catégories
sont « petites », sauf indication contraire.
Quelques exemples de groupes $K$ pour des catégories de modules
sur des anneaux ont été calculés dans
Algèbre, section \ref{algebra-section-K-groups}.

\begin{lemma}
\label{lemma-exact-functor-K-groups}
Soit $F : \mathcal{A} \to \mathcal{B}$ un foncteur exact entre
catégories abéliennes. Alors $F$ induit un homomorphisme de groupes $K$
$K_0(F) : K_0(\mathcal{A}) \to K_0(\mathcal{B})$, simplement défini par
$K_0(F)([A]) = [F(A)]$.
\end{lemma}

\begin{proof}
La preuve est immédiate.
\end{proof}

\noindent
Supposons donnés un objet $M$ d'une catégorie abélienne $\mathcal{A}$
et un complexe de la forme
\begin{equation}
\label{equation-cyclic-complex}
\xymatrix{
\ldots \ar[r] &
M \ar[r]^\varphi &
M \ar[r]^\psi &
M \ar[r]^\varphi &
M \ar[r] & \ldots
}
\end{equation}
Dans cette situation, on définit
$$
H^0(M, \varphi, \psi) = \Ker(\psi)/\Im(\varphi)
, \quad\text{et}\quad
H^1(M, \varphi, \psi) = \Ker(\varphi)/\Im(\psi).
$$

\begin{lemma}
\label{lemma-serre-subcategory-K-groups}
Soit $\mathcal{A}$ une catégorie abélienne.
Soit $\mathcal{C} \subset \mathcal{A}$ une sous-catégorie de Serre et
posons $\mathcal{B} = \mathcal{A}/\mathcal{C}$.
\begin{enumerate}
\item Les foncteurs exacts $\mathcal{C} \to \mathcal{A}$ et
$\mathcal{A} \to \mathcal{B}$ induisent une suite exacte
$$
K_0(\mathcal{C}) \to
K_0(\mathcal{A}) \to
K_0(\mathcal{B}) \to
0
$$
de groupes $K$, et
\item le noyau de $K_0(\mathcal{C}) \to K_0(\mathcal{A})$ est égal
à l'ensemble des éléments de la forme
$$
[H^0(M, \varphi, \psi)] - [H^1(M, \varphi, \psi)]
$$
où $(M, \varphi, \psi)$ est un complexe comme dans (\ref{equation-cyclic-complex})
qui devient exact dans $\mathcal{B}$ ; autrement dit,
$H^0(M, \varphi, \psi)$ et $H^1(M, \varphi, \psi)$ sont des
objets de $\mathcal{C}$.
\end{enumerate}
\end{lemma}

\begin{proof}
Prouvons (1). Il est clair que $K_0(\mathcal{A}) \to K_0(\mathcal{B})$
est surjectif et que la composition $K_0(\mathcal{C}) \to
K_0(\mathcal{A}) \to K_0(\mathcal{B})$ est nulle. Soit $x \in K_0(\mathcal{A})$
un élément envoyé sur zéro dans $K_0(\mathcal{B})$. On peut écrire
$x = [A] - [A']$, avec $A, A'$ dans $\mathcal{A}$ (exercice amusant).
Notons $B, B'$ les objets correspondants de $\mathcal{B}$. Le fait que
$x$ soit envoyé sur zéro dans $K_0(\mathcal{B})$
signifie qu'il existe un ensemble fini $I = I^+ \amalg I^{-}$,
et, pour chaque $i \in I$, une suite exacte courte
$$
0 \to B_i \to B'_i \to B''_i \to 0
$$
dans $\mathcal{B}$, tels que l'on ait
$$
[B] - [B'] = \sum\nolimits_{i \in I^{+}} ([B'_i] - [B_i] - [B''_i])
-
\sum\nolimits_{i \in I^{-}} ([B'_i] - [B_i] - [B''_i])
$$
dans le groupe abélien libre sur les classes d'isomorphisme d'objets de $\mathcal{B}$.
On peut réécrire cette égalité sous la forme
$$
[B]
+ \sum\nolimits_{i \in I^{+}} ([B_i] + [B''_i])
+ \sum\nolimits_{i \in I^{-}} [B'_i]
=
[B']
+ \sum\nolimits_{i \in I^{-}} ([B_i] + [B''_i])
+ \sum\nolimits_{i \in I^{+}} [B'_i].
$$
Comme les membres de droite et de gauche doivent contenir les mêmes classes d'isomorphisme
d'objets de $\mathcal{B}$, comptées avec multiplicité, il doit donc
exister une bijection
$$
\tau :
\{B\} \amalg \{B_i, B''_i; i \in I^+\} \amalg \{B'_i; i \in I^-\}
\longrightarrow
\{B'\} \amalg \{B_i, B''_i; i \in I^-\} \amalg \{B'_i; i \in I^+\}
$$
telle que $N$ et $\tau(N)$ soient isomorphes dans $\mathcal{B}$.
Les preuves des lemmes \ref{lemma-serre-subcategory-is-kernel} et
\ref{lemma-localization-abelian} montrent que l'on peut choisir, pour $i \in I$,
une suite exacte courte
$$
0 \to A_i \to A'_i \to A''_i \to 0
$$
dans $\mathcal{A}$ telle que $B_i, B'_i, B''_i$ soient isomorphes aux
images de $A_i, A'_i, A''_i$ dans $\mathcal{B}$. Cela implique que la
bijection correspondante
$$
\tau :
\{A\} \amalg \{A_i, A''_i; i \in I^+\} \amalg \{A'_i; i \in I^-\}
\longrightarrow
\{A'\} \amalg \{A_i, A''_i; i \in I^-\} \amalg \{A'_i; i \in I^+\}
$$
vérifie la propriété suivante : $M$ et $\tau(M)$ sont des objets de $\mathcal{A}$
qui deviennent isomorphes dans $\mathcal{B}$. Cela signifie que $[M] - [\tau(M)]$
appartient à l'image de $K_0(\mathcal{C}) \to K_0(\mathcal{A})$. En effet,
l'isomorphisme dans $\mathcal{B}$ est donné par un diagramme
$M \leftarrow M' \rightarrow \tau(M)$ dans $\mathcal{A}$, où
$M' \to M$ et $M' \to \tau(M)$ ont tous deux noyau et conoyau dans $\mathcal{C}$.
En remontant les étapes, on conclut que $x = [A] - [A']$ appartient à l'image
de $K_0(\mathcal{C}) \to K_0(\mathcal{A})$, ce qui achève la preuve de (1).

\medskip\noindent
Prouvons (2). La preuve est analogue à celle de (1), mais demande un peu
plus de tenue des comptes. Remarquons d'abord que toute classe de la forme
$[H^0(M, \varphi, \psi)] - [H^1(M, \varphi, \psi)]$ est nulle
dans $K_0(\mathcal{A})$, comme le montre le calcul suivant :
\begin{align*}
0 & = [M] - [M] \\
& =  [\Ker(\varphi)] + [\Im(\varphi)]
- [\Ker(\psi)] - [\Im(\psi)] \\
& =
[\Ker(\varphi)/\Im(\psi)] -
[\Ker(\psi)/\Im(\varphi)] \\
& = [H^1(M, \varphi, \psi)] - [H^0(M, \varphi, \psi)]
\end{align*}
comme souhaité. Il suffit donc de montrer que tout élément du noyau
de $K_0(\mathcal{C}) \to K_0(\mathcal{A})$ est de cette forme.

\medskip\noindent
Tout élément $x$ de $K_0(\mathcal{C})$ peut être représenté comme la
différence $x = [P] - [Q]$ de deux objets de $\mathcal{C}$ (exercice amusant).
Supposons que cet élément soit envoyé sur zéro dans $K_0(\mathcal{A})$.
Cela signifie qu'il existe
\begin{enumerate}
\item un ensemble fini $I = I^{+} \amalg I^{-}$,
\item pour $i \in I$, une suite exacte courte $0 \to A_i \to B_i \to C_i \to 0$
dans $\mathcal{A}$,
\end{enumerate}
tels que
$$
[P] - [Q] =
\sum\nolimits_{i \in I^{+}} ([B_i] - [A_i] - [C_i])
-
\sum\nolimits_{i \in I^{-}} ([B_i] - [A_i] - [C_i])
$$
dans le groupe abélien libre sur les objets de $\mathcal{A}$.
On peut réécrire cette égalité sous la forme
$$
[P]
+ \sum\nolimits_{i \in I^{+}} ([A_i] + [C_i])
+ \sum\nolimits_{i \in I^{-}} [B_i]
=
[Q]
+ \sum\nolimits_{i \in I^{-}} ([A_i] + [C_i])
+ \sum\nolimits_{i \in I^{+}} [B_i].
$$
Comme les membres de droite et de gauche doivent contenir les mêmes objets
de $\mathcal{A}$, comptés avec multiplicité, il doit donc exister
une bijection $\tau$ entre les termes qui figurent ci-dessus. Posons
$$
T^{+} =
\{p\}\ \amalg\ \{a, c\} \times I^{+}\ \amalg\ \{b\} \times I^{-}
$$
et
$$
T^{-} =
\{q\}\ \amalg\ \{a, c\} \times I^{-}\ \amalg\ \{b\} \times I^{+}.
$$
Posons $T = T^{+} \amalg T^{-} = \{p, q\} \amalg \{a, b, c\} \times I$.
Pour $t \in T$, définissons
$$
O(t)
=
\left\{
\begin{matrix}
P & \text{si} & t = p \\
Q & \text{si} & t = q \\
A_i & \text{si} & t = (a, i) \\
B_i & \text{si} & t = (b, i) \\
C_i & \text{si} & t = (c, i)
\end{matrix}
\right.
$$
On peut ainsi considérer $\tau : T^{+} \to T^{-}$ comme une bijection
telle que $O(t) = O(\tau(t))$ pour tout $t \in T^{+}$.
Soit $t^{-}_0 = \tau(p)$ et soit $t^{+}_0 \in T^{+}$ l'unique
élément tel que $\tau(t^{+}_0) = q$.
Considérons l'objet
$$
M^{+} = \bigoplus\nolimits_{t \in T^{+}} O(t)
$$
En utilisant $\tau$, on voit qu'il est égal à l'objet
$$
M^{-} = \bigoplus\nolimits_{t \in T^{-}} O(t)
$$
Considérons l'application
$$
\varphi : M^{+} \longrightarrow M^{-}
$$
qui, sur le facteur $O(t) = A_i$ correspondant à $t = (a, i)$, $i \in I^{+}$,
utilise l'application $A_i \to B_i$ à valeurs dans le facteur $O((b, i)) = B_i$ de $M^{-}$,
et qui, sur le facteur $O(t) = B_i$ correspondant à $(b, i)$, $i \in I^{-}$,
utilise l'application $B_i \to C_i$ à valeurs dans le facteur $O((c, i)) = C_i$ de $M^{-}$.
L'application est nulle sur les facteurs correspondant à $p$
et à $(c, i)$, $i \in I^{+}$.
De même, considérons l'application
$$
\psi : M^{-} \longrightarrow M^{+}
$$
qui, sur le facteur $O(t) = A_i$ correspondant à $t = (a, i)$, $i \in I^{-}$,
utilise l'application $A_i \to B_i$ à valeurs dans le facteur $O((b, i)) = B_i$ de $M^{+}$,
et qui, sur le facteur $O(t) = B_i$ correspondant à $(b, i)$, $i \in I^{+}$,
utilise l'application $B_i \to C_i$ à valeurs dans le facteur $O((c, i)) = C_i$ de $M^{+}$.
L'application est nulle sur les facteurs correspondant à $q$ et à
$(c, i)$, $i \in I^{-}$.

\medskip\noindent
Remarquons que le noyau de $\varphi$ est égal à la somme directe du
facteur $P$, des facteurs $O((c, i)) = C_i$, $i \in I^{+}$, et
des sous-objets $A_i$ à l'intérieur des facteurs $O((b, i)) = B_i$, $i \in I^{-}$.
L'image de $\psi$ est égale à la somme directe des
facteurs $O((c, i)) = C_i$, $i \in I^{+}$, et
des sous-objets $A_i$ à l'intérieur des facteurs $O((b, i)) = B_i$, $i \in I^{-}$.
Autrement dit, on voit que
$$
P \cong \Ker(\varphi)/\Im(\psi).
$$
De manière exactement analogue, on voit que
$$
Q \cong \Ker(\psi)/\Im(\varphi).
$$
Puisque, comme nous l'avons remarqué ci-dessus, l'existence de la bijection
$\tau$ montre que $M^{+} = M^{-}$, le lemme en découle.
\end{proof}



\section{Foncteurs cohomologiques}
\label{section-cohomological-delta-functor}

\begin{definition}
\label{definition-cohomological-delta-functor}
Soient $\mathcal{A}, \mathcal{B}$ des catégories abéliennes.
Un {\it foncteur cohomologique $\delta$}, ou simplement un
{\it foncteur $\delta$}, de $\mathcal{A}$
vers $\mathcal{B}$ est constitué des données suivantes :
\begin{enumerate}
\item une famille $F^n : \mathcal{A} \to \mathcal{B}$, $n \geq 0$, de foncteurs
additifs, et
\item pour toute suite exacte courte $0 \to A \to B \to C \to 0$
de $\mathcal{A}$,
une famille $\delta_{A \to B \to C} : F^n(C) \to F^{n + 1}(A)$, $n \geq 0$,
de morphismes de $\mathcal{B}$.
\end{enumerate}
On suppose que ces données satisfont aux axiomes suivants :
\begin{enumerate}
\item pour toute suite exacte courte comme ci-dessus, la suite
$$
\xymatrix{
0 \ar[r] &
F^0(A) \ar[r] &
F^0(B) \ar[r] &
F^0(C) \ar[lld]^{\delta_{A \to B \to C}} \\
 &
F^1(A) \ar[r] &
F^1(B) \ar[r] &
F^1(C) \ar[lld]^{\delta_{A \to B \to C}} \\
 &
F^2(A) \ar[r] &
F^2(B) \ar[r] &
\ldots
}
$$
est exacte, et
\item pour tout morphisme $(A \to B \to C) \to (A' \to B' \to C')$
de suites exactes courtes de $\mathcal{A}$, les diagrammes
$$
\xymatrix{
F^n(C) \ar[d] \ar[rr]_{\delta_{A \to B \to C}} & & F^{n + 1}(A) \ar[d] \\
F^n(C') \ar[rr]^{\delta_{A' \to B' \to C'}} & & F^{n + 1}(A')
}
$$
sont commutatifs.
\end{enumerate}
\end{definition}

\noindent
Notons que ceci implique en particulier que $F^0$ est exact à gauche.

\begin{definition}
\label{definition-morphism-delta-functors}
Soient $\mathcal{A}, \mathcal{B}$ des catégories abéliennes.
Soient $(F^n, \delta_F)$ et $(G^n, \delta_G)$ des foncteurs $\delta$
de $\mathcal{A}$ vers $\mathcal{B}$. Un {\it morphisme de foncteurs $\delta$
de $F$ vers $G$} est une famille de
transformations de foncteurs $t^n : F^n \to G^n$, $n \geq 0$, telle
que, pour toute suite exacte courte $0 \to A \to B \to C \to 0$
de $\mathcal{A}$, les diagrammes
$$
\xymatrix{
F^n(C) \ar[d]_{t^n} \ar[rr]_{\delta_{F, A \to B \to C}} &
& F^{n + 1}(A) \ar[d]^{t^{n + 1}} \\
G^n(C) \ar[rr]^{\delta_{G, A \to B \to C}} & & G^{n + 1}(A)
}
$$
sont commutatifs.
\end{definition}

\begin{definition}
\label{definition-universal-delta-functor}
Soient $\mathcal{A}, \mathcal{B}$ des catégories abéliennes.
Soit $F = (F^n, \delta_F)$ un foncteur $\delta$
de $\mathcal{A}$ vers $\mathcal{B}$.
On dit que $F$ est un {\it foncteur $\delta$ universel} si et seulement
si, pour tout foncteur $\delta$ $G = (G^n, \delta_G)$ et tout
morphisme de foncteurs $t : F^0 \to G^0$, il existe
un unique morphisme de foncteurs $\delta$ $\{t^n\}_{n \geq 0} : F \to G$
tel que $t = t^0$.
\end{definition}

\begin{lemma}
\label{lemma-efface-implies-universal}
Soient $\mathcal{A}, \mathcal{B}$ des catégories abéliennes.
Soit $F = (F^n, \delta_F)$ un foncteur $\delta$
de $\mathcal{A}$ vers $\mathcal{B}$.
Supposons que, pour tout $n > 0$ et tout $A \in \Ob(\mathcal{A})$,
il existe un morphisme injectif $u : A \to B$ (dépendant de $A$ et de $n$)
tel que $F^n(u) : F^n(A) \to F^n(B)$ soit nul. Alors $F$ est un foncteur
$\delta$ universel.
\end{lemma}

\begin{proof}
Soit $G = (G^n, \delta_G)$ un foncteur $\delta$
de $\mathcal{A}$ vers $\mathcal{B}$ et soit $t : F^0 \to G^0$
un morphisme de foncteurs. Il faut montrer qu'il existe
un unique morphisme de foncteurs $\delta$ $\{t^n\}_{n \geq 0} : F \to G$
tel que $t = t^0$. Construisons $t^n$ par récurrence sur $n$.
Pour $n = 0$, posons $t^0 = t$.
Supposons déjà construite une unique famille de
transformations de foncteurs $t^i$ pour $i \leq n$, compatible avec
les applications $\delta$ en degrés $\leq n$.

\medskip\noindent
Soit $A \in \Ob(\mathcal{A})$. Par hypothèse, nous pouvons choisir
un plongement $u : A \to B$ tel que $F^{n + 1}(u) = 0$.
Soit $C = B/u(A)$. La suite exacte longue de cohomologie associée à
la suite exacte courte $0 \to A \to B \to C \to 0$ et au
foncteur $\delta$ $F$ donne
$F^{n + 1}(A) = \Coker(F^n(B) \to F^n(C))$, par notre choix de $u$.
Comme $t^n$ est déjà défini, nous pouvons définir
$$
t^{n + 1}_A : F^{n + 1}(A) \to G^{n + 1}(A)
$$
comme l'unique application rendant commutatif le diagramme
$$
\xymatrix{
\Coker(F^n(B) \to F^n(C)) \ar[r]_{t^n}
\ar[d]_{\delta_{F, A \to B \to C}} &
\Coker(G^n(B) \to G^n(C))
\ar[d]^{\delta_{G, A \to B \to C}} \\
F^{n + 1}(A) \ar[r]^{t^{n + 1}_A} &
G^{n + 1}(A)
}
$$
Ceci est manifestement déterminé de manière unique par les conditions
imposées. Nous omettons la vérification que cela définit une transformation
de foncteurs.
\end{proof}

\begin{lemma}
\label{lemma-uniqueness-universal-delta-functor}
Soient $\mathcal{A}, \mathcal{B}$ des catégories abéliennes.
Soit $F : \mathcal{A} \to \mathcal{B}$ un foncteur.
S'il existe un foncteur $\delta$ universel
$(F^n, \delta_F)$ de $\mathcal{A}$ vers $\mathcal{B}$
tel que $F^0 = F$, alors il est déterminé à un unique isomorphisme
de foncteurs $\delta$ près.
\end{lemma}

\begin{proof}
Cela résulte immédiatement des définitions.
\end{proof}





\section{Complexes}
\label{section-complexes}

\noindent
Bien entendu, les notions de complexe de chaînes et de complexe de cochaînes
sont duales et il suffit de lire l'une des deux parties de
cette section. Choisissez donc celle que vous préférez. (En fait, cela ne
fonctionne pas tout à fait, car les conventions de numérotation
ne sont pas adaptées à un passage facile des complexes de chaînes aux complexes de
cochaînes.)

\medskip\noindent
Un {\it complexe de chaînes $A_\bullet$} dans une catégorie additive $\mathcal{A}$
est un complexe
$$
\ldots \to
A_{n + 1} \xrightarrow{d_{n + 1}}
A_n \xrightarrow{d_n}
A_{n - 1} \to
\ldots
$$
de $\mathcal{A}$. Autrement dit, on se donne un objet $A_i$ de
$\mathcal{A}$ pour tout $i \in \mathbf{Z}$ et, pour
tout $i \in \mathbf{Z}$, un morphisme $d_i : A_i \to A_{i - 1}$ tel que
$d_{i - 1} \circ d_i = 0$ pour tout $i$. Un {\it morphisme de complexes
de chaînes $f : A_\bullet \to B_\bullet$} est constitué d'une
famille de morphismes $f_i : A_i \to B_i$ telle que tous
les diagrammes
$$
\xymatrix{
A_i \ar[r]_{d_i} \ar[d]_{f_i} & A_{i - 1} \ar[d]^{f_{i - 1}} \\
B_i \ar[r]^{d_i} & B_{i - 1}
}
$$
soient commutatifs. La {\it catégorie des complexes de chaînes de $\mathcal{A}$}
est notée $\text{Ch}(\mathcal{A})$. La sous-catégorie pleine formée
des objets de la forme
$$
\ldots \to A_2 \to A_1 \to A_0 \to 0 \to 0 \to \ldots
$$
est notée $\text{Ch}_{\geq 0}(\mathcal{A})$.
Autrement dit, un complexe de chaînes $A_\bullet$ appartient à
$\text{Ch}_{\geq 0}(\mathcal{A})$ si et seulement si
$A_i = 0$ pour tout $i < 0$.

\medskip\noindent
Étant donnée une catégorie additive $\mathcal{A}$, on identifie $\mathcal{A}$
à la sous-catégorie pleine de $\text{Ch}(\mathcal{A})$ formée
des complexes de chaînes nuls sauf en degré $0$, au moyen du foncteur
$$
\mathcal{A} \longrightarrow \text{Ch}(\mathcal{A}),\quad
A \longmapsto (\ldots \to 0 \to A \to 0 \to \ldots)
$$
Par abus de notation, on note souvent simplement $A$ l'objet du membre de droite.
Si l'on veut souligner que l'on considère $A$ comme un complexe de chaînes,
on emploiera parfois la notation $A[0]$ ; voir la
section \ref{section-homotopy-shift}.

\medskip\noindent
Une {\it homotopie $h$} entre deux morphismes
de complexes de chaînes $f, g : A_\bullet \to B_\bullet$
est une famille de morphismes $h_i : A_i \to B_{i + 1}$
telle que l'on ait
$$
f_i - g_i = d_{i + 1} \circ h_i + h_{i - 1} \circ d_i
$$
pour tout $i$.
Deux morphismes $f, g : A_\bullet \to B_\bullet$ sont dits
{\it homotopes} s'il existe une homotopie entre $f$
et $g$.
Il est clair que les notions de complexe de chaînes, de morphisme de
complexes de chaînes et d'homotopie entre morphismes de complexes de chaînes
ont un sens même dans une catégorie préadditive.

\begin{lemma}
\label{lemma-compose-homotopy}
\begin{slogan}
Les foncteurs Hom de $\text{Ch}(\mathcal{A})$ respectent la relation d'homotopie.
\end{slogan}
Soit $\mathcal{A}$ une catégorie additive.
Soient $f, g : B_\bullet \to C_\bullet$ des morphismes
de complexes de chaînes. Supposons donnés des morphismes de complexes
de chaînes $a : A_\bullet \to B_\bullet$ et
$c : C_\bullet \to D_\bullet$.
Si $\{h_i : B_i \to C_{i + 1}\}$ définit une homotopie
entre $f$ et $g$, alors $\{c_{i + 1} \circ h_i \circ a_i\}$
définit une homotopie entre $c \circ f \circ a$ et
$c \circ g \circ a$.
\end{lemma}

\begin{proof}
La démonstration est omise.
\end{proof}

\noindent
En particulier, cela signifie qu'il est légitime de définir
la catégorie des complexes de chaînes dont les applications sont prises à homotopie près.
Nous y reviendrons plus loin.

\begin{definition}
\label{definition-homotopy-equivalent}
Soit $\mathcal{A}$ une catégorie additive.
On dit qu'un morphisme $a : A_\bullet \to B_\bullet$
est un {\it homotopisme} s'il existe
un morphisme $b : B_\bullet \to A_\bullet$
tel qu'il existe une homotopie entre
$b \circ a$ et $\text{id}_A$
et une homotopie entre $a \circ b$ et $\text{id}_B$.
S'il existe un tel morphisme entre $A_\bullet$ et $B_\bullet$, on
dit que $A_\bullet$ et $B_\bullet$ sont {\it homotopiquement équivalents}.
\end{definition}

\noindent
Autrement dit, deux complexes sont homotopiquement équivalents s'ils deviennent
isomorphes dans la catégorie des complexes à homotopie près.

\begin{lemma}
\label{lemma-cat-chain-abelian}
Soit $\mathcal{A}$ une catégorie abélienne.
\begin{enumerate}
\item La catégorie des complexes de chaînes dans $\mathcal{A}$ est
abélienne.
\item Un morphisme de complexes
$f : A_\bullet \to B_\bullet$ est injectif
si et seulement si chaque $f_n : A_n \to B_n$ est injectif.
\item Un morphisme de complexes
$f : A_\bullet \to B_\bullet$ est surjectif
si et seulement si chaque $f_n : A_n \to B_n$ est surjectif.
\item Une suite de complexes de chaînes
$$
A_\bullet \xrightarrow{f} B_\bullet \xrightarrow{g} C_\bullet
$$
est exacte en $B_\bullet$ si et seulement si chaque suite
$$
A_i \xrightarrow{f_i} B_i \xrightarrow{g_i} C_i
$$
est exacte en $B_i$.
\end{enumerate}
\end{lemma}

\begin{proof}
La démonstration est omise.
\end{proof}

\noindent
Pour tout $i \in \mathbf{Z}$, le $i$-ième {\it objet d'homologie}
d'un complexe de chaînes $A_\bullet$ dans une catégorie abélienne est défini par
la formule suivante :
$$
H_i(A_\bullet) = \Ker(d_i)/\Im(d_{i + 1}).
$$
Si $f : A_\bullet \to B_\bullet$ est un morphisme de complexes de chaînes
de $\mathcal{A}$, on obtient un morphisme induit
$H_i(f) : H_i(A_\bullet) \to H_i(B_\bullet)$
car, manifestement,
$f_i(\Ker(d_i : A_i \to A_{i - 1})) \subset
\Ker(d_i : B_i \to B_{i - 1})$, et de même
pour $\Im(d_{i + 1})$.
On obtient ainsi un foncteur
$$
H_i : \text{Ch}(\mathcal{A}) \longrightarrow \mathcal{A}.
$$

\begin{definition}
\label{definition-quasi-isomorphism}
Soit $\mathcal{A}$ une catégorie abélienne.
\begin{enumerate}
\item Un morphisme de complexes de chaînes $f : A_\bullet \to B_\bullet$
est appelé un {\it quasi-isomorphisme} si l'application induite
$H_i(f) : H_i(A_\bullet) \to H_i(B_\bullet)$
est un isomorphisme pour tout $i \in \mathbf{Z}$.
\item Un complexe de chaînes $A_\bullet$ est dit
{\it acyclique} si tous ses objets d'homologie
$H_i(A_\bullet)$ sont nuls.
\end{enumerate}
\end{definition}


\begin{lemma}
\label{lemma-map-homology-homotopy}
Soit $\mathcal{A}$ une catégorie abélienne.
\begin{enumerate}
\item Si les applications $f, g : A_\bullet \to B_\bullet$ sont
homotopes, alors les applications induites $H_i(f)$ et $H_i(g)$
sont égales.
\item Si l'application $f : A_\bullet \to B_\bullet$ est un
homotopisme, alors $f$ est un quasi-isomorphisme.
\end{enumerate}
\end{lemma}

\begin{proof}
La démonstration est omise.
\end{proof}

\begin{lemma}
\label{lemma-long-exact-sequence-chain}
Soit $\mathcal{A}$ une catégorie abélienne.
Supposons que
$$
0 \to
A_\bullet \to
B_\bullet \to
C_\bullet \to
0
$$
soit une suite exacte courte de complexes de chaînes de $\mathcal{A}$.
Il existe alors une suite exacte longue canonique d'homologie
$$
\xymatrix{
\ldots & \ldots & \ldots \ar[lld] \\
H_i(A_\bullet) \ar[r] & H_i(B_\bullet) \ar[r] & H_i(C_\bullet) \ar[lld] \\
H_{i - 1}(A_\bullet) \ar[r] &
H_{i - 1}(B_\bullet) \ar[r] &
H_{i - 1}(C_\bullet) \ar[lld] \\
\ldots & \ldots & \ldots \\
}
$$
\end{lemma}

\begin{proof}
La démonstration est omise. Les applications proviennent du lemme du serpent
\ref{lemma-snake} appliqué aux diagrammes
$$
\xymatrix{
&
A_i/\Im(d_{A, i + 1}) \ar[r] \ar[d]^{d_{A, i}} &
B_i/\Im(d_{B, i + 1}) \ar[r] \ar[d]^{d_{B, i}} &
C_i/\Im(d_{C, i + 1}) \ar[r] \ar[d]^{d_{C, i}} &
0 \\
0 \ar[r] &
\Ker(d_{A, i - 1}) \ar[r] &
\Ker(d_{B, i - 1}) \ar[r] &
\Ker(d_{C, i - 1}) &
}
$$
\end{proof}

\noindent
Un {\it complexe de cochaînes $A^\bullet$} dans une catégorie additive $\mathcal{A}$
est un complexe
$$
\ldots \to
A^{n - 1} \xrightarrow{d^{n - 1}}
A^n \xrightarrow{d^n}
A^{n + 1} \to
\ldots
$$
de $\mathcal{A}$. Autrement dit, on se donne un objet $A^i$ de
$\mathcal{A}$ pour tout $i \in \mathbf{Z}$ et, pour
tout $i \in \mathbf{Z}$, un morphisme $d^i : A^i \to A^{i + 1}$ tel que
$d^{i + 1} \circ d^i = 0$ pour tout $i$. Un {\it morphisme de complexes
de cochaînes $f : A^\bullet \to B^\bullet$} est constitué d'une
famille de morphismes $f^i : A^i \to B^i$ telle que tous
les diagrammes
$$
\xymatrix{
A^i \ar[r]_{d^i} \ar[d]_{f^i} & A^{i + 1} \ar[d]^{f^{i + 1}} \\
B^i \ar[r]^{d^i} & B^{i + 1}
}
$$
soient commutatifs. La {\it catégorie des complexes de cochaînes de $\mathcal{A}$}
est notée $\text{CoCh}(\mathcal{A})$. La sous-catégorie pleine formée
des objets de la forme
$$
\ldots \to 0 \to 0 \to A^0 \to A^1 \to A^2 \to \ldots
$$
est notée $\text{CoCh}_{\geq 0}(\mathcal{A})$.
Autrement dit, un complexe de cochaînes $A^\bullet$ appartient à la sous-catégorie
$\text{CoCh}_{\geq 0}(\mathcal{A})$ si et seulement si
$A^i = 0$ pour tout $i < 0$.

\medskip\noindent
Étant donnée une catégorie additive $\mathcal{A}$, on identifie $\mathcal{A}$
à la sous-catégorie pleine de $\text{CoCh}(\mathcal{A})$ formée
des complexes de cochaînes nuls sauf en degré $0$, au moyen du foncteur
$$
\mathcal{A} \longrightarrow \text{CoCh}(\mathcal{A}),\quad
A \longmapsto (\ldots \to 0 \to A \to 0 \to \ldots)
$$
Par abus de notation, on note souvent simplement $A$ l'objet du membre de droite.
Si l'on veut souligner que l'on considère $A$ comme un complexe de cochaînes,
on emploiera parfois la notation $A[0]$ ; voir la
section \ref{section-homotopy-shift}.

\medskip\noindent
Une {\it homotopie $h$} entre deux morphismes
de complexes de cochaînes $f, g : A^\bullet \to B^\bullet$
est une famille de morphismes $h^i : A^i \to B^{i - 1}$
telle que l'on ait
$$
f^i - g^i = d^{i - 1} \circ h^i + h^{i + 1} \circ d^i
$$
pour tout $i$.
Deux morphismes $f, g : A^\bullet \to B^\bullet$ sont dits
{\it homotopes} s'il existe une homotopie entre $f$
et $g$.
Il est clair que les notions de complexe de cochaînes, de morphisme de
complexes de cochaînes et d'homotopie entre morphismes de complexes de cochaînes
ont un sens même dans une catégorie préadditive.

\begin{lemma}
\label{lemma-compose-homotopy-cochain}
Soit $\mathcal{A}$ une catégorie additive.
Soient $f, g : B^\bullet \to C^\bullet$ des morphismes
de complexes de cochaînes. Supposons donnés des morphismes de complexes de
cochaînes $a : A^\bullet \to B^\bullet$ et
$c : C^\bullet \to D^\bullet$.
Si $\{h^i : B^i \to C^{i - 1}\}$ définit une homotopie
entre $f$ et $g$, alors $\{c^{i - 1} \circ h^i \circ a^i\}$
définit une homotopie entre $c \circ f \circ a$ et
$c \circ g \circ a$.
\end{lemma}

\begin{proof}
La démonstration est omise.
\end{proof}

\noindent
En particulier, cela signifie qu'il est légitime de définir
la catégorie des complexes de cochaînes dont les applications sont prises à homotopie près.
Nous y reviendrons plus loin.

\begin{definition}
\label{definition-homotopy-equivalent-cochain}
Soit $\mathcal{A}$ une catégorie additive.
On dit qu'un morphisme $a : A^\bullet \to B^\bullet$
est un {\it homotopisme} s'il existe
un morphisme $b : B^\bullet \to A^\bullet$
tel qu'il existe une homotopie entre
$b \circ a$ et $\text{id}_A$
et une homotopie entre $a \circ b$ et $\text{id}_B$.
S'il existe un tel morphisme entre $A^\bullet$ et $B^\bullet$, on
dit que $A^\bullet$ et $B^\bullet$ sont {\it homotopiquement équivalents}.
\end{definition}

\noindent
Autrement dit, deux complexes sont homotopiquement équivalents s'ils deviennent
isomorphes dans la catégorie des complexes à homotopie près.

\begin{lemma}
\label{lemma-cat-cochain-abelian}
Soit $\mathcal{A}$ une catégorie abélienne.
\begin{enumerate}
\item La catégorie des complexes de cochaînes dans $\mathcal{A}$ est
abélienne.
\item Un morphisme de complexes de cochaînes
$f : A^\bullet \to B^\bullet$ est injectif
si et seulement si chaque $f^n : A^n \to B^n$ est injectif.
\item Un morphisme de complexes de cochaînes
$f : A^\bullet \to B^\bullet$ est surjectif
si et seulement si chaque $f^n : A^n \to B^n$ est surjectif.
\item Une suite de complexes de cochaînes
$$
A^\bullet \xrightarrow{f} B^\bullet \xrightarrow{g} C^\bullet
$$
est exacte en $B^\bullet$ si et seulement si chaque suite
$$
A^i \xrightarrow{f^i} B^i \xrightarrow{g^i} C^i
$$
est exacte en $B^i$.
\end{enumerate}
\end{lemma}

\begin{proof}
La démonstration est omise.
\end{proof}

\noindent
Pour tout $i \in \mathbf{Z}$, le $i$-ième {\it objet de cohomologie}
d'un complexe de cochaînes $A^\bullet$ est défini par
la formule suivante :
$$
H^i(A^\bullet) = \Ker(d^i)/\Im(d^{i - 1}).
$$
Si $f : A^\bullet \to B^\bullet$ est un morphisme de complexes de cochaînes
de $\mathcal{A}$, on obtient un
morphisme induit $H^i(f) : H^i(A^\bullet) \to H^i(B^\bullet)$
car, manifestement,
$f^i(\Ker(d^i : A^i \to A^{i + 1})) \subset
\Ker(d^i : B^i \to B^{i + 1})$, et de même
pour $\Im(d^{i - 1})$.
On obtient ainsi un foncteur
$$
H^i : \text{CoCh}(\mathcal{A}) \longrightarrow \mathcal{A}.
$$

\begin{definition}
\label{definition-quasi-isomorphism-cochain}
Soit $\mathcal{A}$ une catégorie abélienne.
\begin{enumerate}
\item Un morphisme de complexes de cochaînes $f : A^\bullet \to B^\bullet$
de $\mathcal{A}$ est appelé un {\it quasi-isomorphisme} si les applications
induites $H^i(f) : H^i(A^\bullet) \to H^i(B^\bullet)$
sont des isomorphismes pour tout $i \in \mathbf{Z}$.
\item Un complexe de cochaînes $A^\bullet$ est dit
{\it acyclique} si tous ses objets de cohomologie
$H^i(A^\bullet)$ sont nuls.
\end{enumerate}
\end{definition}

\begin{lemma}
\label{lemma-map-cohomology-homotopy-cochain}
Soit $\mathcal{A}$ une catégorie abélienne.
\begin{enumerate}
\item Si les applications $f, g : A^\bullet \to B^\bullet$ sont
homotopes, alors les applications induites $H^i(f)$ et $H^i(g)$
sont égales.
\item Si $f : A^\bullet \to B^\bullet$ est un homotopisme,
alors $f$ est un quasi-isomorphisme.
\end{enumerate}
\end{lemma}

\begin{proof}
La démonstration est omise.
\end{proof}

\begin{lemma}
\label{lemma-long-exact-sequence-cochain}
\begin{slogan}
Les suites exactes courtes de complexes donnent naissance à des suites exactes longues
d'(co)homologie.
\end{slogan}
Soit $\mathcal{A}$ une catégorie abélienne. Supposons que
$$
0 \to
A^\bullet \to
B^\bullet \to
C^\bullet \to
0
$$
soit une suite exacte courte de complexes de cochaînes de $\mathcal{A}$.
Il existe alors une suite exacte longue de cohomologie
$$
\xymatrix{
\ldots & \ldots & \ldots \ar[lld] \\
H^i(A^\bullet) \ar[r] &
H^i(B^\bullet) \ar[r] &
H^i(C^\bullet) \ar[lld] \\
H^{i + 1}(A^\bullet) \ar[r] &
H^{i + 1}(B^\bullet) \ar[r] &
H^{i + 1}(C^\bullet) \ar[lld] \\
\ldots & \ldots & \ldots \\
}
$$
La construction produit des suites exactes longues de cohomologie
qui sont fonctorielles en la suite exacte courte
et compatibles aux décalages au sens de la
définition \ref{definition-cohomology-shift}.
\end{lemma}

\begin{proof}
Pour les applications horizontales $H^i(A^\bullet) \to H^i(B^\bullet)$ et
$H^i(B^\bullet) \to H^i(C^\bullet)$, on utilise le fait que $H^i$
est un foncteur, comme on l'a vu ci-dessus. Pour l'« homomorphisme bord »
$H^i(C^\bullet) \to H^{i + 1}(A^\bullet)$, on utilise l'application $\delta$
du lemme du serpent \ref{lemma-snake}
appliqué au diagramme
$$
\xymatrix{
&
A^i/\Im(d_A^{i - 1}) \ar[r] \ar[d]^{d_A^i} &
B^i/\Im(d_B^{i - 1}) \ar[r] \ar[d]^{d_B^i} &
C^i/\Im(d_C^{i - 1}) \ar[r] \ar[d]^{d_C^i} &
0 \\
0 \ar[r] &
\Ker(d_A^{i + 1}) \ar[r] &
\Ker(d_B^{i + 1}) \ar[r] &
\Ker(d_C^{i + 1}) &
}
$$
Cela fonctionne car le noyau de l'application verticale de droite est égal à
$H^i(C^\bullet)$ et le conoyau de l'application verticale de gauche est
$H^{i + 1}(A^\bullet)$. L'exactitude de la suite longue est
l'exactitude de la partie (2) du lemme \ref{lemma-snake}.
La fonctorialité est le lemme \ref{lemma-snake-natural}. La compatibilité
aux décalages résulte immédiatement des définitions.
\end{proof}

\section{Homotopie et foncteur de translation}
\label{section-homotopy-shift}

\noindent
Il est fâcheux que signes et indices doivent
intervenir dans toute discussion d'algèbre
homologique\footnote{Merci de nous signaler toute erreur de signe que vous
remarqueriez, ainsi que toute amélioration possible de nos conventions.}.

\begin{definition}
\label{definition-shift}
Soit $\mathcal{A}$ une catégorie additive.
Soit $A_\bullet$ un complexe de chaînes
de différentielles $d_{A, n} : A_n \to A_{n - 1}$.
Pour tout $k \in \mathbf{Z}$, on définit le
{\it complexe de chaînes décalé de $k$, $A[k]_\bullet$}
comme suit :
\begin{enumerate}
\item on pose $A[k]_n = A_{n + k}$, et
\item on définit $d_{A[k], n} : A[k]_n \to A[k]_{n - 1}$
par $d_{A[k], n} = (-1)^k d_{A, n + k}$.
\end{enumerate}
Si $f : A_\bullet \to B_\bullet$ est un morphisme de
complexes de chaînes, on note
$f[k] : A[k]_\bullet \to B[k]_\bullet$ le
morphisme de complexes de chaînes donné par
$f[k]_n = f_{k + n}$.
\end{definition}

\noindent
Bien entendu, cela signifie que l'on dispose de foncteurs
$[k] : \text{Ch}(\mathcal{A}) \to \text{Ch}(\mathcal{A})$
qui commutent deux à deux (au sens strict, sans
qu'interviennent d'isomorphismes de foncteurs),
tels que $A[k][l]_\bullet = A[k + l]_\bullet$ et
$[0] = \text{id}_{\text{Ch}(\mathcal{A})}$.

\medskip\noindent
Rappelons que l'on considère $\mathcal{A}$ comme une sous-catégorie pleine de
$\text{Ch}(\mathcal{A})$ ; voir la section \ref{section-complexes}.
Ainsi, pour tout objet $A$ de $\mathcal{A}$, la notation $A[k]$
désigne l'unique complexe de chaînes nul en tout degré sauf à avoir
$A$ en degré $-k$.

\begin{definition}
\label{definition-homology-shift}
Soit $\mathcal{A}$ une catégorie abélienne.
Soit $A_\bullet$ un complexe de chaînes
de différentielles $d_{A, n} : A_n \to A_{n - 1}$.
Pour tout $k \in \mathbf{Z}$, on identifie
{\it $H_{i + k}(A_\bullet) \rightarrow H_i(A[k]_\bullet)$}
au moyen de l'identification
$A_{i + k} = A[k]_i$.
\end{definition}

\noindent
Cette identification est fonctorielle en $A_\bullet$.
Notons que, puisqu'aucun signe n'intervient dans cette
définition, on obtient en fait un système compatible
d'identifications de tous les objets d'homologie
$H_{i - k}(A[k]_\bullet)$, qui sont en outre
compatibles aux identifications
$A[k][l]_\bullet = A[k + l]_\bullet$ et à
$[0] = \text{id}_{\text{Ch}(\mathcal{A})}$.

\medskip\noindent
Soit $\mathcal{A}$ une catégorie additive.
Supposons que $A_\bullet$ et $B_\bullet$ soient des
complexes de chaînes, que $a, b : A_\bullet \to B_\bullet$ soient des
morphismes de complexes de chaînes et que $\{h_i : A_i \to B_{i + 1}\}$
soit une homotopie entre $a$ et $b$. Rappelons que cela signifie
que
$a_i - b_i = d_{i + 1} \circ h_i + h_{i - 1} \circ d_i$.
Que se passe-t-il si $a = b$ ? On obtient alors la formule
$0 = d_{i + 1} \circ h_i + h_{i - 1} \circ d_i$,
autrement dit $ - d_{i + 1} \circ h_i = h_{i - 1} \circ d_i$.
D'après la définition ci-dessus, cela signifie que la famille $\{h_i\}$
définit un morphisme de complexes de chaînes
$$
A_\bullet \longrightarrow B[1]_\bullet.
$$
Une telle donnée équivaut à un morphisme $A[-1]_\bullet \to B_\bullet$
d'après les remarques précédentes. Ceci démontre le lemme suivant.

\begin{lemma}
\label{lemma-homotopy-shift}
Soit $\mathcal{A}$ une catégorie additive.
Supposons que $A_\bullet$ et $B_\bullet$ soient des
complexes de chaînes. Pour tout morphisme de complexes de chaînes
$a : A_\bullet \to B_\bullet$, il existe
une bijection entre l'ensemble des homotopies
de $a$ à $a$ et
$\Mor_{\text{Ch}(\mathcal{A})}(A_\bullet, B[1]_\bullet)$.
Plus généralement, l'ensemble des homotopies entre
$a$ et $b$ est soit vide, soit un espace principal homogène
sous le groupe
$\Mor_{\text{Ch}(\mathcal{A})}(A_\bullet, B[1]_\bullet)$.
\end{lemma}

\begin{proof}
Voir ci-dessus.
\end{proof}

\begin{lemma}
\label{lemma-ses-termwise-split}
Soit $\mathcal{A}$ une catégorie abélienne.
Soit
$$
0 \to A_\bullet \to B_\bullet \to C_\bullet \to 0
$$
une suite exacte courte de complexes.
Supposons que $\{s_n : C_n \to B_n\}$ soit une famille
de morphismes scindant les suites exactes courtes
$0 \to A_n \to B_n \to C_n \to 0$. Soient
$\pi_n : B_n \to A_n$ les
projections associées ; voir le lemme \ref{lemma-ses-split}.
Alors la famille de morphismes
$$
\pi_{n - 1} \circ d_{B, n} \circ s_n
:
C_n \to A_{n - 1}
$$
définit un morphisme de complexes $\delta(s) : C_\bullet \to A[-1]_\bullet$.
\end{lemma}

\begin{proof}
Notons $i : A_\bullet \to B_\bullet$ et $q : B_\bullet \to C_\bullet$
les applications de complexes de la suite exacte courte. Alors
$i_{n - 1} \circ \pi_{n - 1} \circ d_{B, n} \circ s_n =
d_{B, n} \circ s_n - s_{n - 1} \circ d_{C, n}$. Par conséquent,
$i_{n - 2} \circ d_{A, n - 1} \circ \pi_{n - 1} \circ d_{B, n} \circ s_n =
d_{B, n - 1} \circ (d_{B, n} \circ s_n - s_{n - 1} \circ d_{C, n}) =
- d_{B, n - 1} \circ s_{n - 1} \circ d_{C, n}$, comme désiré.
\end{proof}

\begin{lemma}
\label{lemma-ses-termwise-split-long}
Notations et hypothèses comme dans le lemme \ref{lemma-ses-termwise-split} ci-dessus.
Le morphisme de complexes $\delta(s) : C_\bullet \to A[-1]_\bullet$
induit les applications
$$
H_i(\delta(s)) :
H_i(C_\bullet) \longrightarrow H_i(A[-1]_\bullet) = H_{i - 1}(A_\bullet)
$$
qui interviennent dans la suite exacte longue d'homologie associée
à la suite exacte courte de complexes de chaînes par le
lemme \ref{lemma-long-exact-sequence-chain}.
\end{lemma}

\begin{proof}
La démonstration est omise.
\end{proof}

\begin{lemma}
\label{lemma-ses-termwise-split-homotopy}
Notations et hypothèses comme dans le lemme \ref{lemma-ses-termwise-split} ci-dessus.
Supposons que $\{s'_n : C_n \to B_n\}$ soit un second choix de scindages.
Écrivons $s'_n = s_n + i_n \circ h_n$ pour d'uniques
morphismes $h_n : C_n \to A_n$. La famille d'applications
$\{h_n : C_n \to A[-1]_{n + 1}\}$ est une homotopie entre
les morphismes associés
$\delta(s), \delta(s') : C_\bullet \to A[-1]_\bullet$.
\end{lemma}

\begin{proof}
La démonstration est omise.
\end{proof}

\begin{definition}
\label{definition-shift-cochain}
Soit $\mathcal{A}$ une catégorie additive.
Soit $A^\bullet$ un complexe de cochaînes
de différentielles $d_A^n : A^n \to A^{n + 1}$.
Pour tout $k \in \mathbf{Z}$, on définit le
{\it complexe de cochaînes décalé de $k$, $A[k]^\bullet$}
comme suit :
\begin{enumerate}
\item on pose $A[k]^n = A^{n + k}$, et
\item on définit $d_{A[k]}^n : A[k]^n \to A[k]^{n + 1}$
par $d_{A[k]}^n = (-1)^k d_A^{n + k}$.
\end{enumerate}
Si $f : A^\bullet \to B^\bullet$ est un morphisme de
complexes de cochaînes, on note
$f[k] : A[k]^\bullet \to B[k]^\bullet$ le
morphisme de complexes de cochaînes donné par
$f[k]^n = f^{k + n}$.
\end{definition}

\noindent
Bien entendu, cela signifie que l'on dispose de foncteurs
$[k] : \text{CoCh}(\mathcal{A}) \to \text{CoCh}(\mathcal{A})$
qui commutent deux à deux (au sens strict, sans
qu'interviennent d'isomorphismes de foncteurs) et
tels que $A[k][l]^\bullet = A[k + l]^\bullet$ et
$[0] = \text{id}_{\text{CoCh}(\mathcal{A})}$.

\medskip\noindent
Rappelons que l'on considère $\mathcal{A}$ comme une sous-catégorie pleine de
$\text{CoCh}(\mathcal{A})$ ; voir la section \ref{section-complexes}.
Ainsi, pour tout objet $A$ de $\mathcal{A}$, la notation $A[k]$
désigne l'unique complexe de cochaînes nul en tout degré sauf à avoir
$A$ en degré $-k$.

\begin{definition}
\label{definition-cohomology-shift}
Soit $\mathcal{A}$ une catégorie abélienne.
Soit $A^\bullet$ un complexe de cochaînes
de différentielles $d_A^n : A^n \to A^{n + 1}$.
Pour tout $k \in \mathbf{Z}$, on identifie
{\it $H^{i + k}(A^\bullet) \longrightarrow H^i(A[k]^\bullet)$}
au moyen de l'identification $A^{i + k} = A[k]^i$.
\end{definition}

\noindent
Cette identification est fonctorielle en $A^\bullet$.
Notons que, puisqu'aucun signe n'intervient dans cette
définition, on obtient en fait un système compatible
d'identifications de tous les objets d'homologie
$H^{i - k}(A[k]^\bullet)$, qui sont en outre
compatibles aux identifications
$A[k][l]^\bullet = A[k + l]^\bullet$ et à
$[0] = \text{id}_{\text{CoCh}(\mathcal{A})}$.

\medskip\noindent
Soit $\mathcal{A}$ une catégorie additive.
Supposons que $A^\bullet$ et $B^\bullet$ soient des
complexes de cochaînes, que $a, b : A^\bullet \to B^\bullet$ soient des
morphismes de complexes de cochaînes et que $\{h^i : A^i \to B^{i - 1}\}$
soit une homotopie entre $a$ et $b$. Rappelons que cela signifie
que
$a^i - b^i = d^{i - 1} \circ h^i + h^{i + 1} \circ d^i$.
Que se passe-t-il si $a = b$ ? On obtient alors la formule
$0 = d^{i - 1} \circ h^i + h^{i + 1} \circ d^i$,
autrement dit $ - d^{i - 1} \circ h^i = h^{i + 1} \circ d^i$.
D'après la définition ci-dessus, cela signifie que la famille $\{h^i\}$
définit un morphisme de complexes de cochaînes
$$
A^\bullet \longrightarrow B[-1]^\bullet.
$$
Une telle donnée équivaut à un morphisme $A[1]^\bullet \to B^\bullet$
d'après les remarques précédentes. Ceci démontre le lemme suivant.

\begin{lemma}
\label{lemma-homotopy-shift-cochain}
Soit $\mathcal{A}$ une catégorie additive.
Supposons que $A^\bullet$ et $B^\bullet$ soient des
complexes de cochaînes. Pour tout morphisme de complexes de cochaînes
$a : A^\bullet \to B^\bullet$, il existe
une bijection entre l'ensemble des homotopies
de $a$ à $a$ et
$\Mor_{\text{CoCh}(\mathcal{A})}(A^\bullet, B[-1]^\bullet)$.
Plus généralement, l'ensemble des homotopies entre
$a$ et $b$ est soit vide, soit un espace principal homogène
sous le groupe
$\Mor_{\text{CoCh}(\mathcal{A})}(A^\bullet, B[-1]^\bullet)$.
\end{lemma}

\begin{proof}
Voir ci-dessus.
\end{proof}

\begin{lemma}
\label{lemma-ses-termwise-split-cochain}
Soit $\mathcal{A}$ une catégorie additive.
Soit
$$
0 \to A^\bullet \to B^\bullet \to C^\bullet \to 0
$$
un complexe (!) de complexes.
Supposons donnés des scindages $B^n = A^n \oplus C^n$
compatibles avec les applications de la suite affichée.
Soient $s^n : C^n \to B^n$ et $\pi^n : B^n \to A^n$ les
applications correspondantes. Alors la famille de morphismes
$$
\pi^{n + 1} \circ d_B^n \circ s^n
:
C^n \to A^{n + 1}
$$
définit un morphisme de complexes $\delta : C^\bullet \to A[1]^\bullet$.
\end{lemma}

\begin{proof}
Notons $i : A^\bullet \to B^\bullet$ et $q : B^\bullet \to C^\bullet$
les applications de complexes de la suite exacte courte. Alors
$i^{n + 1} \circ \pi^{n + 1} \circ d_B^n \circ s^n =
d_B^n \circ s^n - s^{n + 1} \circ d_C^n$. Par conséquent,
$i^{n + 2} \circ d_A^{n + 1} \circ \pi^{n + 1} \circ d_B^n \circ s^n =
d_B^{n + 1} \circ (d_B^n \circ s^n - s^{n + 1} \circ d_C^n) =
- d_B^{n + 1} \circ s^{n + 1} \circ d_C^n$, comme désiré.
\end{proof}

\begin{lemma}
\label{lemma-ses-termwise-split-long-cochain}
Notations et hypothèses comme dans le
lemme \ref{lemma-ses-termwise-split-cochain} ci-dessus.
Supposons en outre que $\mathcal{A}$ soit abélienne.
Le morphisme de complexes $\delta : C^\bullet \to A[1]^\bullet$
induit les applications
$$
H^i(\delta) :
H^i(C^\bullet) \longrightarrow H^i(A[1]^\bullet) = H^{i + 1}(A^\bullet)
$$
qui interviennent dans la suite exacte longue de cohomologie associée
à la suite exacte courte de complexes de cochaînes par le
lemme \ref{lemma-long-exact-sequence-cochain}.
\end{lemma}

\begin{proof}
La démonstration est omise.
\end{proof}

\begin{lemma}
\label{lemma-ses-termwise-split-homotopy-cochain}
Notations et hypothèses comme dans le
lemme \ref{lemma-ses-termwise-split-cochain}.
Soient $\alpha : A^\bullet \to B^\bullet$ et
$\beta : B^\bullet \to C^\bullet$ les
morphismes de complexes donnés.
Supposons que $(s')^n : C^n \to B^n$ et $(\pi')^n : B^n \to A^n$
constituent un second choix de scindages.
Écrivons $(s')^n = s^n + \alpha^n \circ h^n$ et
$(\pi')^n = \pi^n + g^n \circ \beta^n$ pour d'uniques
morphismes $h^n : C^n \to A^n$ et $g^n : C^n \to A^n$. Alors :
\begin{enumerate}
\item $g^n = - h^n$, et
\item la famille d'applications $\{g^n : C^n \to A[1]^{n - 1}\}$ est une homotopie
entre $\delta, \delta' : C^\bullet \to A[1]^\bullet$ ; plus précisément,
$(\delta')^n = \delta^n + g^{n + 1} \circ d_C^n + d_{A[1]}^{n - 1} \circ g^n$.
\end{enumerate}
\end{lemma}

\begin{proof}
Comme $(s')^n$ et $(\pi')^n$ sont des scindages, on a $(\pi')^n \circ (s')^n = 0$.
Ainsi
$$
0 = ( \pi^n + g^n \circ \beta^n ) \circ ( s^n + \alpha^n \circ h^n ) =
g^n \circ \beta^n \circ s^n + \pi^n \circ \alpha^n \circ h^n =
g^n + h^n
$$
ce qui démontre (1). Calculons $(\delta')^n$ comme suit :
$$
( \pi^{n + 1} + g^{n + 1} \circ \beta^{n + 1} )
\circ d_B^n \circ
( s^n + \alpha^n \circ h^n )
= \delta^n + g^{n + 1} \circ d_C^n + d_A^n \circ h^n
$$
Puisque $h^n = -g^n$ et $d_{A[1]}^{n - 1} = -d_A^n$, on conclut que (2)
est vérifié.
\end{proof}




\section{Troncatures des complexes}
\label{section-truncations}

\noindent
Soit $\mathcal{A}$ une catégorie abélienne.
Soit $A_\bullet$ un complexe de chaînes. Il existe
plusieurs façons de {\it tronquer} le complexe $A_\bullet$.
\begin{enumerate}
\item La {\it troncature « bête » $\sigma_{\leq n}$}
est le sous-complexe $\sigma_{\leq n} A_\bullet$ défini
par la règle $(\sigma_{\leq n} A_\bullet)_i = 0$ si
$i > n$ et $(\sigma_{\leq n} A_\bullet)_i = A_i$ si
$i \leq n$. En image :
$$
\xymatrix{
\sigma_{\leq n}A_\bullet \ar[d]  &
\ldots \ar[r] &
0 \ar[r] \ar[d] &
A_n \ar[r] \ar[d] &
A_{n - 1} \ar[r] \ar[d] &
\ldots \\
A_\bullet  &
\ldots \ar[r] &
A_{n + 1} \ar[r] &
A_n \ar[r] &
A_{n - 1} \ar[r] &
\ldots
}
$$
Notons la propriété
$\sigma_{\leq n}A_\bullet / \sigma_{\leq n - 1}A_\bullet = A_n[-n]$.
\item La {\it troncature « bête » $\sigma_{\geq n}$}
est le complexe quotient $\sigma_{\geq n} A_\bullet$ défini
par la règle $(\sigma_{\geq n} A_\bullet)_i = A_i$ si
$i \geq n$ et $(\sigma_{\geq n} A_\bullet)_i = 0$ si
$i < n$. En image :
$$
\xymatrix{
A_\bullet \ar[d]  &
\ldots \ar[r] &
A_{n + 1} \ar[r] \ar[d] &
A_n \ar[r] \ar[d] &
A_{n - 1} \ar[r] \ar[d] &
\ldots \\
\sigma_{\geq n}A_\bullet  &
\ldots \ar[r] &
A_{n + 1} \ar[r] &
A_n \ar[r] &
0 \ar[r] &
\ldots
}
$$
L'application de complexes
$\sigma_{\geq n}A_\bullet \to \sigma_{\geq n + 1}A_\bullet$ est surjective
et son noyau est $A_n[-n]$.
\item La {\it troncature canonique} $\tau_{\geq n}A_\bullet$
est définie par l'image
$$
\xymatrix{
\tau_{\geq n}A_\bullet \ar[d]  &
\ldots \ar[r] &
A_{n + 1} \ar[r] \ar[d] &
\Ker(d_n) \ar[r] \ar[d] &
0 \ar[r] \ar[d] &
\ldots \\
A_\bullet  &
\ldots \ar[r] &
A_{n + 1} \ar[r] &
A_n \ar[r] &
A_{n - 1} \ar[r] &
\ldots
}
$$
Notons que ces complexes ont la propriété suivante :
$$
H_i(\tau_{\geq n}A_\bullet) =
\left\{
\begin{matrix}
H_i(A_\bullet) & \text{si} & i \geq n \\
0 & \text{si} & i < n
\end{matrix}
\right.
$$
\item La {\it troncature canonique} $\tau_{\leq n}A_\bullet$
est définie par l'image
$$
\xymatrix{
A_\bullet \ar[d]  &
\ldots \ar[r] &
A_{n + 1} \ar[r] \ar[d] &
A_n \ar[r] \ar[d] &
A_{n - 1} \ar[r] \ar[d] &
\ldots \\
\tau_{\leq n}A_\bullet  &
\ldots \ar[r] &
0 \ar[r] &
\Coker(d_{n + 1}) \ar[r] &
A_{n - 1} \ar[r] &
\ldots
}
$$
Notons que ces complexes ont la propriété suivante :
$$
H_i(\tau_{\leq n}A_\bullet) =
\left\{
\begin{matrix}
H_i(A_\bullet) & \text{si} & i \leq n \\
0 & \text{si} & i > n
\end{matrix}
\right.
$$
\end{enumerate}

\noindent
Soit $\mathcal{A}$ une catégorie abélienne.
Soit $A^\bullet$ un complexe de cochaînes. Il existe
quatre façons de tronquer le complexe $A^\bullet$.
\begin{enumerate}
\item La {\it troncature « bête » $\sigma_{\geq n}$} est le sous-complexe
$\sigma_{\geq n} A^\bullet$ défini par la règle
$(\sigma_{\geq n} A^\bullet)^i = 0$ si
$i < n$ et $(\sigma_{\geq n} A^\bullet)^i = A^i$ si
$i \geq n$. En image :
$$
\xymatrix{
\sigma_{\geq n}A^\bullet \ar[d]  &
\ldots \ar[r] &
0 \ar[r] \ar[d] &
A^n \ar[r] \ar[d] &
A^{n + 1} \ar[r] \ar[d] &
\ldots \\
A^\bullet  &
\ldots \ar[r] &
A^{n - 1} \ar[r] &
A^n \ar[r] &
A^{n + 1} \ar[r] &
\ldots
}
$$
Notons la propriété
$\sigma_{\geq n}A^\bullet / \sigma_{\geq n + 1}A^\bullet
= A^n[-n]$.
\item La {\it troncature « bête » $\sigma_{\leq n}$}
est le complexe quotient $\sigma_{\leq n} A^\bullet$ défini
par la règle $(\sigma_{\leq n} A^\bullet)^i = 0$ si
$i > n$ et $(\sigma_{\leq n} A^\bullet)^i = A^i$ si
$i \leq n$. En image :
$$
\xymatrix{
A^\bullet \ar[d]  &
\ldots \ar[r] &
A^{n - 1} \ar[r] \ar[d] &
A^n \ar[r] \ar[d] &
A^{n + 1} \ar[r] \ar[d] &
\ldots \\
\sigma_{\leq n}A^\bullet &
\ldots \ar[r] &
A^{n - 1} \ar[r] &
A^n \ar[r] &
0 \ar[r] &
\ldots \\
}
$$
L'application de complexes
$\sigma_{\leq n}A^\bullet \to \sigma_{\leq n - 1}A^\bullet$ est surjective
et son noyau est $A^n[-n]$.
\item La {\it troncature canonique} $\tau_{\leq n}A^\bullet$
est définie par l'image
$$
\xymatrix{
\tau_{\leq n}A^\bullet \ar[d]  &
\ldots \ar[r] &
A^{n - 1} \ar[r] \ar[d] &
\Ker(d^n) \ar[r] \ar[d] &
0 \ar[r] \ar[d] &
\ldots \\
A^\bullet  &
\ldots \ar[r] &
A^{n - 1} \ar[r] &
A^n \ar[r] &
A^{n + 1} \ar[r] &
\ldots
}
$$
Notons que ces complexes ont la propriété suivante :
$$
H^i(\tau_{\leq n}A^\bullet) =
\left\{
\begin{matrix}
H^i(A^\bullet) & \text{si} & i \leq n \\
0 & \text{si} & i > n
\end{matrix}
\right.
$$
\item La {\it troncature canonique} $\tau_{\geq n}A^\bullet$
est définie par l'image
$$
\xymatrix{
A^\bullet \ar[d] &
\ldots \ar[r] &
A^{n - 1} \ar[r] \ar[d] &
A^n \ar[r] \ar[d] &
A^{n + 1} \ar[r] \ar[d] &
\ldots \\
\tau_{\geq n}A^\bullet &
\ldots \ar[r] &
0 \ar[r] &
\Coker(d^{n - 1}) \ar[r] &
A^{n + 1} \ar[r] &
\ldots
}
$$
Notons que ces complexes ont la propriété suivante :
$$
H^i(\tau_{\geq n}A^\bullet) =
\left\{
\begin{matrix}
0 & \text{si} & i < n \\
H^i(A^\bullet) & \text{si} & i \geq n
\end{matrix}
\right.
$$
\end{enumerate}


















\section{Objets gradués}
\label{section-graded}

\noindent
Posons la définition suivante.

\begin{definition}
\label{definition-graded}
Soit $\mathcal{A}$ une catégorie additive. La {\it catégorie des objets
gradués de $\mathcal{A}$}, notée $\text{Gr}(\mathcal{A})$, est
la catégorie dont
\begin{enumerate}
\item les objets $A = (A^i)$ sont les familles d'objets $A^i$, $i \in \mathbf{Z}$,
de $\mathcal{A}$, et
\item les morphismes $f : A = (A^i) \to B = (B^i)$ sont les familles de
morphismes $f^i : A^i \to B^i$ de $\mathcal{A}$.
\end{enumerate}
\end{definition}

\noindent
Si $\mathcal{A}$ admet les sommes directes dénombrables, on peut associer à
un objet $A = (A^i)$ de $\text{Gr}(\mathcal{A})$ l'objet
$$
A = \bigoplus\nolimits_{i \in \mathbf{Z}} A^i
$$
et poser $k^iA = A^i$. Dans ce cas, $\text{Gr}(\mathcal{A})$ est équivalente
à la catégorie des couples $(A, k)$ constitués d'un objet $A$ de
$\mathcal{A}$ et d'une décomposition en somme directe
$$
A = \bigoplus\nolimits_{i \in \mathbf{Z}} k^iA
$$
par des facteurs directs indexés par $\mathbf{Z}$ ; un morphisme $(A, k) \to (B, k)$
entre de tels objets est donné par un morphisme $\varphi : A \to B$ de $\mathcal{A}$
tel que $\varphi(k^iA) \subset k^iB$ pour tout $i \in \mathbf{Z}$. Chaque fois
que notre catégorie additive $\mathcal{A}$ admet les sommes directes dénombrables,
nous utiliserons cette équivalence sans autre mention.

\medskip\noindent
Cependant, avec nos définitions, une catégorie additive ou abélienne n'admet pas
nécessairement toutes les sommes directes (dénombrables). Dans ce cas, notre définition
garde un sens. Par exemple, si $\mathcal{A} = \text{Vect}_k$ est la
catégorie des espaces vectoriels de dimension finie sur un corps $k$, alors
$\text{Gr}(\text{Vect}_k)$ est la catégorie des espaces
vectoriels munis d'une graduation dont toutes les composantes graduées sont de dimension
finie, et non la catégorie des espaces vectoriels de dimension finie
munis d'une graduation.

\begin{lemma}
\label{lemma-graded}
Soit $\mathcal{A}$ une catégorie abélienne. La catégorie des objets gradués
$\text{Gr}(\mathcal{A})$ est abélienne.
\end{lemma}

\begin{proof}
Soit $f : A = (A^i) \to B = (B^i)$ un morphisme d'objets gradués
de $\mathcal{A}$ donné par des
morphismes $f^i : A^i \to B^i$ de $\mathcal{A}$.
On a alors $\Ker(f) = (\Ker(f^i))$ et $\Coker(f) = (\Coker(f^i))$
dans la catégorie $\text{Gr}(\mathcal{A})$.
Comme $\Im = \Coim$ dans $\mathcal{A}$,
on voit qu'il en est de même dans $\text{Gr}(\mathcal{A})$.
\end{proof}

\begin{remark}[Avertissement]
\label{remark-direct-sums-not-exact}
Il existe des catégories abéliennes $\mathcal{A}$ qui admettent les sommes directes
dénombrables, mais dans lesquelles celles-ci ne sont pas exactes. Un exemple
est la catégorie opposée à celle des faisceaux abéliens sur $\mathbf{R}$.
En effet, la catégorie des faisceaux abéliens sur $\mathbf{R}$ admet les
produits dénombrables, mais ceux-ci ne sont pas exacts.
Pour une telle catégorie, le foncteur $\text{Gr}(\mathcal{A}) \to \mathcal{A}$,
$(A^i) \mapsto \bigoplus A^i$,
décrit ci-dessus n'est pas exact. Il reste vrai que
$\text{Gr}(\mathcal{A})$ est équivalente à la catégorie des
objets gradués $(A, k)$ de $\mathcal{A}$, mais le noyau, dans la catégorie
des objets gradués, d'une application $\varphi : (A, k) \to (B, k)$ n'est pas égal à
$\Ker(\varphi)$ muni d'une décomposition en somme directe : c'est plutôt
la somme directe des noyaux des applications $k^iA \to k^iB$.
\end{remark}

\begin{definition}
\label{definition-graded-shift}
Soit $\mathcal{A}$ une catégorie additive. Si $A = (A^i)$ est un objet gradué,
alors son $k$-ième {\it décalé} $A[k]$ est l'objet gradué défini par
$A[k]^i = A^{k + i}$.
\end{definition}

\noindent
Si $A$ et $B$ sont des objets gradués de $\mathcal{A}$, on a
\begin{equation}
\label{equation-hom-into-shift}
\Hom_{\text{Gr}(\mathcal{A})}(A, B[k]) =
\Hom_{\text{Gr}(\mathcal{A})}(A[-k], B)
\end{equation}
et un élément de ce groupe est parfois appelé une application d'objets gradués
{\it homogène de degré $k$}.

\medskip\noindent
Pour tout ensemble $G$, on peut définir les objets $G$-gradués de $\mathcal{A}$
comme formant la catégorie dont les objets sont les familles $A = (A^g)_{g \in G}$
d'objets paramétrées par les éléments de $G$. Les morphismes
$f : A \to B$ sont définis comme les familles d'applications $f^g : A^g \to B^g$
où $g$ parcourt les éléments de $G$. Si $G$ est un groupe abélien,
on peut définir sans ambiguïté des foncteurs de translation $[g]$ sur la catégorie
des objets $G$-gradués par la règle $(A[g])^{g_0} = A^{g + g_0}$.
Un cas particulier de ce type de construction est celui où
$G = \mathbf{Z} \times \mathbf{Z}$. Les objets de la
catégorie sont alors appelés des objets {\it bigradués} de $\mathcal{A}$.
La composante $(p, q)$ d'un objet bigradué $A$ est habituellement notée
$A^{p, q}$. Pour $(a, b) \in \mathbf{Z} \times \mathbf{Z}$, on écrit
$A[a, b]$ au lieu de $A[(a, b)]$.
Un morphisme $A \to A[a, b]$ est parfois appelé une
{\it application de bidegré $(a, b)$}.






\section{Catégories monoïdales additives}
\label{section-monoidal}

\noindent
Voici quelques éléments sur l'interaction entre une structure monoïdale
et une structure additive sur une catégorie.

\begin{definition}
\label{definition-additive-monoidal}
Une {\it catégorie monoïdale additive} est une catégorie additive $\mathcal{A}$
munie d'une structure monoïdale $\otimes, \phi$
(Catégories, définition \ref{categories-definition-monoidal-category})
telle que $\otimes$ soit un foncteur additif en chacune des variables.
\end{definition}

\begin{lemma}
\label{lemma-additive-dual}
Soit $\mathcal{A}$ une catégorie monoïdale additive.
Si $Y_i$, $i = 1, 2$, sont des duaux à gauche de $X_i$, $i = 1, 2$, alors
$Y_1 \oplus Y_2$ est un dual à gauche de $X_1 \oplus X_2$.
\end{lemma}

\begin{proof}
Cela résulte de l'unicité des adjoints et de
Catégories, remarque \ref{categories-remark-left-dual-adjoint}.
\end{proof}

\begin{lemma}
\label{lemma-Karoubian-dual}
Dans une catégorie monoïdale additive karoubienne, tout facteur direct
d'un objet qui admet un dual à gauche admet un dual à gauche.
\end{lemma}

\begin{proof}
Nous utiliserons Catégories, lemme \ref{categories-lemma-left-dual},
sans autre mention.
Soit $X$ un objet qui admet un dual à gauche $Y$. On a
$$
\Hom(X, X) = \Hom(\mathbf{1}, X \otimes Y) = \Hom(Y, Y)
$$
Si $a : X \to X$ correspond à $b : Y \to Y$, alors $b$ est l'unique
endomorphisme de $Y$ tel que précomposer par $a$ dans
$$
\Hom(Z' \otimes X, Z) = \Hom(Z', Z \otimes Y)
$$
revienne à postcomposer par $1 \otimes b$.
La bijection $\Hom(X, X) \to \Hom(Y, Y)$, $a \mapsto b$,
est donc un isomorphisme de l'algèbre opposée à $\Hom(X, X)$ sur
l'algèbre $\Hom(Y, Y)$. En particulier, si $X = X_1 \oplus X_2$, alors
les projecteurs correspondants $e_1, e_2$ sont envoyés sur des idempotents
de $\Hom(Y, Y)$. Si $Y = Y_1 \oplus Y_2$ est la décomposition en somme
directe correspondante de $Y$ (section \ref{section-karoubian}),
on voit alors que, sous la bijection
$\Hom(Z' \otimes X, Z) = \Hom(Z', Z \otimes Y)$,
on a $\Hom(Z' \otimes X_i, Z) = \Hom(Z', Z \otimes Y_i)$
fonctoriellement comme sous-groupes, pour $i = 1, 2$.
Il s'ensuit que $Y_i$ est le dual à gauche de
$X_i$, d'après la discussion de
Catégories, remarque \ref{categories-remark-left-dual-adjoint}.
\end{proof}

\begin{example}
\label{example-graded-vector-spaces}
Soit $F$ un corps. Soit $\mathcal{C}$ la catégorie des espaces vectoriels
$F$-gradués. Étant donnés des espaces vectoriels gradués $V$ et $W$, on
note $V \otimes W$ l'espace vectoriel $F$-gradué dont la
partie de degré $n$ est
$$
(V \otimes W)^n = \bigoplus\nolimits_{n = p + q} V^p \otimes_F W^q
$$
Pour un troisième espace vectoriel gradué $U$, on prend comme contrainte d'associativité
$\phi : U \otimes (V \otimes W) \to (U \otimes V) \otimes W$
les isomorphismes « usuels »
$$
U^p \otimes_F (V^q \otimes_F W^r) \to (U^p \otimes_F V^q) \otimes_F W^r
$$
d'espaces vectoriels. Comme unité, on prend l'espace vectoriel $F$-gradué $\mathbf{1}$
qui vaut $F$ en degré $0$ et est nul dans les autres degrés.
Il existe deux contraintes de commutativité sur $\mathcal{C}$
qui font de $\mathcal{C}$ une catégorie monoïdale symétrique : l'une fait
intervenir des signes et l'autre non. Nous utiliserons habituellement
la première. Explicitement, si $V$ et $W$ sont des espaces vectoriels $F$-gradués,
nous utiliserons l'isomorphisme $\psi : V \otimes W \to W \otimes V$
qui, en degré $n$, emploie
$$
V^p \otimes_F W^q \to W^q \otimes_F V^p,\quad
v \otimes w \mapsto (-1)^{pq} w \otimes v
$$
Nous omettons la vérification que cela convient.
\end{example}

\begin{lemma}
\label{lemma-left-dual-graded-vector-spaces}
Soit $F$ un corps. Soit $\mathcal{C}$ la catégorie des espaces vectoriels
$F$-gradués, considérée comme catégorie monoïdale ainsi que dans
l'exemple \ref{example-graded-vector-spaces}. Si $V$ dans $\mathcal{C}$
admet un dual à gauche $W$, alors $\sum_n \dim_F V^n < \infty$
et l'application $\epsilon$ définit des accouplements non dégénérés
$W^{-n} \times V^n \to F$.
\end{lemma}

\begin{proof}
Pour unité, nous prenons l'objet $\mathbf{1}$ décrit ci-dessus.
D'après Catégories, définition \ref{categories-definition-dual}, on a des
applications
$$
\eta : \mathbf{1} \to V \otimes W\quad
\epsilon : W \otimes V \to \mathbf{1}
$$
Comme $\mathbf{1} = F$ placé en degré $0$, on peut
considérer $\epsilon$ comme une famille d'accouplements
$W^{-n} \times V^n \to F$ comme dans l'énoncé du lemme.
Choisissons des bases $\{e_{n, i}\}_{i \in I_n}$ de $V^n$ pour tout $n$.
Écrivons
$$
\eta(1) = \sum e_{n, i} \otimes w_{-n, i}
$$
pour des éléments $w_{-n, i} \in W^{-n}$ presque tous nuls !
La condition que $(\epsilon \otimes 1) \circ (1 \otimes \eta)$ soit
l'identité sur $W$ signifie que
$$
\sum\nolimits_{n, i} \epsilon(w, e_{n, i})w_{-n, i} = w
$$
On voit ainsi que $W$ est engendré, comme espace vectoriel gradué,
par les vecteurs non nuls $w_{-n, i}$, qui sont en nombre fini.
La condition que $(1 \otimes \epsilon) \circ (\eta \otimes 1)$
soit l'identité de $V$ signifie que
$$
\sum\nolimits_{n, i} e_{n, i}\ \epsilon(w_{-n, i}, v) = v
$$
En particulier, en posant $v = e_{n, i}$, on conclut que
$\epsilon(w_{-n, i}, e_{n, i'}) = \delta_{ii'}$. Ainsi,
l'énoncé du lemme est vérifié et
$\{w_{-n, i}\}_{i \in I_n}$ est la base duale de $W^{-n}$ associée à
la base choisie de $V^n$.
\end{proof}
\section{Complexes doubles et complexes simples associés}
\label{section-double-complexes}

\noindent
Nous étudions les complexes doubles et les complexes simples associés.

\begin{definition}
\label{definition-double-complex}
Soit $\mathcal{A}$ une catégorie additive.
Un {\it complexe double} dans $\mathcal{A}$ est la donnée
d'un système $(\{A^{p, q}, d_1^{p, q}, d_2^{p, q}\}_{p, q\in \mathbf{Z}})$,
où chaque $A^{p, q}$ est un objet de $\mathcal{A}$ et où
$d_1^{p, q} : A^{p, q} \to A^{p + 1, q}$ et
$d_2^{p, q} : A^{p, q} \to A^{p, q + 1}$ sont des morphismes de $\mathcal{A}$
qui satisfont aux règles suivantes :
\begin{enumerate}
\item $d_1^{p + 1, q} \circ d_1^{p, q} = 0$
\item $d_2^{p, q + 1} \circ d_2^{p, q} = 0$
\item $d_1^{p, q + 1} \circ d_2^{p, q} = d_2^{p + 1, q} \circ d_1^{p, q}$
\end{enumerate}
pour tous $p, q \in \mathbf{Z}$.
\end{definition}

\noindent
C'est simplement la version cochaînique de la définition.
Elle dit que chaque $A^{p, \bullet}$ est un complexe de cochaînes
et que chaque $d_1^{p, \bullet}$ est un morphisme de complexes
$A^{p, \bullet} \to A^{p + 1, \bullet}$ tel que
$d_1^{p + 1, \bullet} \circ d_1^{p, \bullet} = 0$ comme morphismes
de complexes. Autrement dit, un complexe double peut être considéré comme
un complexe de complexes. Ainsi, dans le diagramme
$$
\xymatrix{
\ldots &
\ldots &
\ldots &
\ldots \\
\ldots \ar[r] &
A^{p, q + 1} \ar[r]^{d_1^{p, q + 1}} \ar[u] &
A^{p + 1, q + 1} \ar[r] \ar[u] &
\ldots \\
\ldots \ar[r] &
A^{p, q} \ar[r]^{d_1^{p, q}} \ar[u]^{d_2^{p, q}} &
A^{p + 1, q} \ar[r] \ar[u]_{d_2^{p + 1, q}} &
\ldots \\
\ldots &
\ldots \ar[u] &
\ldots \ar[u] &
\ldots
}
$$
tout carré est commutatif.
Avertissement : on rencontre dans la littérature une autre définition
selon laquelle les carrés du diagramme anticommutent dans un « bicomplexe »
ou un « complexe double ».

\begin{example}
\label{example-double-complex-as-tensor-product-of}
Soient $\mathcal{A}$, $\mathcal{B}$, $\mathcal{C}$ des catégories additives.
Supposons que
$$
\otimes : \mathcal{A} \times \mathcal{B} \longrightarrow \mathcal{C},
\quad
(X, Y) \longmapsto X \otimes Y
$$
soit un foncteur bilinéaire sur les morphismes ; voir
Catégories, définition \ref{categories-definition-product-category},
pour la définition de $\mathcal{A} \times \mathcal{B}$.
Étant donnés un complexe $X^\bullet$ de $\mathcal{A}$ et un complexe $Y^\bullet$
de $\mathcal{B}$, on obtient un complexe double
$$
K^{\bullet, \bullet} = X^\bullet \otimes Y^\bullet
$$
dans $\mathcal{C}$. Ici, la première différentielle
$K^{p, q} \to K^{p + 1, q}$ est le morphisme
$X^p \otimes Y^q \to X^{p + 1} \otimes Y^q$ induit par
le morphisme $X^p \to X^{p + 1}$ et l'identité de $Y^q$.
De même pour la seconde différentielle.
\end{example}

\begin{definition}
\label{definition-associated-simple-complex}
Soit $\mathcal{A}$ une catégorie additive.
Soit $A^{\bullet, \bullet}$ un complexe double.
Le {\it complexe simple associé}, noté $sA^\bullet$, également
appelé souvent le {\it complexe total associé}, noté
$\text{Tot}(A^{\bullet, \bullet})$, est
donné par
$$
sA^n = \text{Tot}^n(A^{\bullet, \bullet}) =
\bigoplus\nolimits_{n = p + q} A^{p, q}
$$
(s'il existe), avec pour différentielle
$$
d_{sA^\bullet}^n = d_{\text{Tot}(A^{\bullet, \bullet})}^n =
\sum\nolimits_{n = p + q} (d_1^{p, q} + (-1)^p d_2^{p, q})
$$
\end{definition}

\noindent
Si les sommes directes dénombrables existent dans $\mathcal{A}$, ou si, pour chaque $n$,
seuls un nombre fini de $A^{p, n - p}$ sont non nuls, alors
$\text{Tot}(A^{\bullet, \bullet})$ existe. Notons que
la définition n'est {\it pas} symétrique en les indices $(p, q)$.

\begin{remark}
\label{remark-triple-complex}
Soit $\mathcal{A}$ une catégorie additive. Soit $A^{\bullet, \bullet, \bullet}$
un complexe triple. Le complexe simple associé est le complexe de termes
$$
\text{Tot}^n(A^{\bullet, \bullet, \bullet}) =
\bigoplus\nolimits_{p + q + r = n} A^{p, q, r}
$$
et de différentielle
$$
d^n_{\text{Tot}(A^{\bullet, \bullet, \bullet})} =
\sum\nolimits_{p + q + r = n}\left(
d_1^{p, q, r} + (-1)^pd_2^{p, q, r} + (-1)^{p + q}d_3^{p, q, r}\right)
$$
Avec cette définition, un calcul simple montre que le complexe simple associé
est égal à
$$
\text{Tot}(A^{\bullet, \bullet, \bullet}) =
\text{Tot}(\text{Tot}_{12}(A^{\bullet, \bullet, \bullet})) =
\text{Tot}(\text{Tot}_{23}(A^{\bullet, \bullet, \bullet}))
$$
Autrement dit, on peut d'abord combiner les deux premières variables,
puis combiner leur somme avec la dernière, ou bien combiner d'abord les
deux dernières variables, puis la première avec la somme des deux dernières.
\end{remark}

\begin{remark}
\label{remark-shift-double-complex}
Soit $\mathcal{A}$ une catégorie additive. Soit $A^{\bullet, \bullet}$
un complexe double de différentielles $d_1^{p, q}$ et $d_2^{p, q}$.
Notons $A^{\bullet, \bullet}[a, b]$ le complexe double défini par
$$
(A^{\bullet, \bullet}[a, b])^{p, q} = A^{p + a, q + b}
$$
et par les différentielles
$$
d_{A^{\bullet, \bullet}[a, b], 1}^{p, q} = (-1)^a d_1^{p + a, q + b}
\quad\text{et}\quad
d_{A^{\bullet, \bullet}[a, b], 2}^{p, q} = (-1)^b d_2^{p + a, q + b}
$$
Dans cette situation, il existe un isomorphisme bien défini
$$
\gamma :
\text{Tot}(A^{\bullet, \bullet})[a + b]
\longrightarrow
\text{Tot}(A^{\bullet, \bullet}[a, b])
$$
qui, en degré $n$, est donné par l'application
$$
\xymatrix{
(\text{Tot}(A^{\bullet, \bullet})[a + b])^n =
\bigoplus_{p + q = n + a + b} A^{p, q}
\ar[d]^{\epsilon(p, q, a, b)\text{id}_{A^{p, q}}} \\
\text{Tot}(A^{\bullet, \bullet}[a, b])^n =
\bigoplus_{p' + q' = n} A^{p' + a, q' + b}
}
$$
pour un certain signe $\epsilon(p, q, a, b)$. Bien entendu, le facteur $A^{p, q}$
s'envoie sur le facteur $A^{p' + a, q' + b}$ lorsque $p = p' + a$ et $q = q' + b$.
Pour déterminer les conditions sur ces signes, observons que, sur la source, on a
$$
d|_{A^{p, q}} = (-1)^{a + b}\left(d_1^{p, q} + (-1)^pd_2^{p, q}\right)
$$
tandis que, sur le but, on a
$$
d|_{A^{p' + a, q' + b}} =
(-1)^ad_1^{p' + a, q' + b} + (-1)^{p'}(-1)^bd_2^{p' + a, q' + b}
$$
Nos contraintes sont donc
$$
(-1)^a \epsilon(p, q, a, b) = \epsilon(p + 1, q, a, b)(-1)^{a + b}
\Leftrightarrow
\epsilon(p + 1, q, a, b) = (-1)^b \epsilon(p, q, a, b)
$$
et
$$
(-1)^{p' + b}\epsilon(p, q, a, b) =
\epsilon(p, q + 1, a, b) (-1)^{a + b + p}
\Leftrightarrow
\epsilon(p, q, a, b) = \epsilon(p, q + 1, a, b)
$$
On choisit donc $\epsilon(p, q, a, b) = (-1)^{pb}$.
\end{remark}

\begin{remark}
\label{remark-double-complex-complex-of-complexes-first}
Soit $\mathcal{A}$ une catégorie additive admettant les sommes directes dénombrables.
Notons $\text{DoubleComp}(\mathcal{A})$ la catégorie des complexes doubles.
On peut considérer un objet $A^{\bullet, \bullet}$ de
$\text{DoubleComp}(\mathcal{A})$ comme un complexe de complexes
de la manière suivante :
$$
\ldots \to A^{\bullet, -1} \to A^{\bullet, 0} \to A^{\bullet, 1} \to \ldots
$$
Pour la variante où les rôles des indices sont échangés, voir la
remarque \ref{remark-double-complex-complex-of-complexes-second}.
Dans cette remarque, nous montrons que prendre le complexe simple associé
est compatible avec toutes les structures sur les complexes étudiées
jusqu'ici dans le chapitre.

\medskip\noindent
Premièrement, observons que le foncteur de translation sur les complexes doubles considérés
comme complexes de complexes de la manière indiquée ci-dessus est le foncteur
$[0, 1]$ défini dans la remarque \ref{remark-shift-double-complex}.
D'après la remarque \ref{remark-shift-double-complex}, le foncteur
$$
\text{Tot} : \text{DoubleComp}(\mathcal{A}) \to \text{Comp}(\mathcal{A})
$$
est compatible aux foncteurs de translation, au sens où l'on dispose d'un isomorphisme
fonctoriel $\gamma : \text{Tot}(A^{\bullet, \bullet})[1] \to
\text{Tot}(A^{\bullet, \bullet}[0, 1])$.

\medskip\noindent
Deuxièmement, si
$$
f, g : A^{\bullet, \bullet} \to B^{\bullet, \bullet}
$$
sont homotopes lorsque $f$ et $g$ sont considérés comme des morphismes de complexes
de complexes de la manière indiquée ci-dessus, alors
$$
\text{Tot}(f), \text{Tot}(g) :
\text{Tot}(A^{\bullet, \bullet}) \to \text{Tot}(B^{\bullet, \bullet})
$$
sont des applications homotopes de complexes. En effet, soit $h = (h^q)$
une homotopie entre $f$ et $g$. Si l'on note
$h^{p, q} : A^{p, q} \to B^{p, q - 1}$ la composante en degré $p$ de $h^q$,
alors cela signifie que
$$
f^{p, q} - g^{p, q} = d_2^{p, q - 1} \circ h^{p, q} +
h^{p, q + 1} \circ d_2^{p, q}
$$
Le fait que $h^q : A^{\bullet, q} \to B^{\bullet, q - 1}$ soit une application de
complexes signifie que
$$
d_1^{p, q - 1} \circ h^{p, q} = h^{p + 1, q} \circ d_1^{p, q}
$$
Définissons $h' = ((h')^n)$ comme l'homotopie donnée par les applications
$(h')^n : \text{Tot}^n(A^{\bullet, \bullet}) \to
\text{Tot}^{n - 1}(B^{\bullet, \bullet})$
qui valent $(-1)^ph^{p, q}$ sur le facteur $A^{p, q}$ pour $p + q = n$.
On voit alors que
$$
d_{\text{Tot}(B^{\bullet, \bullet})} \circ h' +
h' \circ d_{\text{Tot}(A^{\bullet, \bullet})}
$$
restreint au facteur $A^{p, q}$ est égal à
$$
d_1^{p, q - 1} \circ (-1)^p h^{p, q} +
(-1)^p d_2^{p, q - 1} \circ (-1)^p h^{p, q} +
(-1)^{p + 1} h^{p + 1, q} \circ d_1^{p, q} +
(-1)^p h^{p, q + 1} \circ (-1)^p d_2^{p, q}
$$
ce qui donne $f^{p, q} - g^{p, q}$ d'après les équations ci-dessus.
Ceci démontre la deuxième compatibilité.

\medskip\noindent
Troisièmement, supposons que $f = g$ dans le paragraphe précédent.
Alors l'affectation $h \leadsto h'$ ci-dessus est compatible à
l'identification du lemme \ref{lemma-homotopy-shift-cochain}.
Plus précisément, si l'on considère $h$ comme un morphisme de complexes
de complexes $A^{\bullet, \bullet} \to B^{\bullet, \bullet}[0, -1]$
au moyen de ce lemme, alors
$$
\text{Tot}(A^{\bullet, \bullet})
\xrightarrow{\text{Tot}(h)}
\text{Tot}(B^{\bullet, \bullet}[0, -1])
\xrightarrow{\gamma^{-1}}
\text{Tot}(B^{\bullet, \bullet})[-1]
$$
est égal à $h'$ considéré, au moyen du lemme, comme un morphisme de complexes.
Ici, $\gamma$ est l'identification de la
remarque \ref{remark-shift-double-complex}.
La vérification de ce troisième point est immédiate.

\medskip\noindent
Quatrièmement, soit
$$
0 \to A^{\bullet, \bullet} \to B^{\bullet, \bullet} \to
C^{\bullet, \bullet} \to 0
$$
un complexe de complexes doubles et supposons donnés des scindages
$s^q : C^{\bullet, q} \to B^{\bullet, q}$ et
$\pi^q : B^{\bullet, q} \to A^{\bullet, q}$
comme dans le lemme \ref{lemma-ses-termwise-split-cochain},
lorsque l'on considère les complexes doubles
comme des complexes de complexes de la manière indiquée ci-dessus.
D'une part, ceci produit une application
$$
\delta : C^{\bullet, \bullet} \longrightarrow A^{\bullet, \bullet}[0, 1]
$$
par le procédé du lemme \ref{lemma-ses-termwise-split-cochain}.
D'autre part, en prenant $\text{Tot}$, on obtient un complexe
$$
0 \to \text{Tot}(A^{\bullet, \bullet}) \to
\text{Tot}(B^{\bullet, \bullet}) \to
\text{Tot}(C^{\bullet, \bullet}) \to 0
$$
qui est scindé degré par degré (voir ci-dessous) et qui est donc muni d'un morphisme
$$
\delta' :
\text{Tot}(C^{\bullet, \bullet})
\longrightarrow
\text{Tot}(A^{\bullet, \bullet})[1]
$$
bien défini à homotopie près par les lemmes \ref{lemma-ses-termwise-split-cochain}
et \ref{lemma-ses-termwise-split-homotopy-cochain}. Affirmation :
ces applications coïncident au sens où
$$
\text{Tot}(C^{\bullet, \bullet})
\xrightarrow{\text{Tot}(\delta)}
\text{Tot}(A^{\bullet, \bullet}[0, 1]) \xrightarrow{\gamma^{-1}}
\text{Tot}(A^{\bullet, \bullet})[1]
$$
est égal à $\delta'$, où $\gamma$ est comme dans la
remarque \ref{remark-shift-double-complex}. Pour le voir, notons
$s^{p, q} : C^{p, q} \to B^{p, q}$ et
$\pi^{p, q} : B^{p, q} \to A^{p, q}$ les composantes de $s^q$ et $\pi^q$.
Comme scindages
$(s')^n : \text{Tot}^n(C^{\bullet, \bullet}) \to
\text{Tot}^n(B^{\bullet, \bullet})$
et
$(\pi')^n : \text{Tot}^n(B^{\bullet, \bullet}) \to
\text{Tot}^n(A^{\bullet, \bullet})$,
on prend les applications dont les composantes sont $s^{p, q}$ et $\pi^{p, q}$
pour $p + q = n$. Rappelons que
$$
(\delta')^n =
(\pi')^{n + 1} \circ d_{\text{Tot}(B^{\bullet, \bullet})}^n \circ (s')^n :
\text{Tot}^n(C^{\bullet, \bullet}) \to
\text{Tot}^{n + 1}(A^{\bullet, \bullet})
$$
La restriction de ceci au facteur $C^{p, q}$ est égale à
$$
\pi^{p + 1, q} \circ
d_1^{p, q} \circ
s^{p, q} +
\pi^{p, q + 1} \circ
(-1)^p d_2^{p, q} \circ
s^{p, q} =
\pi^{p, q + 1} \circ
(-1)^p d_2^{p, q} \circ
s^{p, q}
$$
L'égalité vaut parce que $s^q$ est un morphisme de complexes (de différentielle $d_1$)
et parce que $\pi^{p + 1, q} \circ s^{p + 1, q} = 0$,
puisque $s$ et $\pi$ correspondent à une décomposition en somme directe de $B$
en chaque bidegré. D'autre part, pour $\delta$, on a
$$
\delta^q =  \pi^q \circ d_2 \circ s^q :
C^{\bullet, q} \to A^{\bullet, q + 1}
$$
dont la restriction au facteur $C^{p, q}$ est égale à
$\pi^{p, q + 1} \circ d_2^{p, q} \circ s^{p, q}$.
La différence de signes est exactement compensée par le signe
$(-1)^p$ dans l'isomorphisme $\gamma$, ce qui démontre la quatrième affirmation.
\end{remark}

\begin{remark}
\label{remark-double-complex-complex-of-complexes-second}
Soit $\mathcal{A}$ une catégorie additive admettant les sommes directes dénombrables.
Notons $\text{DoubleComp}(\mathcal{A})$ la catégorie des complexes doubles.
On peut considérer un objet $A^{\bullet, \bullet}$ de
$\text{DoubleComp}(\mathcal{A})$ comme un complexe de complexes
de la manière suivante :
$$
\ldots \to A^{-1, \bullet} \to A^{0, \bullet} \to A^{1, \bullet} \to \ldots
$$
Pour la variante où les rôles des indices sont échangés, voir la
remarque \ref{remark-double-complex-complex-of-complexes-first}.
Dans cette remarque, nous montrons que prendre le complexe simple associé
est compatible avec toutes les structures sur les complexes étudiées
jusqu'ici dans le chapitre.

\medskip\noindent
Premièrement, observons que le foncteur de translation sur les complexes doubles considérés
comme complexes de complexes de la manière indiquée ci-dessus est le foncteur
$[1, 0]$ défini dans la remarque \ref{remark-shift-double-complex}.
D'après la remarque \ref{remark-shift-double-complex}, le foncteur
$$
\text{Tot} : \text{DoubleComp}(\mathcal{A}) \to \text{Comp}(\mathcal{A})
$$
est compatible aux foncteurs de translation, au sens où l'on dispose d'un isomorphisme
fonctoriel $\gamma : \text{Tot}(A^{\bullet, \bullet})[1] \to
\text{Tot}(A^{\bullet, \bullet}[1, 0])$.

\medskip\noindent
Deuxièmement, si
$$
f, g : A^{\bullet, \bullet} \to B^{\bullet, \bullet}
$$
sont homotopes lorsque $f$ et $g$ sont considérés comme des morphismes de complexes
de complexes de la manière indiquée ci-dessus, alors
$$
\text{Tot}(f), \text{Tot}(g) :
\text{Tot}(A^{\bullet, \bullet}) \to \text{Tot}(B^{\bullet, \bullet})
$$
sont des applications homotopes de complexes. En effet, soit $h = (h^p)$
une homotopie entre $f$ et $g$. Si l'on note
$h^{p, q} : A^{p, q} \to B^{p - 1, q}$ la composante en degré $q$ de $h^p$,
alors cela signifie que
$$
f^{p, q} - g^{p, q} = d_1^{p - 1, q} \circ h^{p, q} +
h^{p + 1, q} \circ d_1^{p, q}
$$
Le fait que $h^p : A^{p, \bullet} \to B^{p - 1, \bullet}$ soit une application de
complexes signifie que
$$
d_2^{p - 1, q} \circ h^{p, q} = h^{p, q + 1} \circ d_2^{p, q}
$$
Définissons $h' = ((h')^n)$ comme l'homotopie donnée par les applications
$(h')^n : \text{Tot}^n(A^{\bullet, \bullet}) \to
\text{Tot}^{n - 1}(B^{\bullet, \bullet})$
qui valent $h^{p, q}$ sur le facteur $A^{p, q}$ pour $p + q = n$.
On voit alors que
$$
d_{\text{Tot}(B^{\bullet, \bullet})} \circ h' +
h' \circ d_{\text{Tot}(A^{\bullet, \bullet})}
$$
restreint au facteur $A^{p, q}$ est égal à
$$
d_1^{p - 1, q} \circ h^{p, q} +
(-1)^{p - 1} d_2^{p - 1, q} \circ h^{p, q} +
h^{p + 1, q} \circ d_1^{p, q} +
h^{p, q + 1} \circ (-1)^p d_2^{p, q}
$$
ce qui donne $f^{p, q} - g^{p, q}$ d'après les équations ci-dessus.
Ceci démontre la deuxième compatibilité.

\medskip\noindent
Troisièmement, supposons que $f = g$ dans le paragraphe précédent.
Alors l'affectation $h \leadsto h'$ ci-dessus est compatible à
l'identification du lemme \ref{lemma-homotopy-shift-cochain}.
Plus précisément, si l'on considère $h$ comme un morphisme de complexes
de complexes $A^{\bullet, \bullet} \to B^{\bullet, \bullet}[-1, 0]$
au moyen de ce lemme, alors
$$
\text{Tot}(A^{\bullet, \bullet})
\xrightarrow{\text{Tot}(h)}
\text{Tot}(B^{\bullet, \bullet}[-1, 0])
\xrightarrow{\gamma^{-1}}
\text{Tot}(B^{\bullet, \bullet})[-1]
$$
est égal à $h'$ considéré, au moyen du lemme, comme un morphisme de complexes.
Ici, $\gamma$ est l'identification de la
remarque \ref{remark-shift-double-complex}.
La vérification de ce troisième point est immédiate.

\medskip\noindent
Quatrièmement, soit
$$
0 \to A^{\bullet, \bullet} \to B^{\bullet, \bullet} \to
C^{\bullet, \bullet} \to 0
$$
un complexe de complexes doubles et supposons donnés des scindages
$s^p : C^{p, \bullet} \to B^{p, \bullet}$ et
$\pi^p : B^{p, \bullet} \to A^{p, \bullet}$
comme dans le lemme \ref{lemma-ses-termwise-split-cochain},
lorsque l'on considère les complexes doubles
comme des complexes de complexes de la manière indiquée ci-dessus.
D'une part, ceci produit une application
$$
\delta : C^{\bullet, \bullet} \longrightarrow A^{\bullet, \bullet}[1, 0]
$$
par le procédé du lemme \ref{lemma-ses-termwise-split-cochain}.
D'autre part, en prenant $\text{Tot}$, on obtient un complexe
$$
0 \to \text{Tot}(A^{\bullet, \bullet}) \to
\text{Tot}(B^{\bullet, \bullet}) \to
\text{Tot}(C^{\bullet, \bullet}) \to 0
$$
qui est scindé degré par degré (voir ci-dessous) et qui est donc muni d'un morphisme
$$
\delta' :
\text{Tot}(C^{\bullet, \bullet})
\longrightarrow
\text{Tot}(A^{\bullet, \bullet})[1]
$$
bien défini à homotopie près par les lemmes \ref{lemma-ses-termwise-split-cochain}
et \ref{lemma-ses-termwise-split-homotopy-cochain}. Affirmation :
ces applications coïncident au sens où
$$
\text{Tot}(C^{\bullet, \bullet})
\xrightarrow{\text{Tot}(\delta)}
\text{Tot}(A^{\bullet, \bullet}[1, 0]) \xrightarrow{\gamma^{-1}}
\text{Tot}(A^{\bullet, \bullet})[1]
$$
est égal à $\delta'$, où $\gamma$ est comme dans la
remarque \ref{remark-shift-double-complex}. Pour le voir, notons
$s^{p, q} : C^{p, q} \to B^{p, q}$ et
$\pi^{p, q} : B^{p, q} \to A^{p, q}$ les composantes de $s^p$ et $\pi^p$.
Comme scindages
$(s')^n : \text{Tot}^n(C^{\bullet, \bullet}) \to
\text{Tot}^n(B^{\bullet, \bullet})$
et
$(\pi')^n : \text{Tot}^n(B^{\bullet, \bullet}) \to
\text{Tot}^n(A^{\bullet, \bullet})$,
on prend les applications dont les composantes sont $s^{p, q}$ et $\pi^{p, q}$
pour $p + q = n$. Rappelons que
$$
(\delta')^n =
(\pi')^{n + 1} \circ d_{\text{Tot}(B^{\bullet, \bullet})}^n \circ (s')^n :
\text{Tot}^n(C^{\bullet, \bullet}) \to
\text{Tot}^{n + 1}(A^{\bullet, \bullet})
$$
La restriction de ceci au facteur $C^{p, q}$ est égale à
$$
\pi^{p + 1, q} \circ
d_1^{p, q} \circ
s^{p, q} +
\pi^{p, q + 1} \circ
(-1)^p d_2^{p, q} \circ
s^{p, q} =
\pi^{p + 1, q} \circ
d_1^{p, q} \circ
s^{p, q}
$$
L'égalité vaut parce que $s^p$ est un morphisme de complexes (de différentielle $d_2$)
et parce que $\pi^{p, q + 1} \circ s^{p, q + 1} = 0$,
puisque $s$ et $\pi$ correspondent à une décomposition en somme directe de $B$
en chaque bidegré. D'autre part, pour $\delta$, on a
$$
\delta^p =  \pi^p \circ d_1 \circ s^p :
C^{p, \bullet} \to A^{p + 1, \bullet}
$$
dont la restriction au facteur $C^{p, q}$ est égale à
$\pi^{p + 1, q} \circ d_1^{p, q} \circ s^{p, q}$.
On obtient donc la même chose que précédemment, conformément au fait que
l'isomorphisme
$\gamma : \text{Tot}(A^{\bullet, \bullet})[1] \to
\text{Tot}(A^{\bullet, \bullet}[1, 0])$
est défini sans intervention de signes.
\end{remark}








\section{Filtrations}
\label{section-filtrations}

\noindent
Une bonne référence pour ces notions est \cite[section 1]{HodgeII}.
(Notons que nos conventions sur les catégories abéliennes sont différentes.)

\begin{definition}
\label{definition-filtered}
Soit $\mathcal{A}$ une catégorie abélienne.
\begin{enumerate}
\item Une {\it filtration décroissante} $F$ sur un objet $A$
est une famille $(F^nA)_{n \in \mathbf{Z}}$ de sous-objets de $A$ telle que
$$
A \supset \ldots \supset F^nA \supset F^{n + 1}A \supset \ldots \supset 0
$$
\item Un {\it objet filtré de $\mathcal{A}$} est
un couple $(A, F)$ constitué d'un objet $A$ de $\mathcal{A}$
et d'une filtration décroissante $F$ sur $A$.
\item Un {\it morphisme $(A, F) \to (B, F)$ d'objets filtrés}
est donné par un morphisme $\varphi : A \to B$ de $\mathcal{A}$
tel que $\varphi(F^iA) \subset F^iB$ pour tout $i \in \mathbf{Z}$.
\item La catégorie des objets filtrés est notée $\text{Fil}(\mathcal{A})$.
\item Étant donnés un objet filtré $(A, F)$ et un sous-objet $X \subset A$, la
{\it filtration induite} sur $X$ est la filtration définie par $F^nX = X \cap F^nA$.
\item Étant donnés un objet filtré $(A, F)$ et une surjection
$\pi : A \to Y$, la {\it filtration quotient} est la filtration définie par
$F^nY = \pi(F^nA)$.
\item Une filtration $F$ sur un objet $A$ est dite {\it finie}
s'il existe $n, m$ tels que $F^nA = A$ et $F^mA = 0$.
\item Étant donné un objet filtré $(A, F)$, on dit que $\bigcap F^iA$ existe
s'il existe un plus grand sous-objet de $A$ contenu dans tous les $F^iA$.
On dit que $\bigcup F^iA$ existe s'il existe un plus petit sous-objet
de $A$ contenant tous les $F^iA$.
\item La filtration d'un objet filtré $(A, F)$ est dite
{\it séparée} si $\bigcap F^iA = 0$ et
{\it exhaustive} si $\bigcup F^iA = A$.
\end{enumerate}
\end{definition}

\noindent
Par abus de notation, on dit qu'un morphisme $f : (A, F) \to (B, F)$
d'objets filtrés est {\it injectif} si $f : A \to B$ est injectif
dans la catégorie abélienne $\mathcal{A}$. De même, on dit que $f$ est
{\it surjectif} si $f : A \to B$ est surjectif dans la catégorie
$\mathcal{A}$. Être injectif (resp.\ surjectif)
équivaut à être un monomorphisme (resp.\ un épimorphisme)
dans $\text{Fil}(\mathcal{A})$. D'après le
lemme \ref{lemma-filtered},
cela équivaut aussi à avoir un noyau nul (resp.\ un conoyau nul).

\begin{lemma}
\label{lemma-filtered}
Soit $\mathcal{A}$ une catégorie abélienne.
La catégorie des objets filtrés $\text{Fil}(\mathcal{A})$
possède les propriétés suivantes :
\begin{enumerate}
\item C'est une catégorie additive.
\item Elle possède un objet nul.
\item Elle possède des noyaux et des conoyaux, des images et des coimages.
\item En général, ce n'est pas une catégorie abélienne.
\end{enumerate}
\end{lemma}

\begin{proof}
Il est clair que $\text{Fil}(\mathcal{A})$ est additive, la somme
directe étant donnée par $(A, F) \oplus (B, F) = (A \oplus B, F)$, où
$F^p(A \oplus B) = F^pA \oplus F^pB$.
Le noyau d'un morphisme $f : (A, F) \to (B, F)$ d'objets
filtrés est l'injection $\Ker(f) \subset A$, où $\Ker(f)$
est muni de la filtration induite.
Le conoyau d'un morphisme $f : A \to B$ d'objets
filtrés est la surjection $B \to \Coker(f)$, où $\Coker(f)$
est muni de la filtration quotient. Puisque tous les noyaux et conoyaux
existent, il en va de même de toutes les coimages et images. Voir
l'exemple \ref{example-not-abelian}
pour la dernière assertion.
\end{proof}

\begin{definition}
\label{definition-strict}
Soit $\mathcal{A}$ une catégorie abélienne.
Un morphisme $f : A \to B$ d'objets filtrés de $\mathcal{A}$ est
dit {\it strict} si $f(F^iA) = f(A) \cap F^iB$ pour
tout $i \in \mathbf{Z}$.
\end{definition}

\noindent
Cela équivaut aussi à exiger que $f^{-1}(F^iB) = F^iA + \Ker(f)$
pour tout $i \in \mathbf{Z}$. Caractérisons les morphismes stricts
comme suit.

\begin{lemma}
\label{lemma-characterize-strict-general}
Soit $\mathcal{A}$ une catégorie abélienne.
Soit $f : A \to B$ un morphisme d'objets filtrés de $\mathcal{A}$.
Les conditions suivantes sont équivalentes :
\begin{enumerate}
\item $f$ est strict,
\item le morphisme $\Coim(f) \to \Im(f)$ du
lemme \ref{lemma-coim-im-map}
est un isomorphisme.
\end{enumerate}
\end{lemma}

\begin{proof}
Notons que $\Coim(f) \to \Im(f)$ est un isomorphisme entre les
objets de $\mathcal{A}$, et que la partie (2) signifie que c'est
un isomorphisme d'objets filtrés.
D'après la description des noyaux et conoyaux dans la démonstration du
lemme \ref{lemma-filtered},
on voit que la filtration sur $\Coim(f)$ est la
filtration quotient provenant de $A \to \Coim(f)$.
De même, la filtration sur $\Im(f)$ est la filtration
induite provenant de l'injection $\Im(f) \to B$.
Dire que le morphisme est strict signifie exactement que la filtration quotient
est la filtration induite.
\end{proof}

\begin{lemma}
\label{lemma-add-summand-strict-monomorphism}
Soit $\mathcal{A}$ une catégorie abélienne.
Soit $f : A \to B$ un monomorphisme strict d'objets filtrés.
Soit $g : A \to C$ un morphisme d'objets filtrés.
Alors $f \oplus g : A \to B \oplus C$ est un monomorphisme strict.
\end{lemma}

\begin{proof}
Cela résulte clairement des définitions.
\end{proof}

\begin{lemma}
\label{lemma-add-summand-strict-epimorphism}
Soit $\mathcal{A}$ une catégorie abélienne.
Soit $f : B \to A$ un épimorphisme strict d'objets filtrés.
Soit $g : C \to A$ un morphisme d'objets filtrés.
Alors $f \oplus g : B \oplus C \to A$ est un épimorphisme strict.
\end{lemma}

\begin{proof}
Cela résulte clairement des définitions.
\end{proof}

\begin{lemma}
\label{lemma-induced-and-quotient-strict}
Soit $\mathcal{A}$ une catégorie abélienne.
Soient $(A, F)$, $(B, F)$ des objets filtrés.
Soit $u : A \to B$ un morphisme d'objets filtrés.
Si $u$ est injectif, alors $u$ est strict si et seulement si la filtration
sur $A$ est la filtration induite.
Si $u$ est surjectif, alors $u$ est strict si et seulement si la filtration
sur $B$ est la filtration quotient.
\end{lemma}

\begin{proof}
Cela résulte immédiatement de la définition.
\end{proof}

\begin{lemma}
\label{lemma-composition-strict}
Soit $\mathcal{A}$ une catégorie abélienne. Soient $f : A \to B$, $g : B \to C$
des morphismes stricts d'objets filtrés.
\begin{enumerate}
\item En général, la composée $g \circ f$ n'est pas stricte.
\item Si $g$ est injectif, alors $g \circ f$ est stricte.
\item Si $f$ est surjectif, alors $g \circ f$ est stricte.
\end{enumerate}
\end{lemma}

\begin{proof}
Soit $B$ un espace vectoriel sur un corps $k$, de base $e_1, e_2$, muni de la
filtration $F^nB = B$ pour $n < 0$, $F^0B = ke_1$, et $F^nB = 0$ pour
$n > 0$. Prenons $A = k(e_1 + e_2)$ et $C = B/ke_2$ avec les filtrations
induites par $B$, c'est-à-dire telles que $A \to B$ et $B \to C$ soient stricts
(lemme \ref{lemma-induced-and-quotient-strict}).
Alors $F^n(A) = A$ pour $n < 0$ et $F^n(A) = 0$ pour $n \geq 0$. 
De plus, $F^n(C) = C$ pour $n \leq 0$ et $F^n(C) = 0$ pour $n > 0$.
La composée (non nulle) $A \to C$ n'est donc pas stricte.

\medskip\noindent
Supposons $g$ injectif. Alors
\begin{align*}
g(f(F^pA)) & = g(f(A) \cap F^pB) \\
& = g(f(A)) \cap g(F^p(B)) \\
& = (g \circ f)(A) \cap (g(B) \cap F^pC) \\
& = (g \circ f)(A) \cap F^pC.
\end{align*}
La première égalité vient de ce que $f$ est strict, la deuxième de ce que $g$ est injectif,
la troisième de ce que $g$ est strict, et la quatrième de ce que
$(g \circ f)(A) \subset g(B)$.

\medskip\noindent
Supposons $f$ surjectif. Alors
\begin{align*}
(g \circ f)^{-1}(F^iC) & = f^{-1}(F^iB + \Ker(g)) \\
& = f^{-1}(F^iB) + f^{-1}(\Ker(g)) \\
& = F^iA + \Ker(f) + \Ker(g \circ f) \\
& = F^iA + \Ker(g \circ f)
\end{align*}
La première égalité vient de ce que $g$ est strict, la deuxième de ce que $f$ est
surjectif, la troisième de ce que $f$ est strict, et la dernière de ce que
$\Ker(f) \subset \Ker(g \circ f)$.
\end{proof}

\noindent
Le lemme suivant dit que les sous-objets d'un objet filtré possèdent une filtration
bien définie, indépendante du choix d'une présentation de l'objet comme
conoyau.

\begin{lemma}
\label{lemma-filtration-subobject}
Soit $\mathcal{A}$ une catégorie abélienne.
Soit $(A, F)$ un objet filtré de $\mathcal{A}$.
Soient $X \subset Y \subset A$ des sous-objets de $A$.
Sur l'objet
$$
Y/X = \Ker(A/X \to A/Y)
$$
la filtration quotient provenant de la filtration induite sur $Y$ et la
filtration induite provenant de la filtration quotient sur $A/X$ coïncident.
Chacun des morphismes $X \to Y$, $X \to A$, $Y \to A$, $Y \to A/X$,
$Y \to Y/X$, $Y/X \to A/X$ est strict (pour les filtrations induites/quotients).
\end{lemma}

\begin{proof}
La filtration quotient sur $Y/X$ est donnée par
$F^p(Y/X) = F^pY/(X \cap F^pY) = F^pY/F^pX$
car $F^pY = Y \cap F^pA$ et $F^pX = X \cap F^pA$.
La filtration induite par l'injection $Y/X \to A/X$ est donnée par
\begin{align*}
F^p(Y/X) & = Y/X \cap F^p(A/X) \\
& = Y/X \cap (F^pA + X)/X \\
& = (Y \cap F^pA)/(X \cap F^pA) \\
& = F^pY/F^pX.
\end{align*}
D'où la première assertion du lemme.
La démonstration des autres cas est analogue.
\end{proof}

\begin{lemma}
\label{lemma-pushout-filtered}
Soit $\mathcal{A}$ une catégorie abélienne.
Soient $A, B, C \in \text{Fil}(\mathcal{A})$.
Soient $f : A \to B$ et $g : A \to C$ des morphismes.
Il existe alors une somme amalgamée
$$
\xymatrix{
A \ar[r]_f \ar[d]_g & B \ar[d]^{g'} \\
C \ar[r]^{f'} & C \amalg_A B
}
$$
dans $\text{Fil}(\mathcal{A})$. Si $f$ est strict, alors $f'$ l'est aussi.
\end{lemma}

\begin{proof}
Posons $C \amalg_A B$ égal à $\Coker((g, -f) : A \to C \oplus B)$
dans $\text{Fil}(\mathcal{A})$. Ce conoyau existe d'après le
lemme \ref{lemma-filtered}.
C'est une somme amalgamée ; voir
l'exemple \ref{example-fibre-product-pushouts}.
Notons que $F^p(C \amalg_A B)$ est l'image de $F^pC \oplus F^pB$.
Ainsi
$$
(f')^{-1}(F^p(C \amalg_A B)) = g(f^{-1}(F^pB)) + F^pC
$$
D'où la dernière assertion.
\end{proof}

\begin{lemma}
\label{lemma-fibre-product-filtered}
Soit $\mathcal{A}$ une catégorie abélienne.
Soient $A, B, C \in \text{Fil}(\mathcal{A})$.
Soient $f : B \to A$ et $g : C \to A$ des morphismes.
Il existe alors un produit fibré
$$
\xymatrix{
B \times_A C \ar[r]_{g'} \ar[d]_{f'} & B \ar[d]^f \\
C \ar[r]^g & A
}
$$
dans $\text{Fil}(\mathcal{A})$. Si $f$ est strict, alors $f'$ l'est aussi.
\end{lemma}

\begin{proof}
Ce lemme est dual du
lemme \ref{lemma-pushout-filtered}.
\end{proof}

\noindent
Soit $\mathcal{A}$ une catégorie abélienne. Soit $(A, F)$ un objet
filtré de $\mathcal{A}$. On note $\text{gr}^p_F(A) = \text{gr}^p(A)$
l'objet $F^pA/F^{p + 1}A$ de $\mathcal{A}$. Cela définit un
foncteur additif
$$
\text{gr}^p :
\text{Fil}(\mathcal{A})
\longrightarrow
\mathcal{A}, \quad
(A, F)
\longmapsto
\text{gr}^p(A).
$$
Rappelons que nous avons défini la catégorie $\text{Gr}(\mathcal{A})$
des objets gradués de $\mathcal{A}$ dans la section \ref{section-graded}.
Pour $(A, F)$ dans $\text{Fil}(\mathcal{A})$, on peut poser
$$
\text{gr}(A) = \text{l'objet gradué de }\mathcal{A}\text{ dont la }
p\text{-ième composante graduée est }\text{gr}^p(A)
$$
et, si $\mathcal{A}$ admet les sommes directes dénombrables, on a simplement
$$
\text{gr}(A) = \bigoplus \text{gr}^p(A)
$$
Cela définit un foncteur additif
$$
\text{gr} :
\text{Fil}(\mathcal{A})
\longrightarrow
\text{Gr}(\mathcal{A}), \quad
(A, F) \longmapsto \text{gr}(A).
$$

\begin{lemma}
\label{lemma-ses-gr}
Soit $\mathcal{A}$ une catégorie abélienne.
\begin{enumerate}
\item Soient $A$ un objet filtré et $X \subset A$. Alors, pour tout $p$,
la suite
$$
0 \to \text{gr}^p(X) \to \text{gr}^p(A) \to \text{gr}^p(A/X) \to 0
$$
est exacte (pour la filtration induite sur $X$ et la filtration quotient sur $A/X$).
\item Soit $f : A \to B$ un morphisme d'objets filtrés de $\mathcal{A}$.
Alors, pour tout $p$, les suites
$$
0 \to \text{gr}^p(\Ker(f)) \to \text{gr}^p(A) \to
\text{gr}^p(\Coim(f)) \to 0
$$
et
$$
0 \to \text{gr}^p(\Im(f)) \to \text{gr}^p(B) \to
\text{gr}^p(\Coker(f)) \to 0
$$
sont exactes.
\end{enumerate}
\end{lemma}

\begin{proof}
On a $F^{p + 1}X = X \cap F^{p + 1}A$, donc l'application
$\text{gr}^p(X) \to \text{gr}^p(A)$ est injective. Dualement, l'application
$\text{gr}^p(A) \to \text{gr}^p(A/X)$ est surjective.
Le noyau de l'application
$F^pA/F^{p + 1}A \to (F^pA + X)/(F^{p + 1}A + X)$ est
$(F^{p + 1}A + (X \cap F^pA))/F^{p + 1}A \cong F^pX/F^{p + 1}X$,
d'où l'exactitude au milieu.
Les deux suites exactes courtes de (2) sont des cas particuliers de la
suite exacte courte de (1).
\end{proof}

\begin{lemma}
\label{lemma-characterize-strict}
Soit $\mathcal{A}$ une catégorie abélienne.
Soit $f : A \to B$ un morphisme d'objets
filtrés finis de $\mathcal{A}$. Les conditions suivantes sont équivalentes :
\begin{enumerate}
\item $f$ est strict,
\item le morphisme $\Coim(f) \to \Im(f)$ est un isomorphisme,
\item $\text{gr}(\Coim(f)) \to \text{gr}(\Im(f))$ est un
isomorphisme,
\item la suite
$\text{gr}(\Ker(f)) \to \text{gr}(A) \to \text{gr}(B)$
est exacte,
\item la suite $\text{gr}(A) \to \text{gr}(B) \to
\text{gr}(\Coker(f))$ est exacte, et
\item la suite
$$
0 \to
\text{gr}(\Ker(f)) \to
\text{gr}(A) \to
\text{gr}(B) \to
\text{gr}(\Coker(f)) \to 0
$$
est exacte.
\end{enumerate}
\end{lemma}

\begin{proof}
L'équivalence de (1) et (2) est le
lemme \ref{lemma-characterize-strict-general}.
D'après le
lemme \ref{lemma-ses-gr},
on voit que (4), (5), (6) impliquent (3) et que (3) implique (4), (5), (6).
Il suffit donc de montrer que (3) implique (2).
Il faut ainsi montrer que si $f : A \to B$ est une application injective et surjective
d'objets filtrés finis qui induit un isomorphisme
$\text{gr}(A) \to \text{gr}(B)$, alors $f$ induit un isomorphisme d'objets
filtrés. Autrement dit, il faut montrer que
$f(F^pA) = F^pB$ pour tout $p$.
Comme les filtrations sont finies, on peut le démontrer par récurrence descendante
sur $p$. Supposons que $f(F^{p + 1}A) = F^{p + 1}B$.
Alors le diagramme commutatif
$$
\xymatrix{
0 \ar[r] &
F^{p + 1}A \ar[r] \ar[d]^f &
F^pA \ar[r] \ar[d]^f &
\text{gr}^p(A) \ar[r] \ar[d]^{\text{gr}^p(f)} &
0 \\
0 \ar[r] &
F^{p + 1}B \ar[r] &
F^pB \ar[r] &
\text{gr}^p(B) \ar[r] &
0
}
$$
et le lemme des cinq impliquent que $f(F^pA) = F^pB$.
\end{proof}

\begin{lemma}
\label{lemma-filtered-complex}
Soit $\mathcal{A}$ une catégorie abélienne. Soit $A \to B \to C$ un complexe
d'objets filtrés de $\mathcal{A}$. Supposons que $\alpha : A \to B$ et
$\beta : B \to C$ soient des morphismes stricts d'objets filtrés. Alors
$\text{gr}(\Ker(\beta)/\Im(\alpha)) =
\Ker(\text{gr}(\beta))/\Im(\text{gr}(\alpha))$.
\end{lemma}

\begin{proof}
Cela résulte formellement du
lemme \ref{lemma-ses-gr}
et du fait que
$\Coim(\alpha) \cong \Im(\alpha)$ et
$\Coim(\beta) \cong \Im(\beta)$ d'après le
lemme \ref{lemma-characterize-strict-general}.
\end{proof}

\begin{lemma}
\label{lemma-filtered-acyclic}
Soit $\mathcal{A}$ une catégorie abélienne.
Soit $A \to B \to C$ un complexe d'objets filtrés de $\mathcal{A}$.
Supposons que $A, B, C$ aient des filtrations finies et que
$\text{gr}(A) \to \text{gr}(B) \to \text{gr}(C)$ soit exacte.
Alors :
\begin{enumerate}
\item pour tout $p \in \mathbf{Z}$, la suite
$\text{gr}^p(A) \to \text{gr}^p(B) \to \text{gr}^p(C)$ est exacte,
\item pour tout $p \in \mathbf{Z}$, la suite
$F^p(A) \to F^p(B) \to F^p(C)$ est exacte,
\item pour tout $p \in \mathbf{Z}$, la suite
$A/F^p(A) \to B/F^p(B) \to C/F^p(C)$ est exacte,
\item les applications $A \to B$ et $B \to C$ sont strictes, et
\item $A \to B \to C$ est exacte (comme suite dans $\mathcal{A}$).
\end{enumerate}
\end{lemma}

\begin{proof}
La partie (1) résulte immédiatement des définitions.
Démontrons (3) par récurrence sur la longueur des filtrations.
Si chacun de $A$, $B$, $C$ ne possède qu'une seule
partie graduée non nulle, alors (3) est vrai puisque $\text{gr}(A) = A$, etc.
Soit $n$ le plus grand entier tel que l'un au moins de
$F^nA, F^nB, F^nC$ soit non nul. Posons $A' = A/F^nA$, $B' = B/F^nB$,
$C' = C/F^nC$, avec les filtrations induites. Notons que
$\text{gr}(A) = F^nA \oplus \text{gr}(A')$
et de même pour $B$ et $C$. L'hypothèse de récurrence
s'applique à $A' \to B' \to C'$, ce qui implique que
$A/F^p(A) \to B/F^p(B) \to C/F^p(C)$ est exacte pour $p \leq n$.
Pour conclure de même pour $p = n + 1$, c'est-à-dire pour démontrer que $A \to B \to C$
est exacte, on utilise le diagramme commutatif
$$
\xymatrix{
0 \ar[r] & F^nA \ar[r] \ar[d] & A \ar[r] \ar[d] & A' \ar[r] \ar[d] & 0 \\
0 \ar[r] & F^nB \ar[r] \ar[d] & B \ar[r] \ar[d] & B' \ar[r] \ar[d] & 0 \\
0 \ar[r] & F^nC \ar[r] & C \ar[r] & C' \ar[r] & 0
}
$$
dont les lignes sont des suites exactes courtes d'objets de $\mathcal{A}$.
La démonstration de (2) est duale. Bien entendu, (5) résulte de (2).

\medskip\noindent
Pour démontrer (4), notons $f : A \to B$ et $g : B \to C$ les morphismes donnés.
On sait par (2) que $f(F^p(A)) = \Ker(F^p(B) \to F^p(C))$ et,
par (5), que $f(A) = \Ker(g)$. Par conséquent,
$f(F^p(A)) =  \Ker(F^p(B) \to F^p(C)) =
\Ker(g) \cap F^p(B) = f(A) \cap F^p(B)$, ce qui démontre que
$f$ est strict. La démonstration que $g$ est strict est duale.
\end{proof}








\section{Suites spectrales}
\label{section-spectral-sequence}

\noindent
On trouvera une présentation agréable des suites spectrales dans
\cite{Eisenbud}. Voir aussi \cite{McCleary}, \cite{Lang}, etc.

\begin{definition}
\label{definition-spectral-sequence}
Soit $\mathcal{A}$ une catégorie abélienne.
\begin{enumerate}
\item Une {\it suite spectrale dans $\mathcal{A}$} est la donnée d'un
système $(E_r, d_r)_{r \geq 1}$ où chaque $E_r$ est un objet
de $\mathcal{A}$ et chaque $d_r : E_r \to E_r$ est un morphisme tel
que $d_r \circ d_r = 0$ et $E_{r + 1} = \Ker(d_r)/\Im(d_r)$
pour $r \geq 1$.
\item Un {\it morphisme de suites spectrales}
$f : (E_r, d_r)_{r \geq 1} \to (E'_r, d'_r)_{r \geq 1}$ est
la donnée d'une famille de morphismes $f_r : E_r \to E'_r$ tels que
$f_r \circ d_r = d'_r \circ f_r$ et tels que $f_{r + 1}$
soit le morphisme induit par $f_r$ au moyen des identifications
$E_{r + 1} = \Ker(d_r)/\Im(d_r)$
et
$E'_{r + 1} = \Ker(d'_r)/\Im(d'_r)$.
\end{enumerate}
\end{definition}

\noindent
Nous assouplirons parfois quelque peu cette définition en autorisant $E_{r + 1}$
à être un objet muni d'un isomorphisme donné
$E_{r + 1} \to \Ker(d_r)/\Im(d_r)$.
Il arrive également que l'on ait un système $(E_r, d_r)_{r \geq r_0}$
pour un certain $r_0 \in \mathbf{Z}$, satisfaisant les propriétés de la définition ci-dessus
pour les indices $\geq r_0$. Nous l'appellerons aussi une suite spectrale puisque,
par une simple renumérotation, il relève de toute façon de la définition.
En fait, les cas $r_0 = 0$ et $r_0 = -1$ se rencontrent dans la littérature.

\medskip\noindent
Étant donnée une suite spectrale $(E_r, d_r)_{r \geq 1}$, on définit
$$
0 = B_1 \subset B_2 \subset \ldots \subset B_r \subset \ldots
\subset Z_r \subset \ldots \subset Z_2 \subset Z_1 = E_1
$$
par le procédé simple suivant. Posons $B_2 = \Im(d_1)$
et $Z_2 = \Ker(d_1)$. Il est alors clair que
$d_2 : Z_2/B_2 \to Z_2/B_2$. On peut donc définir $B_3$ comme l'unique
sous-objet de $E_1$ contenant $B_2$ tel que $B_3/B_2$ soit l'image
de $d_2$. De même, on peut définir $Z_3$ comme l'unique sous-objet de
$E_1$ contenant $B_2$ tel que $Z_3/B_2$ soit le noyau de $d_2$.
Et ainsi de suite. En particulier, on a
$$
E_r = Z_r/B_r
$$
pour tout $r \geq 1$. Si la suite spectrale commence en $r = r_0$,
on peut construire de même $B_i$, $Z_i$ comme sous-objets de $E_{r_0}$.
En fait, on rencontre parfois dans la littérature la notation
$$
0 = B_r(E_r) \subset B_{r + 1}(E_r) \subset B_{r + 2}(E_r) \subset \ldots
\subset Z_{r + 2}(E_r) \subset Z_{r + 1}(E_r) \subset Z_r(E_r) = E_r
$$
pour désigner la filtration décrite ci-dessus, mais en partant de $E_r$.

\begin{definition}
\label{definition-limit-spectral-sequence}
Soit $\mathcal{A}$ une catégorie abélienne.
Soit $(E_r, d_r)_{r \geq 1}$ une suite spectrale.
\begin{enumerate}
\item Si les sous-objets $Z_{\infty} = \bigcap Z_r$
et $B_{\infty} = \bigcup B_r$ de $E_1$ existent, on définit
la {\it limite}\footnote{Cette notation n'est pas universellement admise. Dans certaines
références, une paire supplémentaire de sous-objets
$Z_\infty$ et $B_\infty$ de $E_1$ tels que
$0 = B_1 \subset B_2 \subset \ldots \subset B_\infty \subset Z_\infty
\subset \ldots \subset Z_2 \subset Z_1 = E_1$
fait partie des données constituant une suite spectrale !}
de la suite spectrale comme étant l'objet
$E_{\infty} = Z_{\infty}/B_{\infty}$.
\item On dit que la suite spectrale {\it dégénère en $E_r$}
si les différentielles $d_r, d_{r + 1}, \ldots$ sont toutes nulles.
\end{enumerate}
\end{definition}

\noindent
Remarquons que si la suite spectrale dégénère en $E_r$, alors
on a $E_r = E_{r + 1} = \ldots = E_{\infty}$ (et la limite
existe bien entendu). En outre, presque toutes les catégories abéliennes que nous rencontrerons
possèdent des sommes et des intersections dénombrables.

\begin{remark}[Variante]
\label{remark-allow-translation-functors}
Il arrive souvent que les termes d'une suite spectrale possèdent une
structure supplémentaire, par exemple une graduation ou une bigraduation.
Pour prendre cela en compte (et contourner certaines difficultés techniques),
nous introduisons la notion suivante. Soit $\mathcal{A}$ une
catégorie abélienne. Soit $(T_r)_{r \geq 1}$ une
suite de foncteurs de {\it translation} ou de {\it décalage}, c'est-à-dire que
$T_r : \mathcal{A} \to \mathcal{A}$ est un isomorphisme de catégories.
Dans ce cadre, une {\it suite spectrale} est la donnée d'un système
$(E_r, d_r)_{r \geq 1}$ où chaque $E_r$ est un objet de
$\mathcal{A}$ et chaque $d_r : E_r \to T_rE_r$
est un morphisme tel que $T_rd_r \circ d_r = 0$, de sorte que
$$
\xymatrix{
\ldots \ar[r] &
T_r^{-1}E_r \ar[r]^-{T_r^{-1}d_r} &
E_r \ar[r]^-{d_r} &
T_rE_r \ar[r]^{T_r d_r} &
T_r^2E_r \ar[r] & \ldots
}
$$
soit un complexe et que $E_{r + 1} = \Ker(d_r)/\Im(T_r^{-1}d_r)$ pour $r \geq 1$.
La notion de {\it morphisme de suites spectrales}
dans ce cadre est claire. On peut encore y définir
$$
0 = B_1 \subset B_2 \subset \ldots \subset B_r \subset \ldots
\subset Z_r \subset \ldots \subset Z_2 \subset Z_1 = E_1
$$
ainsi que $Z_\infty$ et $B_\infty$ (s'ils existent), comme ci-dessus.
\end{remark}










\section{Suites spectrales : couples exacts}
\label{section-exact-couple}

\begin{definition}
\label{definition-exact-couple}
Soit $\mathcal{A}$ une catégorie abélienne.
\begin{enumerate}
\item Un {\it couple exact} est une donnée $(A, E, \alpha, f, g)$ où
$A$, $E$ sont des objets de $\mathcal{A}$ et $\alpha$, $f$, $g$
sont des morphismes comme dans le diagramme suivant :
$$
\xymatrix{
A \ar[rr]_{\alpha} & & A \ar[ld]^g \\
& E \ar[lu]^f &
}
$$
tels que le noyau de chaque flèche soit l'image
de la précédente. Ainsi, $\Ker(\alpha) = \Im(f)$,
$\Ker(f) = \Im(g)$ et $\Ker(g) = \Im(\alpha)$.
\item Un {\it morphisme de couples exacts}
$t : (A, E, \alpha, f, g) \to (A', E', \alpha', f', g')$
est la donnée de morphismes $t_A : A \to A'$ et
$t_E : E \to E'$ tels que
$\alpha' \circ t_A = t_A \circ \alpha$,
$f' \circ t_E = t_A \circ f$ et
$g' \circ t_A = t_E \circ g$.
\end{enumerate}
\end{definition}

\begin{lemma}
\label{lemma-derived-exact-couple}
Soit $(A, E, \alpha, f, g)$ un couple exact dans une catégorie abélienne
$\mathcal{A}$. Posons :
\begin{enumerate}
\item $d = g \circ f : E \to E$, de sorte que $d \circ d = 0$,
\item $E' = \Ker(d)/\Im(d)$,
\item $A' = \Im(\alpha)$,
\item $\alpha' : A' \to A'$ induit par $\alpha$,
\item $f' : E' \to A'$ induit par $f$,
\item $g' : A' \to E'$ induit par « $g \circ \alpha^{-1}$ ».
\end{enumerate}
Alors :
\begin{enumerate}
\item $\Ker(d) = f^{-1}(\Ker(g)) = f^{-1}(\Im(\alpha))$,
\item $\Im(d) = g(\Im(f)) = g(\Ker(\alpha))$,
\item $(A', E', \alpha', f', g')$ est un couple exact.
\end{enumerate}
\end{lemma}

\begin{proof}
Omis.
\end{proof}

\noindent
Il est donc clair qu'un couple exact $(A, E, \alpha, f, g)$
donne une suite spectrale en posant $E_1 = E$, $d_1 = d$,
$E_2 = E'$, $d_2 = d' = g' \circ f'$, $E_3 = E''$, $d_3 = d'' = g'' \circ f''$,
et ainsi de suite.

\begin{definition}
\label{definition-spectral-sequence-associated-exact-couple}
Soit $\mathcal{A}$ une catégorie abélienne.
Soit $(A, E, \alpha, f, g)$ un couple exact.
La {\it suite spectrale associée au couple exact}
est la suite spectrale $(E_r, d_r)_{r \geq 1}$ avec
$E_1 = E$, $d_1 = d$, $E_2 = E'$, $d_2 = d' = g' \circ f'$,
$E_3 = E''$, $d_3 = d'' = g'' \circ f''$,
et ainsi de suite.
\end{definition}

\begin{lemma}
\label{lemma-spectral-sequence-associated-exact-couple}
Soit $\mathcal{A}$ une catégorie abélienne.
Soit $(A, E, \alpha, f, g)$ un couple exact.
Soit $(E_r, d_r)_{r \geq 1}$ la suite spectrale
associée au couple exact.
Dans ce cas, on a
$$
0 = B_1 \subset \ldots \subset
B_{r + 1} = g(\Ker(\alpha^r))
\subset \ldots \subset
Z_{r + 1} = f^{-1}(\Im(\alpha^r))
\subset \ldots \subset Z_1 = E
$$
et l'application $d_{r + 1} : E_{r + 1} \to E_{r + 1}$
est décrite par la règle suivante :
pour tout objet (test) $T$ de $\mathcal{A}$ et tous éléments
$x : T \to Z_{r + 1}$ et $y : T \to A$ tels que
$f \circ x = \alpha^r \circ y$, on a
$$
d_{r + 1} \circ \overline{x} = \overline{g \circ y}
$$
où $\overline{x} : T \to E_{r + 1}$ est le
morphisme induit.
\end{lemma}

\begin{proof}
Omis.
\end{proof}

\noindent
Remarquons que, dans la situation du lemme, on a manifestement
$$
B_\infty = g\left(\bigcup\nolimits_r \Ker(\alpha^r)\right)
\subset
Z_\infty = f^{-1}\left(\bigcap\nolimits_r \Im(\alpha^r)\right)
$$
pourvu que $\bigcup \Ker(\alpha^r)$ et $\bigcap \Im(\alpha^r)$ existent.
Cela donne pour limite $E_\infty = Z_\infty / B_\infty$, voir la
définition \ref{definition-limit-spectral-sequence}.

\begin{remark}[Variante]
\label{remark-shifted-exact-couple}
Soit $\mathcal{A}$ une catégorie abélienne. Soient
$S, T : \mathcal{A} \to \mathcal{A}$ des foncteurs de
décalage, c'est-à-dire des isomorphismes de catégories. Nous noterons
leurs composées $n$ fois par $S^nA$ et $T^nA$, pour
$A \in \Ob(\mathcal{A})$ et $n \in \mathbf{Z}$.
Dans cette situation, un {\it couple exact} est une donnée $(A, E, \alpha, f, g)$
où $A$, $E$ sont des objets de $\mathcal{A}$ et $\alpha : A \to T^{-1}A$,
$f : E \to A$, $g : A \to SE$ sont des morphismes tels que
$$
\xymatrix{
TE \ar[r]^-{Tf} &
TA \ar[r]^-{T\alpha} &
A \ar[r]^-{g} &
SE \ar[r]^{Sf} & SA
}
$$
soit un complexe exact. Représentons cela comme suit :
$$
\xymatrix{
TA \ar[rrr]_{T\alpha} & & &
A \ar[ld]^g \ar[rrr]_\alpha & & &
T^{-1}A \ar[ld]^{T^{-1}g} \\
& TE \ar[lu]^{Tf} \ar@{..}[r] & SE & &
E \ar[lu]^f \ar@{..}[r] & T^{-1}SE
}
$$
Posons $d = g \circ f : E \to SE$. Alors $d \circ S^{-1}d =
g \circ f \circ S^{-1}g \circ S^{-1}f = 0$, car $f \circ S^{-1}g = 0$.
Posons $E' = \Ker(d)/\Im(S^{-1}d)$. Posons $A' = \Im(T\alpha)$.
Soit $\alpha' : A' \to T^{-1}A'$ le morphisme induit par $\alpha$.
Soit $f' : E' \to A'$ induit par $f$, ce qui est possible car
$f(\Ker(d)) \subset \Ker(g) = \Im(T\alpha)$.
Enfin, soit $g' : A' \to TSE'$ le morphisme induit par
« $Tg \circ (T\alpha)^{-1}$ »\footnote{Cela fonctionne car
$TSE' = \Ker(TSd)/\Im(Td)$ et
$Tg(\Ker(T\alpha)) = Tg(\Im(Tf)) = \Im(T(d))$
et $TS(d)(\Im(Tg)) = \Im(TSg \circ TSf \circ Tg) = 0$.}.

\medskip\noindent
Exactement comme ci-dessus, on obtient
\begin{enumerate}
\item $\Ker(d) = f^{-1}(\Ker(g)) = f^{-1}(\Im(T\alpha))$,
\item $\Im(d) = g(\Im(f)) = g(\Ker(\alpha))$,
\item $(A', E', \alpha', f', g')$ est un couple exact
pour les foncteurs de décalage $TS$ et $T$.
\end{enumerate}
On obtient une suite spectrale
(comme dans la remarque \ref{remark-allow-translation-functors})
avec $E_1 = E$, $E_2 = E'$, etc., et $d_r : E_r \to T^{r - 1}SE_r$
pour tout $r \geq 1$. Le lemme \ref{lemma-spectral-sequence-associated-exact-couple}
nous dit que
$$
SB_{r + 1} =
g(\Ker(T^{-r + 1}\alpha \circ \ldots \circ T^{-1}\alpha \circ \alpha))
$$
et
$$
Z_{r + 1} = f^{-1}(\Im(T\alpha \circ T^2\alpha \circ \ldots \circ T^r\alpha))
$$
dans cette situation. La description de l'application $d_{r + 1}$ est analogue
à celle donnée dans le lemme. (Il peut être plus facile d'utiliser ces descriptions
explicites pour démontrer que l'on obtient une suite spectrale à partir d'un tel
couple exact.)
\end{remark}

\section{Suites spectrales : objets différentiels}
\label{section-differential-object}

\begin{definition}
\label{definition-differential-object}
Soit $\mathcal{A}$ une catégorie abélienne.
Un {\it objet différentiel} de $\mathcal{A}$
est un couple $(A, d)$ constitué d'un
objet $A$ de $\mathcal{A}$
muni d'un endomorphisme $d$ tel que $d \circ d = 0$.
Un {\it morphisme d'objets différentiels} $(A, d) \to (B, d)$
est la donnée d'un morphisme $\alpha : A \to B$ tel que
$d \circ \alpha = \alpha \circ d$.
\end{definition}

\begin{lemma}
\label{lemma-differential-objects-abelian}
\begin{slogan}
La catégorie des objets différentiels d'une catégorie abélienne est elle-même
une catégorie abélienne.
\end{slogan}
Soit $\mathcal{A}$ une catégorie abélienne.
La catégorie des objets différentiels de $\mathcal{A}$ est abélienne.
\end{lemma}

\begin{proof}
Omis.
\end{proof}

\begin{definition}
\label{definition-differential-object-homology}
Pour un objet différentiel $(A, d)$, on note
$$
H(A, d) = \Ker(d)/\Im(d)
$$
son {\it homologie}.
\end{definition}

\begin{lemma}
\label{lemma-differential-objects-ses}
Soit $\mathcal{A}$ une catégorie abélienne.
Soit $0 \to (A, d) \to (B, d) \to (C, d) \to 0$ une suite exacte courte
d'objets différentiels. On obtient alors une suite exacte d'homologie
$$
\ldots \to H(C, d) \to H(A, d) \to H(B, d) \to H(C, d) \to \ldots
$$
\end{lemma}

\begin{proof}
Appliquons le lemme \ref{lemma-long-exact-sequence-cochain}
à la suite exacte courte de complexes
$$
\begin{matrix}
0 & \to & A & \to & B & \to & C & \to & 0 \\
& & \downarrow & & \downarrow & & \downarrow \\
0 & \to & A & \to & B & \to & C & \to & 0 \\
& & \downarrow & & \downarrow & & \downarrow \\
0 & \to & A & \to & B & \to & C & \to & 0
\end{matrix}
$$
où les flèches verticales sont $d$.
\end{proof}

\noindent
Nous arrivons à un exemple important de suite spectrale.
Soit $\mathcal{A}$ une catégorie abélienne.
Soit $(A, d)$ un objet différentiel de $\mathcal{A}$.
Soit $\alpha : (A, d) \to (A, d)$ un endomorphisme de cet objet
différentiel. Si l'on suppose $\alpha$ injectif, on obtient une suite exacte courte
$$
0 \to (A, d) \to (A, d) \to (A/\alpha A, d) \to 0
$$
d'objets différentiels. D'après le
lemme \ref{lemma-differential-objects-ses}, on obtient un couple exact
$$
\xymatrix{
H(A, d) \ar[rr]_{\overline{\alpha}} & & H(A, d) \ar[ld]^g \\
& H(A/\alpha A, d) \ar[lu]^f &
}
$$
où $g$ est l'application canonique et $f$ l'application définie dans le lemme du serpent.
La section \ref{section-exact-couple} fournit une suite spectrale associée.
Puisque, dans ce cas, $E_1 = H(A/\alpha A, d)$,
il est naturel de définir $E_0 = A/\alpha A$ et $d_0 = d$.
Autrement dit, on fait commencer la suite spectrale
en $r = 0$. Selon les conventions de la
section \ref{section-spectral-sequence}, on définit une suite de sous-objets
$$
0 = B_0 \subset \ldots \subset B_r \subset \ldots
\subset Z_r \subset \ldots \subset Z_0 = E_0
$$
telle que $E_r = Z_r/B_r$. Il résulte du
lemme \ref{lemma-spectral-sequence-associated-exact-couple}
que, pour $r \geq 1$, on a :
\begin{enumerate}
\item $B_r$ est l'image de $(\alpha^{r - 1})^{-1}(d A)$
par l'application naturelle $A \to A/\alpha A$,
\item $Z_r$ est l'image de $d^{-1}(\alpha^r A)$
par l'application naturelle $A \to A/\alpha A$, et
\item $d_r : E_r \to E_r$ est donnée comme suit : étant donné un élément
$\overline{z} \in Z_r$, choisissons un représentant
$z \in d^{-1}(\alpha^r A)$, puis un élément $y \in A$ tel que
$d(z) = \alpha^r(y)$, et notons $\overline{y}$ son image dans
$A/\alpha A$. Alors
$d_r(\overline{z} + B_r) = \overline{y} + B_r$.
\end{enumerate}
Attention : il n'est pas nécessairement vrai que
$\alpha A \subset (\alpha^{r - 1})^{-1}(dA)$, ni que
$\alpha A \subset d^{-1}(\alpha^r A)$. En revanche, on a bien
$(\alpha^{r - 1})^{-1}(dA) \subset d^{-1}(\alpha^r A)$.
On a
$$
E_r
=
\frac{d^{-1}(\alpha^r A) + \alpha A}{(\alpha^{r - 1})^{-1}(dA) + \alpha A}.
$$
Il n'est pas difficile de vérifier directement que (1) -- (3) définissent une suite spectrale.

\begin{definition}
\label{definition-differential-object-selfmap}
Soit $\mathcal{A}$ une catégorie abélienne.
Soit $(A, d)$ un objet différentiel de $\mathcal{A}$.
Soit $\alpha : A \to A$ un endomorphisme injectif de $A$ qui
commute à $d$. La {\it suite spectrale associée à
$(A, d, \alpha)$} est la suite spectrale
$(E_r, d_r)_{r \geq 0}$ décrite ci-dessus.
\end{definition}

\begin{remark}[Variante]
\label{remark-differential-object-selfmap}
Soit $\mathcal{A}$ une catégorie abélienne et soient
$S, T : \mathcal{A} \to \mathcal{A}$ des foncteurs de décalage, c'est-à-dire des
isomorphismes de catégories. Supposons que $TS = ST$ comme foncteurs.
Considérons les couples $(A, d)$ constitués d'un objet $A$ de $\mathcal{A}$
et d'un morphisme $d : A \to SA$ tel que $d \circ S^{-1}d = 0$.
La catégorie de ces objets est abélienne.
On définit $H(A, d) = \Ker(d)/\Im(S^{-1}d)$ et l'on observe que
$H(SA, Sd) = SH(A, d)$ (isomorphisme canonique).
Étant donnée une suite exacte courte
$$
0 \to (A, d) \to (B, d) \to (C, d) \to 0
$$
on obtient une suite exacte longue d'homologie
$$
\ldots \to S^{-1}H(C, d) \to
H(A, d) \to H(B, d) \to H(C, d) \to SH(A, d) \to \ldots
$$
(remarquer les décalages dans les morphismes de bord). Puisque $ST = TS$, le foncteur
$T$ définit un foncteur de décalage sur les couples en posant $T(A, d) = (TA, Td)$.
Soit ensuite $\alpha : (A, d) \to T^{-1}(A, d)$ injectif, de
conoyau $(Q, d)$. On obtient alors un couple exact comme dans la
remarque \ref{remark-shifted-exact-couple}, avec les foncteurs de décalage $TS$ et
$T$, donné par
$$
(H(A, d), S^{-1}H(Q, d), \overline{\alpha}, f, g)
$$
Ici, $\overline{\alpha} : H(A, d) \to T^{-1}H(A, d)$ est induit par $\alpha$,
l'application $f : S^{-1}H(Q, d) \to H(A, d)$ est le morphisme
de bord, et $g : H(A, d) \to TH(Q, d) = TS(S^{-1}H(Q, d))$ est induit par
l'application quotient $A \to TQ$. On obtient ainsi une suite spectrale comme ci-dessus,
avec $E_1 = S^{-1}H(Q, d)$ et les différentielles $d_r : E_r \to T^rSE_r$.
Comme ci-dessus, on pose $E_0 = S^{-1}Q$ et $d_0 : E_0 \to SE_0$, donné par
$S^{-1}d : S^{-1}Q \to Q$. Si, selon nos conventions, on définit
$B_r \subset Z_r \subset E_0$, alors, pour $r \geq 1$, on a :
\begin{enumerate}
\item $SB_r$ est l'image de
$$
(T^{-r + 1}\alpha \circ \ldots \circ T^{-1}\alpha)^{-1}
\Im(T^{-r}S^{-1}d)
$$
par l'application naturelle $T^{-1}A \to Q$,
\item $Z_r$ est l'image de
$$
(S^{-1}T^{-1}d)^{-1}
\Im(\alpha \circ \ldots \circ T^{r - 1}\alpha)
$$
par l'application naturelle $S^{-1}T^{-1}A \to S^{-1}Q$.
\end{enumerate}
Les différentielles peuvent être décrites comme suit : si $x \in Z_r$, on
choisit $x' \in S^{-1}T^{-1}A$ qui s'envoie sur $x$. Alors $S^{-1}T^{-1}d(x')$
est égal à $(\alpha \circ \ldots \circ T^{r - 1}\alpha)(y)$ pour un certain
$y \in T^{r - 1}A$. Alors $d_r(x) \in T^rSE_r$ est représenté par la
classe de l'image de $y$ dans $T^rSE_0 = T^rQ$ modulo $T^rSB_r$.
\end{remark}
\section{Suites spectrales : objets différentiels filtrés}
\label{section-filtered-differential}

\noindent
On peut construire une suite spectrale à partir d'un
objet différentiel filtré.

\begin{definition}
\label{definition-filtered-differential}
Soit $\mathcal{A}$ une catégorie abélienne.
Un {\it objet différentiel filtré} $(K, F, d)$ est un objet filtré
$(K, F)$ de $\mathcal{A}$ muni d'un endomorphisme
$d : (K, F) \to (K, F)$ dont le carré est nul : $d \circ d = 0$.
\end{definition}

\noindent
Pour décrire la suite spectrale associée à un tel objet,
supposons pour l'instant que $\mathcal{A}$ soit une catégorie abélienne
qui possède des sommes directes dénombrables et que celles-ci soient exactes
(ce n'est pas automatique, voir la remarque \ref{remark-direct-sums-not-exact}).
Soit $(K, F, d)$ un objet différentiel filtré de $\mathcal{A}$.
Remarquons que chaque $F^nK$ est lui-même un objet différentiel.
Considérons l'objet $A = \bigoplus F^nK$ et munissons-le d'une
différentielle $d$ en utilisant $d$ sur chaque facteur.
Alors $(A, d)$ est un objet différentiel de $\mathcal{A}$
muni d'une graduation. Considérons l'application
$$
\alpha : A \to A
$$
donnée par les inclusions $F^nK \to F^{n - 1}K$.
C'est manifestement un morphisme injectif d'objets différentiels
$\alpha : (A, d) \to (A, d)$. Par conséquent, la
définition \ref{definition-differential-object-selfmap}
fournit une suite spectrale.
Nous l'appellerons {\it la suite spectrale associée à
l'objet différentiel filtré $(K, F, d)$}.

\medskip\noindent
Déterminons les termes de cette suite spectrale.
Remarquons d'abord que $A/\alpha A = \text{gr}(K)$,
muni de sa différentielle $d = \text{gr}(d)$. On voit donc que
$$
E_0 = \text{gr}(K), \quad d_0 = \text{gr}(d).
$$
L'homologie de l'objet différentiel gradué $\text{gr}(K)$
est donc le terme suivant :
$$
E_1 = H(\text{gr}(K), \text{gr}(d)).
$$
En outre, $E_0$ est un objet gradué de $\mathcal{A}$
et $d_0$ est compatible avec la graduation. Il est donc clair que $E_1$
est lui aussi un objet gradué. Il se trouve cependant que la différentielle
$d_1$ ne préserve pas cette graduation : elle décale le degré de $1$.

\medskip\noindent
Pour préciser cela, définissons
$$
Z_r^p =
\frac{F^pK \cap d^{-1}(F^{p + r}K) + F^{p + 1}K}{F^{p + 1}K}
$$
et
$$
B_r^p =
\frac{F^pK \cap d(F^{p - r + 1}K) + F^{p + 1}K}{F^{p + 1}K}.
$$
Cette notation, quoique fort naturelle, semble différer de celle employée
dans la plupart des ouvrages. Cela n'a peut-être guère d'importance,
puisque la littérature ne semble pas non plus adopter de notation cohérente.
Avec ces choix, $B_r \subset E_0$,
resp.\ $Z_r \subset E_0$ (tels qu'ils sont définis dans la
section \ref{section-differential-object}), est égal à
$\bigoplus_p B_r^p$, resp.\ $\bigoplus_p Z_r^p$.
Ainsi, si l'on définit
$$
E_r^p = Z_r^p/B_r^p
$$
pour $r \geq 0$ et $p \in \mathbf{Z}$, alors $E_r = \bigoplus_p E_r^p$.
On peut définir une différentielle $d_r^p : E_r^p \to E_r^{p + r}$
par la règle
$$
z + F^{p + 1}K
\longmapsto
dz + F^{p + r + 1}K
$$
où $z \in F^pK \cap d^{-1}(F^{p + r}K)$.

\begin{lemma}
\label{lemma-spectral-sequence-filtered-differential}
Soit $\mathcal{A}$ une catégorie abélienne. Soit $(K, F, d)$ un
objet différentiel filtré de $\mathcal{A}$. Il existe une
suite spectrale $(E_r, d_r)_{r \geq 0}$ dans $\text{Gr}(\mathcal{A})$
associée à $(K, F, d)$ telle que $d_r : E_r \to E_r[r]$
pour tout $r$, et telle que les composantes graduées
$E_r^p$ et les applications $d_r^p : E_r^p \to E_r^{p + r}$
soient celles données ci-dessus. De plus, $E_0^p = \text{gr}^p K$,
$d_0^p = \text{gr}^p(d)$ et $E_1^p = H(\text{gr}^pK, \text{gr}^p(d))$.
\end{lemma}

\begin{proof}
Si $\mathcal{A}$ possède des sommes directes dénombrables et si celles-ci
sont exactes, cela résulte de la discussion ci-dessus.
Dans le cas général, procédons comme suit ; nous conseillons vivement au lecteur
d'omettre cette preuve. Considérons l'objet $A = (F^{p + 1}K)$ de
$\text{Gr}(\mathcal{A})$, c'est-à-dire que l'on place $F^{p + 1}K$ en degré $p$
(le décalage inhabituel de la numérotation donnera plus tard la bonne numérotation).
Munissons-le d'une différentielle $d$ en utilisant $d$ sur chaque composante.
Alors $(A, d)$ est un objet différentiel de $\text{Gr}(\mathcal{A})$.
Considérons l'application
$$
\alpha : A \to A[-1]
$$
donnée en degré $p$ par les inclusions $F^{p + 1}K \to F^pK$.
C'est manifestement un morphisme injectif d'objets différentiels
$\alpha : (A, d) \to (A, d)[-1]$. On peut donc appliquer la
remarque \ref{remark-differential-object-selfmap}
avec $S = \text{id}$ et $T = [1]$.
La suite spectrale correspondante $(E_r, d_r)_{r \geq 0}$
dans $\text{Gr}(\mathcal{A})$ est celle que nous cherchons.
Développons quelque peu les définitions.
Tout d'abord, $E_r = (E_r^p)$ est un objet de $\text{Gr}(\mathcal{A})$.
Ensuite, puisque $T^rS = [r]$, on a $d_r : E_r \to E_r[r]$, ce qui signifie que
$d_r^p : E_r^p \to E_r^{p + r}$.

\medskip\noindent
Pour vérifier la description des composantes graduées, raisonnons
comme ci-dessus. On a d'abord $E_0 = \Coker(\alpha : A \to A[-1])$,
et notre choix de numérotation donne
$E_0^p = \text{gr}^pK$. La première différentielle est donnée par
$d_0^p = \text{gr}^pd : E_0^p \to E_0^p$.
Ensuite, la description des bords $B_r$ et des cocycles $Z_r$
dans la remarque \ref{remark-differential-object-selfmap}
se traduit directement par les formules
pour $Z_r^p$ et $B_r^p$ données ci-dessus.
\end{proof}

\begin{lemma}
\label{lemma-spectral-sequence-filtered-differential-d1}
Soit $\mathcal{A}$ une catégorie abélienne. Soit $(K, F, d)$ un objet
différentiel filtré de $\mathcal{A}$. La suite spectrale
$(E_r, d_r)_{r \geq 0}$ associée à $(K, F, d)$ possède
$$
d_1^p :
E_1^p = H(\text{gr}^pK)
\longrightarrow
H(\text{gr}^{p + 1}K) = E_1^{p + 1}
$$
égal au morphisme de bord en homologie associé à la suite
exacte courte d'objets différentiels
$$
0 \to \text{gr}^{p + 1}K \to F^pK/F^{p + 2}K \to \text{gr}^pK \to 0.
$$
\end{lemma}

\begin{proof}
Cela résulte immédiatement de la formule de la différentielle $d_1^p$
donnée juste avant le lemme \ref{lemma-spectral-sequence-filtered-differential}.
\end{proof}

\begin{definition}
\label{definition-filtration-cohomology-filtered-differential}
Soit $\mathcal{A}$ une catégorie abélienne.
Soit $(K, F, d)$ un objet différentiel filtré de $\mathcal{A}$.
La {\it filtration induite} sur $H(K, d)$ est la filtration définie
par $F^pH(K, d) = \Im(H(F^pK, d) \to H(K, d))$.
\end{definition}

\noindent
En explicitant cette définition, on voit que
$$
F^pH(K, d) =
\frac{\Ker(d) \cap F^pK + \Im(d)}{\Im(d)}
$$
et donc que
$$
\text{gr}^p H(K) =
\frac{\Ker(d) \cap F^pK + \Im(d)}{\Ker(d) \cap F^{p + 1}K + \Im(d)} =
\frac{\Ker(d) \cap F^pK}{\Ker(d) \cap F^{p + 1}K + \Im(d) \cap F^pK}
$$

\begin{lemma}
\label{lemma-compute-filtered-cohomology}
Soit $\mathcal{A}$ une catégorie abélienne. Soit $(K, F, d)$ un objet
différentiel filtré de $\mathcal{A}$. Si $Z_\infty^p$ et $B_\infty^p$
existent (voir la preuve), alors :
\begin{enumerate}
\item la limite $E_\infty$ existe et est graduée, avec
$E_\infty^p = Z_\infty^p/B_\infty^p$ en degré $p$, et
\item l'objet gradué associé $\text{gr}(H(K))$ de l'homologie de $K$
est un sous-quotient gradué de l'objet limite gradué $E_\infty$.
\end{enumerate}
\end{lemma}

\begin{proof}
Les objets $Z_\infty$, $B_\infty$ et la limite
$E_\infty = Z_\infty/B_\infty$ de la
définition \ref{definition-limit-spectral-sequence}
sont des objets de $\text{Gr}(\mathcal{A})$ par notre construction de
la suite spectrale dans la preuve du
lemme \ref{lemma-spectral-sequence-filtered-differential}.
Puisque $Z_r^p$, resp.\ $B_r^p$, est la composante graduée de degré $p$ de $Z_r$,
resp.\ de $B_r$, si l'on suppose que
$$
Z_\infty^p = \bigcap\nolimits_r Z_r^p =
\frac{\bigcap_r (F^pK \cap d^{-1}(F^{p + r}K) + F^{p + 1}K)}{F^{p + 1}K}
$$
et
$$
B_\infty^p = \bigcup\nolimits_r B_r^p =
\frac{\bigcup_r (F^pK \cap d(F^{p - r + 1}K) + F^{p + 1}K)}{F^{p + 1}K}.
$$
existent, alors $Z_\infty$ et $B_\infty$ existent et ont pour composantes de degré $p$
$Z_\infty^p$ et $B_\infty^p$ (cela résulte d'un argument élémentaire
sur les unions et intersections de sous-objets gradués). Ainsi,
$$
E_\infty^p =
\frac{\bigcap_r (F^pK \cap d^{-1}(F^{p + r}K) + F^{p + 1}K)}
{\bigcup_r (F^pK \cap d(F^{p - r + 1}K) + F^{p + 1}K)}.
$$
où le numérateur et le dénominateur existent. On a
\begin{equation}
\label{equation-on-top}
\Ker(d) \cap F^pK + F^{p + 1}K
\subset
\bigcap\nolimits_r \left(F^pK \cap d^{-1}(F^{p + r}K) + F^{p + 1}K\right)
\end{equation}
et
\begin{equation}
\label{equation-at-bottom}
\bigcup\nolimits_r \left(F^pK \cap d(F^{p - r + 1}K) + F^{p + 1}K\right)
\subset
\Im(d) \cap F^pK + F^{p + 1}K.
\end{equation}
Par conséquent, un sous-quotient de $E_\infty^p$ est
$$
\frac{\Ker(d) \cap F^pK + F^{p + 1}K}{\Im(d) \cap F^pK + F^{p + 1}K} =
\frac{\Ker(d) \cap F^pK}{\Im(d) \cap F^pK + \Ker(d) \cap F^{p + 1}K}
$$
La comparaison avec la formule de $\text{gr}^pH(K)$ donnée dans la discussion
qui suit la
définition \ref{definition-filtration-cohomology-filtered-differential}
permet de conclure.
\end{proof}

\begin{definition}
\label{definition-filtered-differential-ss-converges}
Soit $\mathcal{A}$ une catégorie abélienne.
Soit $(K, F, d)$ un objet différentiel filtré de $\mathcal{A}$.
On dit que la suite spectrale associée à $(K, F, d)$ :
\begin{enumerate}
\item {\it converge faiblement vers $H(K)$} si $\text{gr}H(K) = E_{\infty}$
via le lemme \ref{lemma-compute-filtered-cohomology},
\item {\it aboutit à $H(K)$} si elle converge faiblement vers $H(K)$ et si
l'on a $\bigcap F^pH(K) = 0$ et $\bigcup F^pH(K) = H(K)$.
\end{enumerate}
\end{definition}

\noindent
Malheureusement, il semble difficile de trouver une terminologie cohérente pour ces
notions dans la littérature.

\begin{lemma}
\label{lemma-filtered-differential-ss-converges}
Soit $\mathcal{A}$ une catégorie abélienne.
Soit $(K, F, d)$ un objet différentiel filtré de $\mathcal{A}$.
La suite spectrale associée :
\begin{enumerate}
\item converge faiblement vers $H(K)$ si et seulement si, pour tout
$p \in \mathbf{Z}$, les inclusions des équations
(\ref{equation-at-bottom}) et (\ref{equation-on-top}) sont des égalités,
\item aboutit à $H(K)$ si et seulement si elle converge faiblement vers $H(K)$ et
$\bigcap_p (\Ker(d) \cap F^pK + \Im(d)) = \Im(d)$
et $\bigcup_p (\Ker(d) \cap F^pK + \Im(d)) = \Ker(d)$.
\end{enumerate}
\end{lemma}

\begin{proof}
Cela résulte immédiatement des discussions ci-dessus.
\end{proof}
\section{Suites spectrales : complexes filtrés}
\label{section-filtered-complex}

\begin{definition}
\label{definition-filtered-complex}
Soit $\mathcal{A}$ une catégorie abélienne.
Un {\it complexe filtré $K^\bullet$ de $\mathcal{A}$}
est un complexe de $\text{Fil}(\mathcal{A})$ (voir la
définition \ref{definition-filtered}).
\end{definition}

\noindent
Nous noterons $F$ la filtration sur les objets. Ainsi,
$F^pK^n$ désigne le $p$-ième cran de la filtration du $n$-ième terme du
complexe. Remarquons que chaque $F^pK^\bullet$ est un complexe de $\mathcal{A}$.
On aurait donc aussi pu définir un complexe filtré comme un objet filtré
dans la catégorie (abélienne) des complexes de $\mathcal{A}$.
En particulier, $\text{gr} K^\bullet$ est un objet gradué de la
catégorie des complexes de $\mathcal{A}$.

\medskip\noindent
Pour décrire la suite spectrale associée à un tel objet,
supposons pour l'instant que $\mathcal{A}$ soit une catégorie abélienne
qui possède des sommes directes dénombrables et que celles-ci soient exactes
(ce n'est pas automatique, voir la remarque \ref{remark-direct-sums-not-exact}).
Notons $d$ la différentielle de $K$. En oubliant la graduation,
on peut voir $\bigoplus K^n$ comme un objet différentiel filtré de
$\mathcal{A}$. La section \ref{section-filtered-differential}
fournit donc une suite spectrale $(E_r, d_r)_{r \geq 0}$. Dans cette section,
nous déterminons les termes de cette suite spectrale et nous les munissons
d'une structure supplémentaire provenant de la
graduation de $K$.

\medskip\noindent
Commençons par remarquer que $E_0^p = \text{gr}^p K^\bullet$ est un
complexe, et donc est gradué. Ainsi, $E_0$ est naturellement bigradué.
Il est d'usage d'employer la bigraduation
$$
E_0 = \bigoplus\nolimits_{p, q} E_0^{p, q},
\quad
E_0^{p, q} = \text{gr}^p K^{p + q}
$$
L'idée est de considérer $p + q$ comme le {\it degré total} des
classes de (co)homologie. En outre, $p$ est appelé le {\it degré de filtration},
et $q$ le {\it degré complémentaire}.
La différentielle $d_0$ est compatible avec cette
bigraduation de la manière suivante :
$$
d_0  = \bigoplus d_0^{p, q},
\quad
d_0^{p, q} : E_0^{p, q} \to E_0^{p, q + 1}.
$$
En effet, $d_0^p$ n'est autre que la différentielle du complexe
$\text{gr}^p K^\bullet$ (qui apparaît comme $\text{gr}^pE_0$, simplement décalé
un peu).

\medskip\noindent
Pour aller plus loin, identifions les objets $B_r^p$ et $Z_r^p$ introduits
dans la section \ref{section-filtered-differential} comme des objets gradués,
et déterminons les décompositions correspondantes des différentielles.
Nous le faisons de manière tout à fait directe, mais avertissons de nouveau
le lecteur que notre notation ne coïncide pas avec celles rencontrées
ailleurs. Définissons
$$
Z_r^{p, q} =
\frac{F^pK^{p + q} \cap d^{-1}(F^{p + r}K^{p + q + 1}) + F^{p + 1}K^{p + q}}
{F^{p + 1}K^{p + q}}
$$
et
$$
B_r^{p, q} =
\frac{F^pK^{p + q} \cap d(F^{p - r + 1}K^{p + q - 1}) + F^{p + 1}K^{p + q}}
{F^{p + 1}K^{p + q}}
$$
et, bien entendu, $E_r^{p, q} = Z_r^{p, q}/B_r^{p, q}$.
Avec ces définitions, il est tout à fait clair que
$Z_r^p = \bigoplus_q Z_r^{p, q}$,
$B_r^p = \bigoplus_q B_r^{p, q}$ et
$E_r^p = \bigoplus_q E_r^{p, q}$. De plus, on a
$$
0 \subset \ldots \subset B_r^{p, q} \subset
\ldots
\subset Z_r^{p, q} \subset \ldots \subset E_0^{p, q}
$$
L'application $d_r^p$ se décompose également en somme directe des applications
$$
d_r^{p, q} : E_r^{p, q} \longrightarrow E_r^{p + r, q - r + 1},
\quad
z + F^{p + 1}K^{p + q}
\mapsto
dz + F^{p + r + 1}K^{p + q + 1}
$$
où $z \in F^pK^{p + q} \cap d^{-1}(F^{p + r}K^{p + q + 1})$.

\begin{lemma}
\label{lemma-spectral-sequence-filtered-complex}
Soit $\mathcal{A}$ une catégorie abélienne. Soit $(K^\bullet, F)$ un
complexe filtré de $\mathcal{A}$. Il existe une suite spectrale
$(E_r, d_r)_{r \geq 0}$ dans la catégorie des objets bigradués de $\mathcal{A}$
associée à $(K^\bullet, F)$ telle que $d_r$ soit de bidegré $(r, - r + 1)$
et telle que $E_r$ ait pour composantes bigraduées $E_r^{p, q}$ et pour applications
$d_r^{p, q} : E_r^{p, q} \to E_r^{p + r, q - r + 1}$ celles données ci-dessus.
De plus, on a $E_0^{p, q} = \text{gr}^p(K^{p + q})$,
$d_0^{p, q} = \text{gr}^p(d^{p + q})$
et $E_1^{p, q} = H^{p + q}(\text{gr}^p(K^\bullet))$.
\end{lemma}

\begin{proof}
Si $\mathcal{A}$ possède des sommes directes dénombrables et si celles-ci
sont exactes, cela résulte de la discussion ci-dessus.
Dans le cas général, procédons comme suit ; nous conseillons vivement au lecteur
d'omettre cette preuve. Considérons l'objet bigradué $A = (F^{p + 1}K^{p + 1 + q})$
de $\mathcal{A}$, c'est-à-dire que l'on place $F^{p + 1}K^{p + 1 + q}$ en degré $(p, q)$
(le décalage inhabituel de la numérotation donnera plus tard la bonne numérotation).
Munissons-le d'une différentielle $d : A \to A[0, 1]$ en utilisant $d$
sur chaque composante. Alors $(A, d)$ est un objet différentiel bigradué.
Considérons l'application
$$
\alpha : A \to A[-1, 1]
$$
donnée en degré $(p, q)$ par l'inclusion
$F^{p + 1}K^{p + 1 + q} \to F^pK^{p + 1 + q}$.
C'est un morphisme injectif d'objets différentiels
$\alpha : (A, d) \to (A, d)[-1, 1]$. On peut donc appliquer la
remarque \ref{remark-differential-object-selfmap}
avec $S = [0, 1]$ et $T = [1, -1]$.
La suite spectrale correspondante $(E_r, d_r)_{r \geq 0}$
d'objets bigradués est la suite spectrale cherchée.
Développons quelque peu les définitions. Tout d'abord, on a
$E_r = (E_r^{p, q})$. Ensuite, puisque $T^rS = [r, -r + 1]$,
on a $d_r : E_r \to E_r[r, -r + 1]$, ce qui signifie que
$d_r^{p, q} : E_r^{p, q} \to E_r^{p + r, q - r + 1}$.

\medskip\noindent
Pour vérifier la description des composantes graduées, raisonnons
comme ci-dessus. On a d'abord
$$
E_0 = \Coker(\alpha : A \to A[-1, 1])[0, -1] =
\Coker(\alpha[0, -1] : A[0, -1] \to A[-1, 0])
$$
et notre choix de numérotation donne
$$
E_0^{p, q} = \Coker(F^{p + 1}K^{p + q} \to F^pK^{p + q}) = \text{gr}^pK^{p + q}
$$
La première différentielle est donnée par
$d_0^{p, q} = \text{gr}^pd^{p + q} : E_0^{p, q} \to E_0^{p, q + 1}$.
Ensuite, la description des bords $B_r$ et des cocycles $Z_r$
dans la remarque \ref{remark-differential-object-selfmap}
se traduit directement par les formules
pour $Z_r^{p, q}$ et $B_r^{p, q}$ données ci-dessus.
\end{proof}

\begin{lemma}
\label{lemma-spectral-sequence-filtered-complex-d1}
Soit $\mathcal{A}$ une catégorie abélienne.
Soit $(K^\bullet, F)$ un complexe filtré de $\mathcal{A}$.
Supposons que $\mathcal{A}$ possède des sommes directes dénombrables.
Soit $(E_r, d_r)_{r \geq 0}$ la suite spectrale
associée à $(K^\bullet, F)$.
\begin{enumerate}
\item L'application
$$
d_1^{p, q} :
E_1^{p, q} = H^{p + q}(\text{gr}^p(K^\bullet))
\longrightarrow
E_1^{p + 1, q} = H^{p + q + 1}(\text{gr}^{p + 1}(K^\bullet))
$$
est égale au morphisme de bord en cohomologie associé à la suite
exacte courte de complexes
$$
0 \to \text{gr}^{p + 1}(K^\bullet) \to
F^pK^\bullet/F^{p + 2}K^\bullet \to \text{gr}^p(K^\bullet) \to 0.
$$
\item Supposons que $d(F^pK) \subset F^{p + 1}K$ pour tout $p \in \mathbf{Z}$.
Alors $d$ induit la différentielle nulle sur $\text{gr}^p(K^\bullet)$,
et donc
$E_1^{p, q} = \text{gr}^p(K^\bullet)^{p + q}$.
De plus, dans ce cas,
$$
d_1^{p, q} :
E_1^{p, q} = \text{gr}^p(K^\bullet)^{p + q}
\longrightarrow
E_1^{p + 1, q} = \text{gr}^{p + 1}(K^\bullet)^{p + q + 1}
$$
est le morphisme induit par $d$.
\end{enumerate}
\end{lemma}

\begin{proof}
Cela résulte immédiatement de la formule donnée pour la différentielle
$d_1^{p, q}$ juste avant le lemme \ref{lemma-spectral-sequence-filtered-complex}.
\end{proof}

\begin{lemma}
\label{lemma-spectral-sequence-filtered-complex-functorial}
Soit $\mathcal{A}$ une catégorie abélienne.
Soit $\alpha : (K^\bullet, F) \to (L^\bullet, F)$ un morphisme de
complexes filtrés de $\mathcal{A}$. Soient $(E_r(K), d_r)_{r \geq 0}$,
resp.\ $(E_r(L), d_r)_{r \geq 0}$, les suites spectrales associées
à $(K^\bullet, F)$, resp.\ à $(L^\bullet, F)$.
Le morphisme $\alpha$ induit un morphisme canonique de suites
spectrales $\{\alpha_r : E_r(K) \to E_r(L)\}_{r \geq 0}$ compatible
avec les bigraduations.
\end{lemma}

\begin{proof}
Cela ressort immédiatement de la représentation explicite des termes des
suites spectrales.
\end{proof}

\begin{definition}
\label{definition-filtration-cohomology-filtered-complex}
Soit $\mathcal{A}$ une catégorie abélienne.
Soit $(K^\bullet, F)$ un complexe filtré de $\mathcal{A}$.
La {\it filtration induite} sur $H^n(K^\bullet)$ est la filtration définie
par $F^pH^n(K^\bullet) = \Im(H^n(F^pK^\bullet) \to H^n(K^\bullet))$.
\end{definition}

\noindent
En explicitant cette définition, on voit que
\begin{equation}
\label{equation-filtration-cohomology}
F^pH^n(K^\bullet, d) =
\frac{\Ker(d) \cap F^pK^n + \Im(d) \cap K^n}{\Im(d) \cap K^n}
\end{equation}
et donc que
\begin{equation}
\label{equation-graded-cohomology}
\text{gr}^p H^n(K^\bullet) =
\frac{\Ker(d) \cap F^pK^n}{\Ker(d) \cap F^{p + 1}K^n + \Im(d) \cap F^pK^n}
\end{equation}
(une étape intermédiaire a été omise).

\begin{lemma}
\label{lemma-compute-cohomology-filtered-complex}
Soit $\mathcal{A}$ une catégorie abélienne. Soit $(K^\bullet, F)$ un complexe
filtré de $\mathcal{A}$. Si $Z_\infty^{p, q}$ et $B_\infty^{p, q}$ existent
(voir la preuve), alors :
\begin{enumerate}
\item la limite $E_\infty$ existe et est un objet bigradué, avec
$E_\infty^{p, q} = Z_\infty^{p, q}/B_\infty^{p, q}$ en bidegré $(p, q)$,
\item la $p$-ième composante graduée $\text{gr}^pH^n(K^\bullet)$ du
$n$-ième objet de cohomologie de $K^\bullet$ est un sous-quotient de
$E_\infty^{p, n - p}$.
\end{enumerate}
\end{lemma}

\begin{proof}
Les objets $Z_\infty$, $B_\infty$ et la limite
$E_\infty = Z_\infty/B_\infty$ de la
définition \ref{definition-limit-spectral-sequence}
sont des objets bigradués de $\mathcal{A}$ par notre construction de la
suite spectrale dans le lemme \ref{lemma-spectral-sequence-filtered-complex}.
Puisque $Z_r^{p, q}$, resp.\ $B_r^{p, q}$, est la composante de bidegré $(p, q)$
de $Z_r$, resp.\ de $B_r$, si l'on suppose que
$$
Z_\infty^{p, q} = \bigcap\nolimits_r Z_r^{p, q} =
\bigcap\nolimits_r
\frac{F^pK^{p + q} \cap d^{-1}(F^{p + r}K^{p + q + 1}) + F^{p + 1}K^{p + q}}
{F^{p + 1}K^{p + q}}
$$
et
$$
B_\infty^{p, q} = \bigcup\nolimits_r B_r^{p, q} =
\bigcup\nolimits_r
\frac{F^pK^{p + q} \cap d(F^{p - r + 1}K^{p + q - 1}) + F^{p + 1}K^{p + q}}
{F^{p + 1}K^{p + q}}
$$
existent, alors $Z_\infty$ et $B_\infty$ existent, avec pour composantes de bidegré $(p, q)$
$Z_\infty^{p, q}$ et $B_\infty^{p, q}$ (cela résulte d'un argument élémentaire
sur les unions et intersections d'objets bigradués). Ainsi,
$$
E_\infty^{p, q} =
\frac{\bigcap_r (F^pK^{p + q} \cap d^{-1}(F^{p + r}K^{p + q + 1})
+ F^{p + 1}K^{p + q})}
{\bigcup_r (F^pK^{p + q} \cap d(F^{p - r + 1}K^{p + q - 1})
+ F^{p + 1}K^{p + q})}.
$$
où le numérateur et le dénominateur existent. Avec $n = p + q$, on a
\begin{equation}
\label{equation-on-top-bigraded}
\Ker(d) \cap F^pK^{n} + F^{p + 1}K^{n}
\subset
\bigcap\nolimits_r
\left(
F^pK^{n} \cap d^{-1}(F^{p + r}K^{n + 1}) + F^{p + 1}K^{n}
\right)
\end{equation}
et
\begin{equation}
\label{equation-at-bottom-bigraded}
\bigcup\nolimits_r
\left(
F^pK^{n} \cap d(F^{p - r + 1}K^{n - 1}) + F^{p + 1}K^{n}
\right)
\subset
\Im(d) \cap F^pK^{n} + F^{p + 1}K^{n}.
\end{equation}
Ainsi, un sous-quotient de $E_\infty^{p, q}$ est
$$
\frac{\Ker(d) \cap F^pK^{n} + F^{p + 1}K^n}
{\Im(d) \cap F^pK^{n} + F^{p + 1}K^{n}} =
\frac{\Ker(d) \cap F^pK^n}{\Im(d) \cap F^pK^n + \Ker(d) \cap F^{p + 1}K^n}
$$
La comparaison avec (\ref{equation-graded-cohomology}) permet de conclure.
\end{proof}

\begin{definition}
\label{definition-bounded-ss}
Soit $\mathcal{A}$ une catégorie abélienne. Soit $(E_r, d_r)_{r \geq r_0}$
une suite spectrale d'objets bigradués de
$\mathcal{A}$, avec $d_r$ de bidegré $(r, -r + 1)$.
On dit qu'une telle suite spectrale est :
\begin{enumerate}
\item {\it régulière} si, pour tous $p, q \in \mathbf{Z}$, il existe
un $b = b(p, q)$ tel que les applications
$d_r^{p, q} : E_r^{p, q} \to E_r^{p + r, q - r + 1}$ soient nulles pour $r \geq b$,
\item {\it corégulière} si, pour tous $p, q \in \mathbf{Z}$, il existe un
$b = b(p, q)$ tel que les applications
$d_r^{p - r, q + r - 1} : E_r^{p - r, q + r - 1} \to E_r^{p, q}$
soient nulles pour $r \geq b$,
\item {\it bornée} si, pour tout $n$,
il n'existe qu'un nombre fini de $E_{r_0}^{p, n - p}$ non nuls,
\item {\it bornée inférieurement} si, pour tout $n$, il existe un $b = b(n)$ tel que
$E_{r_0}^{p, n - p} = 0$ pour $p \geq b$.
\item {\it bornée supérieurement} si, pour tout $n$, il existe un $b = b(n)$ tel que
$E_{r_0}^{p, n - p} = 0$ pour $p \leq b$.
\end{enumerate}
\end{definition}

\noindent
Être bornée inférieurement signifie que, si l'on considère $E_r^{p, q}$ sur la droite
$p + q = n$ (de pente $-1$), on obtient des zéros lorsque $(p, q)$
se déplace vers le bas et vers la droite.
Comme nous l'avons déjà indiqué, la littérature n'emploie pas de terminologie
cohérente pour ces notions.

\begin{lemma}
\label{lemma-relate-boundedness}
Dans la situation de la définition \ref{definition-bounded-ss},
soient $Z_r^{p, q}, B_r^{p, q} \subset E_{r_0}^{p, q}$ les
composantes graduées de degré $(p, q)$ de $Z_r, B_r$, définies comme dans la
section \ref{section-spectral-sequence}.
\begin{enumerate}
\item La suite spectrale est régulière si et seulement si, pour tous $p, q$,
il existe un $r = r(p, q)$ tel que
$Z_r^{p, q} = Z_{r + 1}^{p, q} = \ldots$
\item La suite spectrale est corégulière si et seulement si, pour tous $p, q$,
il existe un $r = r(p, q)$ tel que
$B_r^{p, q} = B_{r + 1}^{p, q} = \ldots$
\item La suite spectrale est bornée si et seulement si elle est à la fois
bornée inférieurement et bornée supérieurement.
\item Si la suite spectrale est bornée inférieurement, alors elle est régulière.
\item Si la suite spectrale est bornée supérieurement, alors elle est corégulière.
\end{enumerate}
\end{lemma}

\begin{proof}
Omis. Indication : si $E_r^{p, q} = 0$, alors on a $E_{r'}^{p, q} = 0$
pour tout $r' \geq r$.
\end{proof}

\begin{definition}
\label{definition-filtered-complex-ss-converges}
Soit $\mathcal{A}$ une catégorie abélienne. Soit $(K^\bullet, F)$ un
complexe filtré de $\mathcal{A}$. On dit que la suite spectrale
associée à $(K^\bullet, F)$ :
\begin{enumerate}
\item {\it converge faiblement vers $H^*(K^\bullet)$} si
$\text{gr}^pH^n(K^\bullet) = E_{\infty}^{p, n - p}$
via le lemme \ref{lemma-compute-cohomology-filtered-complex},
pour tous $p, n \in \mathbf{Z}$,
\item {\it aboutit à $H^*(K^\bullet)$} si elle converge faiblement vers
$H^*(K^\bullet)$ et si $\bigcap_p F^pH^n(K^\bullet) = 0$ et
$\bigcup_p F^p H^n(K^\bullet) = H^n(K^\bullet)$ pour tout $n$,
\item {\it converge vers $H^*(K^\bullet)$} si elle est régulière,
aboutit à $H^*(K^\bullet)$ et si
$H^n(K^\bullet) = \lim_p H^n(K^\bullet)/F^pH^n(K^\bullet)$.
\end{enumerate}
\end{definition}

\noindent
La convergence faible, l'aboutissement ou la convergence se notent
$E_r^{p, q} \Rightarrow H^{p + q}(K^\bullet)$.
Comme nous l'avons déjà indiqué, la littérature n'emploie pas de terminologie
cohérente pour ces notions.

\begin{lemma}
\label{lemma-filtered-complex-ss-converges}
Soit $\mathcal{A}$ une catégorie abélienne. Soit $(K^\bullet, F)$ un complexe
filtré de $\mathcal{A}$. La suite spectrale associée :
\begin{enumerate}
\item converge faiblement vers $H^*(K^\bullet)$ si et seulement si, pour tous
$p, q \in \mathbf{Z}$, les inclusions des équations
(\ref{equation-at-bottom-bigraded}) et (\ref{equation-on-top-bigraded}) sont des égalités,
\item aboutit à $H^*(K)$ si et seulement si elle converge faiblement vers $H^*(K^\bullet)$
et si l'on a
$\bigcap_p (\Ker(d) \cap F^pK^n + \Im(d) \cap K^n) = \Im(d) \cap K^n$
et
$\bigcup_p (\Ker(d) \cap F^pK^n + \Im(d) \cap K^n) = \Ker(d) \cap K^n$.
\end{enumerate}
\end{lemma}

\begin{proof}
Cela résulte immédiatement des discussions ci-dessus.
\end{proof}

\begin{lemma}
\label{lemma-biregular-ss-converges}
Soit $\mathcal{A}$ une catégorie abélienne. Soit $(K^\bullet, F)$ un
complexe filtré de $\mathcal{A}$. Supposons que la filtration de chaque $K^n$
soit finie (voir la définition \ref{definition-filtered}). Alors :
\begin{enumerate}
\item la suite spectrale associée à $(K^\bullet, F)$ est bornée,
\item la filtration de chaque $H^n(K^\bullet)$ est finie,
\item la suite spectrale associée à $(K^\bullet, F)$ converge
vers $H^*(K^\bullet)$,
\item si $\mathcal{C} \subset \mathcal{A}$ est une sous-catégorie de Serre faible
et si, pour un certain $r$, on a $E_r^{p, q} \in \mathcal{C}$ pour tous
$p, q \in \mathbf{Z}$, alors $H^n(K^\bullet)$ appartient à $\mathcal{C}$.
\end{enumerate}
\end{lemma}

\begin{proof}
La partie (1) découle de $E_0^{p, n - p} = \text{gr}^p K^n$.
La partie (2) résulte immédiatement de l'équation (\ref{equation-filtration-cohomology}).
Nous allons utiliser le lemme \ref{lemma-filtered-complex-ss-converges} pour montrer
que la suite spectrale converge faiblement. Fixons $p, n \in \mathbf{Z}$.
Le membre de droite de (\ref{equation-on-top-bigraded})
est égal à $F^pK^n \cap \Ker(d) + F^{p + 1}K^n$, car
$F^{p + r}K^{n + 1} = 0$ pour $r \gg 0$. Ainsi, (\ref{equation-on-top-bigraded})
est une égalité. Le membre de gauche de (\ref{equation-at-bottom-bigraded})
est égal à $F^pK^n \cap \Im(d) + F^{p + 1}K^n$, car
$F^{p - r + 1}K^{n - 1} = K^{n - 1}$ pour $r \gg 0$.
Ainsi, (\ref{equation-at-bottom-bigraded}) est une égalité. Puisque la filtration
de $H^n(K^\bullet)$ est finie par (2), on voit que la suite aboutit.
Pour prouver la convergence, il faut montrer que la suite spectrale est
régulière, ce qui découle de ce qu'elle est bornée
(lemme \ref{lemma-relate-boundedness}), et il faut
montrer que $H^n(K^\bullet) = \lim_p H^n(K^\bullet)/F^pH^n(K^\bullet)$,
ce qui découle du fait, prouvé dans la partie (2), que la filtration de $H^*(K^\bullet)$
est finie.

\medskip\noindent
Prouvons (4). Supposons que, pour un certain $r \geq 0$, on ait
$E_r^{p, q} \in \mathcal{C}$ pour tous $p, q \in \mathbf{Z}$, où
$\mathcal{C}$ est une sous-catégorie de Serre faible de $\mathcal{A}$.
Alors $E_{r + 1}^{p, q}$
appartient aussi à $\mathcal{C}$, voir le
lemme \ref{lemma-characterize-weak-serre-subcategory}.
Par la bornitude établie ci-dessus (qui implique que la suite spectrale
est à la fois régulière et corégulière, voir le lemme \ref{lemma-relate-boundedness}),
on peut trouver un $r' \geq r$ tel que $E_\infty^{p, q} = E_{r'}^{p, q}$
pour tous les $p, q$ tels que $p + q = n$. Ainsi, $H^n(K^\bullet)$ est un objet
de $\mathcal{A}$ muni d'une filtration finie dont les composantes graduées
appartiennent à $\mathcal{C}$. Cela implique que $H^n(K^\bullet)$ appartient à $\mathcal{C}$
d'après le lemme \ref{lemma-characterize-weak-serre-subcategory}.
\end{proof}

\begin{lemma}
\label{lemma-biregular-ss-relation-in-K0}
Soit $\mathcal{A}$ une catégorie abélienne. Soit $(K^\bullet, F)$ un
complexe filtré de $\mathcal{A}$. Supposons que la filtration de chaque $K^n$
soit finie (voir la définition \ref{definition-filtered}) et que, pour un certain
$r$, seuls un nombre fini de $E_r^{p, q}$ soient non nuls. Alors
seuls un nombre fini de $H^n(K^\bullet)$ sont non nuls et l'on a
$$
\sum (-1)^n[H^n(K^\bullet)] = \sum (-1)^{p + q} [E_r^{p, q}]
$$
dans $K_0(\mathcal{A}')$, où $\mathcal{A}'$ est la plus petite sous-catégorie de
Serre faible de $\mathcal{A}$ contenant les objets
$E_r^{p, q}$.
\end{lemma}

\begin{proof}
Notons $E_r^{even}$ et $E_r^{odd}$ les parties paire et impaire de $E_r$,
définies comme les sommes directes des composantes $(p, q)$ où $p + q$ est pair
et impair. La différentielle $d_r$ définit des applications
$\varphi : E_r^{even} \to E_r^{odd}$ et $\psi : E_r^{odd} \to E_r^{even}$
dont les compositions dans les deux sens sont nulles.
On voit alors que
\begin{align*}
[E_r^{even}] - [E_r^{odd}] & =
[\Ker(\varphi)] + [\Im(\varphi)] - [\Ker(\psi)] - [\Im(\psi)] \\
& =
[\Ker(\varphi)/\Im(\psi)] - [\Ker(\psi)/\Im(\varphi)] \\
& =
[E_{r + 1}^{even}] - [E_{r + 1}^{odd}]
\end{align*}
Remarquons que tous les objets intermédiaires appartiennent à la plus petite sous-catégorie de Serre faible
contenant les objets $E_r^{p, q}$.
En poursuivant ainsi, on voit que l'on peut augmenter $r$ à volonté.
Puisqu'il n'existe qu'un nombre fini de paires $(p, q)$ pour lesquelles
$E_r^{p, q}$ est non nul, propriété héritée par
$E_{r + 1}, E_{r + 2}, \ldots$, on voit que l'on peut supposer $d_r = 0$.
À ce stade, $H^n(K^\bullet)$ possède une filtration finie
(lemme \ref{lemma-biregular-ss-converges}) dont les composantes graduées
sont exactement les $E_r^{p, n - p}$, et le résultat est clair.
\end{proof}

\noindent
Le lemme suivant constitue surtout une vérification de cohérence de nos définitions.
Il est certain que, si l'on a un complexe filtré tel que, pour tout $n$,
$$
H^n(F^pK^\bullet) = 0\text{ pour }p \gg 0
\quad\text{et}\quad
H^n(F^pK^\bullet) = H^n(K^\bullet)\text{ pour }p \ll 0,
$$
alors la suite spectrale correspondante devrait converger.

\begin{lemma}
\label{lemma-ss-converges-trivial}
Soit $\mathcal{A}$ une catégorie abélienne. Soit $(K^\bullet, F)$ un
complexe filtré de $\mathcal{A}$. Supposons que :
\begin{enumerate}
\item pour tout $n$, il existe $p_0(n)$ tel que
$H^n(F^pK^\bullet) = 0$ pour $p \geq p_0(n)$,
\item pour tout $n$, il existe $p_1(n)$ tel que
$H^n(F^pK^\bullet) \to H^n(K^\bullet)$ soit un isomorphisme
pour $p \leq p_1(n)$.
\end{enumerate}
Alors :
\begin{enumerate}
\item la suite spectrale associée à $(K^\bullet, F)$ est bornée,
\item la filtration de chaque $H^n(K^\bullet)$ est finie,
\item la suite spectrale associée à $(K^\bullet, F)$ converge
vers $H^*(K^\bullet)$.
\end{enumerate}
\end{lemma}

\begin{proof}
Fixons $n$. La suite exacte longue de cohomologie associée à
la suite exacte courte de complexes
$$
0 \to F^{p + 1}K^\bullet \to F^pK^\bullet \to \text{gr}^pK^\bullet \to 0
$$
montre que $E_1^{p, n - p} = 0$ pour $p \geq \max(p_0(n), p_0(n + 1))$ et
$p < \min(p_1(n), p_1(n + 1))$. La suite spectrale est donc bornée
(définition \ref{definition-bounded-ss}). Cela prouve (1).

\medskip\noindent
Il résulte clairement des hypothèses et de la
définition \ref{definition-filtration-cohomology-filtered-complex}
que la filtration de $H^n(K^\bullet)$ est finie. Cela prouve (2).

\medskip\noindent
Montrons ensuite que la suite spectrale converge faiblement vers
$H^*(K^\bullet)$ en utilisant le
lemme \ref{lemma-filtered-complex-ss-converges}.
Montrons que l'on a égalité dans (\ref{equation-on-top-bigraded}).
En effet, pour $p + r > p_0(n + 1)$, l'application
$$
d : F^pK^{n} \cap d^{-1}(F^{p + r}K^{n + 1}) \to F^{p + r}K^{n + 1}
$$
prend ses valeurs dans l'image de $d : F^{p + r}K^n \to F^{p + r}K^{n + 1}$,
car le complexe $F^{p + r}K^\bullet$ est exact en degré $n + 1$.
On en conclut que $F^pK^{n} \cap d^{-1}(F^{p + r}K^{n + 1}) =
d(F^{p + r}K^n) + \Ker(d) \cap F^pK^n$. Pour un tel $r$, on a donc
$$
\Ker(d) \cap F^pK^{n} + F^{p + 1}K^{n} =
F^pK^{n} \cap d^{-1}(F^{p + r}K^{n + 1}) + F^{p + 1}K^{n}
$$
ce qui prouve l'égalité voulue. Pour montrer que l'on a égalité dans
(\ref{equation-at-bottom-bigraded}), on utilise que, pour $p - r + 1 < p_1(n)$,
on a
$$
d(F^{p - r + 1}K^{n - 1}) = \Im(d) \cap F^{p - r + 1}K^n
$$
car l'application $F^{p - r + 1}K^\bullet \to K^\bullet$ induit un
isomorphisme en cohomologie de degré $n$. Cela montre que
l'on a
$$
F^pK^{n} \cap d(F^{p - r + 1}K^{n - 1}) + F^{p + 1}K^{n} =
\Im(d) \cap F^pK^{n} + F^{p + 1}K^{n}
$$
pour un tel $r$, ce qui prouve l'égalité voulue.

\medskip\noindent
Pour voir que la suite spectrale aboutit à $H^*(K^\bullet)$ au moyen du
lemme \ref{lemma-filtered-complex-ss-converges}, il faut montrer que
$\bigcap_p (\Ker(d) \cap F^pK^n + \Im(d) \cap K^n) = \Im(d) \cap K^n$
et
$\bigcup_p (\Ker(d) \cap F^pK^n + \Im(d) \cap K^n) = \Ker(d) \cap K^n$.
Pour $p \geq p_0(n)$, on a
$\Ker(d) \cap F^pK^n + \Im(d) \cap K^n = \Im(d) \cap K^n$,
et pour $p \leq p_1(n)$, on a
$\Ker(d) \cap F^pK^n + \Im(d) \cap K^n = \Ker(d) \cap K^n$.
En combinant convergence faible, aboutissement et bornitude, on voit
que (2) et (3) sont vraies.
\end{proof}

\section{Suites spectrales : bicomplexes}
\label{section-double-complex}

\noindent
Soit $K^{\bullet, \bullet}$ un bicomplexe, voir
section \ref{section-double-complexes}. On note traditionnellement
$H^p_I(K^{\bullet, \bullet})$
le complexe dont les termes sont $\Ker(d_1^{p, q})/\Im(d_1^{p - 1, q})$
($q$ variant) et dont la différentielle est induite par $d_2$.
Alors $H^q_{II}(H^p_I(K^{\bullet, \bullet}))$ désigne sa cohomologie en
degré $q$. On note aussi traditionnellement $H^q_{II}(K^{\bullet, \bullet})$
le complexe dont les termes sont $\Ker(d_2^{p, q})/\Im(d_2^{p, q - 1})$
($p$ variant) et dont la différentielle est induite par $d_1$.
Alors $H^p_I(H^q_{II}(K^{\bullet, \bullet}))$ désigne sa cohomologie en
degré $p$. Il s'avérera que ces groupes de cohomologie apparaissent
comme les termes de la suite spectrale d'une filtration du
complexe simple associé, voir
Définition \ref{definition-associated-simple-complex}.

\medskip\noindent
Il y a deux filtrations naturelles sur le complexe simple
$\text{Tot}(K^{\bullet, \bullet})$
associé au bicomplexe $K^{\bullet, \bullet}$. Plus précisément, nous
posons
$$
F_I^p(\text{Tot}^n(K^{\bullet, \bullet})) =
\bigoplus\nolimits_{i + j = n, \ i \geq p} K^{i, j}
\quad
\text{et}
\quad
F_{II}^p(\text{Tot}^n(K^{\bullet, \bullet})) =
\bigoplus\nolimits_{i + j = n, \ j \geq p} K^{i, j}.
$$
On vérifie immédiatement que $(\text{Tot}(K^{\bullet, \bullet}), F_I)$ et
$(\text{Tot}(K^{\bullet, \bullet}), F_{II})$ sont des complexes filtrés.
Par la section \ref{section-filtered-complex},
nous obtenons deux suites spectrales. On note traditionnellement
$({}'E_r, {}'d_r)_{r \geq 0}$ la suite spectrale associée
à la filtration $F_I$, et $({}''E_r, {}''d_r)_{r \geq 0}$
la suite spectrale associée à la filtration $F_{II}$.
Voici une description de ces suites spectrales.

\begin{lemma}
\label{lemma-ss-double-complex}
Soit $\mathcal{A}$ une catégorie abélienne.
Soit $K^{\bullet, \bullet}$ un bicomplexe.
Les suites spectrales associées à $K^{\bullet, \bullet}$
ont les termes suivants :
\begin{enumerate}
\item ${}'E_0^{p, q} = K^{p, q}$ avec
${}'d_0^{p, q} = (-1)^p d_2^{p, q} : K^{p, q} \to K^{p, q + 1}$,
\item ${}''E_0^{p, q} = K^{q, p}$ avec
${}''d_0^{p, q} = d_1^{q, p} : K^{q, p} \to K^{q + 1, p}$,
\item ${}'E_1^{p, q} = H^q(K^{p, \bullet})$ avec
${}'d_1^{p, q} = H^q(d_1^{p, \bullet})$,
\item ${}''E_1^{p, q} = H^q(K^{\bullet, p})$ avec
${}''d_1^{p, q} = (-1)^q H^q(d_2^{\bullet, p})$,
\item ${}'E_2^{p, q} = H^p_I(H^q_{II}(K^{\bullet, \bullet}))$,
\item ${}''E_2^{p, q} = H^p_{II}(H^q_I(K^{\bullet, \bullet}))$.
\end{enumerate}
\end{lemma}

\begin{proof}
Omis.
\end{proof}

\noindent
Ces suites spectrales définissent deux filtrations sur
$H^n(\text{Tot}(K^{\bullet, \bullet}))$.
Nous les noterons $F_I$ et $F_{II}$.

\begin{definition}
\label{definition-ss-double-complex-converge}
Soit $\mathcal{A}$ une catégorie abélienne.
Soit $K^{\bullet, \bullet}$ un bicomplexe.
Nous disons que la suite spectrale $({}'E_r, {}'d_r)_{r \geq 0}$
{\it converge faiblement vers $H^n(\text{Tot}(K^{\bullet, \bullet}))$},
{\it aboutit à $H^n(\text{Tot}(K^{\bullet, \bullet}))$}, ou
{\it converge vers $H^n(\text{Tot}(K^{\bullet, \bullet}))$}
si la Définition \ref{definition-filtered-complex-ss-converges} s'applique.
De même, nous disons que la suite spectrale $({}''E_r, {}''d_r)_{r \geq 0}$
{\it converge faiblement vers $H^n(\text{Tot}(K^{\bullet, \bullet}))$},
{\it aboutit à $H^n(\text{Tot}(K^{\bullet, \bullet}))$}, ou
{\it converge vers $H^n(\text{Tot}(K^{\bullet, \bullet}))$}
si la Définition \ref{definition-filtered-complex-ss-converges} s'applique.
\end{definition}

\noindent
Comme indiqué plus haut, la littérature n'emploie aucune terminologie
uniforme pour ces notions. Dans la situation de la définition, la première
suite spectrale converge faiblement si, pour tout $n$,
$$
\text{gr}_{F_I}(H^n(\text{Tot}(K^{\bullet, \bullet}))) =
\oplus_{p + q = n} {}'E_\infty^{p, q}
$$
via la comparaison canonique du
Lemme \ref{lemma-compute-cohomology-filtered-complex}.
De même, la seconde suite spectrale $({}''E_r, {}''d_r)_{r \geq 0}$
converge faiblement si, pour tout $n$,
$$
\text{gr}_{F_{II}}(H^n(\text{Tot}(K^{\bullet, \bullet}))) =
\oplus_{p + q = n} {}''E_\infty^{p, q}
$$
via la comparaison canonique du
Lemme \ref{lemma-compute-cohomology-filtered-complex}.

\begin{lemma}
\label{lemma-first-quadrant-ss}
Soit $\mathcal{A}$ une catégorie abélienne. Soit $K^{\bullet, \bullet}$
un bicomplexe. Supposons que, pour tout $n \in \mathbf{Z}$, il n'y ait
qu'un nombre fini de $K^{p, q}$ non nuls tels que $p + q = n$. Alors
\begin{enumerate}
\item les deux suites spectrales associées à $K^{\bullet, \bullet}$
sont bornées,
\item les filtrations $F_I$, $F_{II}$ sur chaque
$H^n(\text{Tot}(K^{\bullet, \bullet}))$ sont finies,
\item les suites spectrales $({}'E_r, {}'d_r)_{r \geq 0}$ et
$({}''E_r, {}''d_r)_{r \geq 0}$ convergent vers
$H^*(\text{Tot}(K^{\bullet, \bullet}))$,
\item si $\mathcal{C} \subset \mathcal{A}$ est une sous-catégorie de Serre faible
et si, pour un certain $r$, on a ${}'E_r^{p, q} \in \mathcal{C}$ pour tous
$p, q \in \mathbf{Z}$, alors $H^n(\text{Tot}(K^{\bullet, \bullet}))$
appartient à $\mathcal{C}$. Il en va de même pour $({}''E_r, {}''d_r)_{r \geq 0}$.
\end{enumerate}
\end{lemma}

\begin{proof}
Cela résulte immédiatement du Lemme \ref{lemma-biregular-ss-converges}.
\end{proof}

\noindent
Voici notre première application des suites spectrales.

\begin{lemma}
\label{lemma-double-complex-gives-resolution}
Soit $\mathcal{A}$ une catégorie abélienne.
Soit $K^\bullet$ un complexe.
Soit $A^{\bullet, \bullet}$ un bicomplexe.
Soient $\alpha^p : K^p \to A^{p, 0}$ des morphismes.
Supposons que
\begin{enumerate}
\item pour tout $n \in \mathbf{Z}$, il n'y ait qu'un nombre fini de
$A^{p, q}$ non nuls tels que $p + q = n$ ;
\item on ait $A^{p, q} = 0$ si $q < 0$ ;
\item les morphismes $\alpha^p$ donnent lieu à un morphisme
de complexes $\alpha : K^\bullet \to A^{\bullet, 0}$ ;
\item le complexe $A^{p, \bullet}$ soit exact en tout degré
$q \not = 0$ et le morphisme $K^p \to A^{p, 0}$ induise
un isomorphisme $K^p \to \Ker(d_2^{p, 0})$.
\end{enumerate}
Alors $\alpha$ induit un quasi-isomorphisme
$$
K^\bullet \longrightarrow \text{Tot}(A^{\bullet, \bullet})
$$
de complexes.
En outre, il existe une variante de ce lemme faisant intervenir la seconde
variable $q$ à la place de $p$.
\end{lemma}

\begin{proof}
L'application est simplement celle donnée par les morphismes
$K^n \to A^{n, 0} \to \text{Tot}^n(A^{\bullet, \bullet})$,
dont on voit facilement qu'ils définissent un morphisme de complexes.
Considérons la suite spectrale $({}'E_r, {}'d_r)_{r \geq 0}$
associée au bicomplexe $A^{\bullet, \bullet}$.
Par le Lemme \ref{lemma-first-quadrant-ss}, cette suite spectrale converge
et la filtration induite sur $H^n(\text{Tot}(A^{\bullet, \bullet}))$
est finie pour tout $n$.
Par le Lemme \ref{lemma-ss-double-complex} et l'hypothèse (4), on a
${}'E_1^{p, q} = 0$ sauf si $q = 0$, et $'E_1^{p, 0} = K^p$,
la différentielle ${}'d_1^{p, 0}$ s'identifiant à $d_K^p$.
Ainsi ${}'E_2^{p, 0} = H^p(K^\bullet)$ et les autres termes sont nuls.
Cela implique clairement, pour des raisons de degré, que
${}'d_2^{p, q} = {}'d_3^{p, q} = \ldots = 0$.
Nous concluons donc que
$H^n(\text{Tot}(A^{\bullet, \bullet})) = H^n(K^\bullet)$.
Nous omettons de vérifier que cette identification est donnée par le
morphisme de complexes $K^\bullet \to \text{Tot}(A^{\bullet, \bullet})$
introduit ci-dessus.
\end{proof}

\begin{lemma}
\label{lemma-homotopy-complex-complexes}
Soit $\mathcal{A}$ une catégorie abélienne.
Soit $M^\bullet$ un complexe de $\mathcal{A}$. Soit
$$
a :
M^\bullet[0]
\longrightarrow
\left(A^{0, \bullet} \to A^{1, \bullet} \to A^{2, \bullet} \to \ldots \right)
$$
une équivalence d'homotopie dans la catégorie des complexes de complexes
de $\mathcal{A}$. Alors l'application
$\alpha : M^\bullet \to \text{Tot}(A^{\bullet, \bullet})$
induite par $M^\bullet \to A^{0, \bullet}$ est une équivalence d'homotopie.
\end{lemma}

\begin{proof}
L'énoncé a un sens, puisqu'un complexe de complexes est la même chose
qu'un bicomplexe. L'hypothèse signifie qu'il existe une application
$$
b :
\left(A^{0, \bullet} \to A^{1, \bullet} \to A^{2, \bullet} \to \ldots \right)
\longrightarrow
M^\bullet[0]
$$
telle que $a \circ b$ et $b \circ a$ soient homotopes à l'identité
dans la catégorie des complexes de complexes. Cela signifie que $b \circ a$
est l'identité de $M^\bullet[0]$ (parce qu'il n'y a qu'un seul terme en
degré $0$). Observons aussi que $b$ est donnée par une application
$b^0 : A^{0, \bullet} \to M^\bullet$ et par zéro en tous les autres degrés.
Ainsi $b$ induit une application
$\beta : \text{Tot}(A^{\bullet, \bullet}) \to M^\bullet$
et $\beta \circ \alpha$ est l'identité sur $M^\bullet$.
Enfin, il reste à montrer que l'application
$\alpha \circ \beta$ est homotope à l'identité.
À cette fin, choisissons des applications de complexes
$h^n : A^{n, \bullet} \to A^{n - 1, \bullet}$
telles que $a \circ b - \text{id} = d_1 \circ h + h \circ d_1$ ;
elles existent par hypothèse. Ici, $d_1 : A^{n, \bullet} \to A^{n + 1, \bullet}$
sont les différentielles du complexe de complexes. Nous noterons aussi
$d_2$ les différentielles des complexes $A^{n, \bullet}$
pour tout $n$. Soient $h^{n, m} : A^{n, m} \to A^{n - 1, m}$ les composantes de
$h^n$. Nous pouvons alors considérer
$$
h' : \text{Tot}(A^{\bullet, \bullet})^k =
\bigoplus\nolimits_{n + m = k} A^{n, m}
\to
\bigoplus\nolimits_{n + m = k - 1} A^{n, m} =
\text{Tot}(A^{\bullet, \bullet})^{k - 1}
$$
donnée par $h^{n, m}$ sur le facteur direct $A^{n, m}$. Nous calculons alors
que l'application
$$
d_{\text{Tot}(A^{\bullet, \bullet})} \circ h' +
h' \circ d_{\text{Tot}(A^{\bullet, \bullet})}
$$
restreinte au facteur direct $A^{n, m}$ est égale à
$$
d_1^{n - 1, m} \circ h^{n, m} +
(-1)^{n - 1} d_2^{n - 1, m} \circ h^{n, m} +
h^{n + 1, m} \circ d_1^{n, m} + h^{n, m + 1} \circ (-1)^nd_2^{n, m}
$$
Puisque $h^n$ est une application de complexes, les termes
$(-1)^{n - 1} d_2^{n - 1, m} \circ h^{n, m}$ et
$h^{n, m + 1} \circ (-1)^nd_2^{n, m}$ s'annulent mutuellement.
Les deux autres termes donnent
$(\alpha \circ \beta)|_{A^{n, m}} - \text{id}_{A^{n, m}}$
car $a \circ b - \text{id} = d_1 \circ h + h \circ d_1$.
Ceci achève la démonstration.
\end{proof}

\section{Bicomplexes de groupes abéliens}
\label{section-double-complexes-abelian-groups}

\noindent
Dans cette section, nous donnons quelques résultats sur les bicomplexes
de groupes abéliens dont nous ne connaissons pas (encore) les analogues
pour les catégories abéliennes générales. Il faut prendre garde à n'utiliser
ces lemmes que lorsque la catégorie abélienne sous-jacente est celle des
groupes abéliens ou une catégorie analogue (par exemple, la catégorie des
modules sur un anneau). Certains des arguments seront difficiles à suivre
sans dessiner des « zigzags » sur un coin de table ; comparer avec la
démonstration de
Algèbre, Lemme \ref{algebra-lemma-no-spectral-sequence}.

\begin{lemma}
\label{lemma-right-resolution-gives-qis}
Soit $M^\bullet$ un complexe de groupes abéliens. Soit
$$
0 \to M^\bullet \to A_0^\bullet \to A_1^\bullet \to A_2^\bullet \to \ldots
$$
un complexe exact de complexes de groupes abéliens. Posons
$A^{p, q} = A_p^q$ pour obtenir un bicomplexe.
Alors l'application $M^\bullet \to \text{Tot}(A^{\bullet, \bullet})$
induite par $M^\bullet \to A_0^\bullet$ est un quasi-isomorphisme.
\end{lemma}

\begin{proof}
S'il existe $t \in \mathbf{Z}$ tel que $A_0^q = 0$ pour $q < t$, cela résulte
immédiatement du
Lemme \ref{lemma-double-complex-gives-resolution}
(en échangeant $p$ et $q$ comme dans la dernière assertion de ce lemme).
Bien, mais pour tout $t \in \mathbf{Z}$, nous avons un complexe
$$
0 \to
\sigma_{\geq t}M^\bullet \to
\sigma_{\geq t}A_0^\bullet \to
\sigma_{\geq t}A_1^\bullet \to
\sigma_{\geq t}A_2^\bullet \to \ldots
$$
de troncations bêtes. Notons $A(t)^{\bullet, \bullet}$ le bicomplexe
correspondant. Tout élément $\xi$ de $H^n(\text{Tot}(A^{\bullet, \bullet}))$
est l'image d'un élément de $H^n(\text{Tot}(A(t)^{\bullet, \bullet}))$
pour un certain $t$ (considérer des représentants explicites des classes
de cohomologie). Ainsi $\xi$ appartient à l'image de
$H^n(\sigma_{\geq t}M^\bullet)$.
L'application $H^n(M^\bullet) \to H^n(\text{Tot}(A^{\bullet, \bullet}))$
est donc surjective. Elle est injective parce que, pour tout $t$, l'application
$H^n(\sigma_{\geq t}M^\bullet) \to H^n(\text{Tot}(A(t)^{\bullet, \bullet}))$
est injective, et par des arguments analogues.
\end{proof}

\begin{lemma}
\label{lemma-good-resolution-gives-qis}
Soit $M^\bullet$ un complexe de groupes abéliens. Soit
$$
\ldots \to A_2^\bullet \to A_1^\bullet \to A_0^\bullet \to M^\bullet \to 0
$$
un complexe exact de complexes de groupes abéliens tel que, pour tout
$p \in \mathbf{Z}$, les complexes
$$
\ldots \to \Ker(d_{A_2^\bullet}^p) \to \Ker(d_{A_1^\bullet}^p)
\to \Ker(d_{A_0^\bullet}^p) \to \Ker(d_{M^\bullet}^p) \to 0
$$
soient également exacts. Posons $A^{p, q} = A_{-p}^q$ pour obtenir un
bicomplexe. Alors $\text{Tot}(A^{\bullet, \bullet}) \to M^\bullet$
induit par $A_0^\bullet \to M^\bullet$ est un quasi-isomorphisme.
\end{lemma}

\begin{proof}
En utilisant les suites exactes courtes
$0 \to \Ker(d^p_{A_n^\bullet}) \to A_n^p \to \Im(d^p_{A_n^\bullet}) \to 0$
et les hypothèses, nous voyons que
$$
\ldots \to \Im(d_{A_2^\bullet}^p) \to \Im(d_{A_1^\bullet}^p)
\to \Im(d_{A_0^\bullet}^p) \to \Im(d_{M^\bullet}^p) \to 0
$$
est exact pour tout $p \in \mathbf{Z}$. En répétant l'argument avec les
suites exactes
$0 \to \Im(d^{p - 1}_{A_n^\bullet}) \to \Ker(d^p_{A_n^\bullet})
\to H^p(A_n^\bullet) \to 0$, nous trouvons que
$$
\ldots \to H^p(A_2^\bullet) \to H^p(A_1^\bullet)
\to H^p(A_0^\bullet) \to H^p(M^\bullet) \to 0
$$
est exact pour tout $p \in \mathbf{Z}$.

\medskip\noindent
Écrivons $T^\bullet = \text{Tot}(A^{\bullet, \bullet})$. Nous allons montrer que
$H^0(T^\bullet) \to H^0(M^\bullet)$ est un isomorphisme. Le même argument
vaut dans les autres degrés. Soit $x \in \Ker(\text{d}_{T^\bullet}^0)$ un
représentant d'un élément $\xi \in H^0(T^\bullet)$.
Écrivons $x = \sum_{i = n, \ldots, 0} x_i$ avec $x_i \in A_i^i$.
Supposons $n > 0$. Alors $x_n$ appartient au noyau de
$d_{A_n^\bullet}^n$ et s'envoie sur zéro dans $H^n(A_{n - 1}^\bullet)$,
car son image est, au signe près, le bord de $x_{n - 1}$.
D'après le premier paragraphe de la démonstration, nous trouvons que
$x_n \bmod \Im(d^{n - 1}_{A_n^\bullet})$
appartient à l'image de $H^n(A_{n + 1}^\bullet) \to H^n(A_n^\bullet)$.
Nous pouvons donc modifier $x$ par un bord et nous ramener au cas où
$x_n$ est un bord. En modifiant encore une fois $x$, nous voyons que
nous pouvons supposer $x_n = 0$. Par récurrence, nous voyons que toute
classe de cohomologie $\xi$ est représentée par un cocycle $x = x_0$.
Enfin, la condition d'exactitude des noyaux nous dit que deux tels cocycles
$x_0$ et $x_0'$ sont cohomologues si et seulement si leurs images dans
$H^0(M^\bullet)$ sont égales.
\end{proof}

\begin{lemma}
\label{lemma-good-right-resolution-gives-qis}
Soit $M^\bullet$ un complexe de groupes abéliens. Soit
$$
0 \to M^\bullet \to A_0^\bullet \to A_1^\bullet \to A_2^\bullet \to \ldots
$$
un complexe exact de complexes de groupes abéliens
tel que, pour tout $p \in \mathbf{Z}$, les complexes
$$
0 \to
\Coker(d_{M^\bullet}^p) \to
\Coker(d_{A_0^\bullet}^p) \to
\Coker(d_{A_1^\bullet}^p) \to
\Coker(d_{A_2^\bullet}^p) \to \ldots
$$
soient également exacts. Posons $A^{p, q} = A_p^q$ pour obtenir un
bicomplexe. Soit $\text{Tot}_\pi(A^{\bullet, \bullet})$ le
complexe simple par produits associé au bicomplexe
(voir la démonstration). Alors l'application
$M^\bullet \to \text{Tot}_\pi(A^{\bullet, \bullet})$
induite par $M^\bullet \to A_0^\bullet$ est un quasi-isomorphisme.
\end{lemma}

\begin{proof}
En abrégeant $T^\bullet = \text{Tot}_\pi(A^{\bullet, \bullet})$,
nous posons
$$
T^n = \prod\nolimits_{p + q = n} A^{p, q} =
\prod\nolimits_{p + q = n} A_p^q
\quad\text{avec}\quad
\text{d}_{T^\bullet}^n =
\prod\nolimits_{n = p + q} (f_p^q + (-1)^pd_{A_p^\bullet}^q)
$$
où $f_p^\bullet : A_p^\bullet \to A_{p + 1}^\bullet$
sont les applications de complexes du lemme.

\medskip\noindent
Nous allons montrer que $H^0(M^\bullet) \to H^0(T^\bullet)$ est un
isomorphisme. Le même argument vaut dans les autres degrés.
Soit $x \in \Ker(\text{d}_{T^\bullet}^0)$ un représentant de
$\xi \in H^0(T^\bullet)$.
Écrivons $x = (x_i)$ avec $x_i \in A_i^{-i}$.
Remarquons que $x_0$ s'envoie sur zéro dans
$\Coker(A_1^{-1} \to A_1^0)$.
Nous voyons donc que $x_0 = m_0 + d_{A_0^\bullet}^{-1}(y)$ pour certains
$m_0 \in M^0$ et $y \in A_0^{-1}$.
Alors $d_{M^\bullet}(m_0) = 0$ parce que
$\text{d}_{A_0^\bullet}(x_0) = 0$, puisque
$\text{d}_{T^\bullet}(x) = 0$.
Ainsi, après avoir remplacé $\xi$ par un élément de l'image de
$H^0(M^\bullet) \to H^0(T^\bullet)$, nous pouvons supposer que $x_0$
appartient à $\Im(d^{-1}_{A_0^\bullet})$.

\medskip\noindent
Supposons $x_0 \in \Im(d^{-1}_{A_0^\bullet})$. Nous affirmons que, dans ce
cas, $\xi = 0$. Pour le prouver, construisons par récurrence sur $n$ des
éléments $y_0, y_1, \ldots, y_n$ avec $y_i \in A_i^{-i - 1}$ tels que
$x_0 = \text{d}_{A_0}^{-1}(y_0)$ et
$x_j = f_{j - 1}^{-j}(y_{j - 1}) + (-1)^j d^{-j - 1}_{A_j^\bullet}(y_j)$
pour $j = 1, \ldots, n$. Cela est clair pour $n = 0$. Démontrons l'étape
de récurrence : supposons avoir trouvé $y_0, \ldots, y_{n - 1}$. Alors
$w_n = x_n - f_{n - 1}^{-n}(y_{n - 1})$ appartient au noyau de
$d^{-n}_{A_n^\bullet}$ et s'envoie sur zéro dans $H^{-n}(A_{n + 1}^\bullet)$
(car son image est, au signe près, le bord de $x_{n + 1}$). Exactement
comme dans la démonstration du
Lemme \ref{lemma-good-resolution-gives-qis},
les hypothèses du lemme impliquent que
$$
0 \to
H^p(M^\bullet) \to
H^p(A_0^\bullet) \to
H^p(A_1^\bullet) \to
H^p(A_2^\bullet) \to \ldots
$$
est exact pour tout $p \in \mathbf{Z}$. Ainsi, après avoir modifié
$y_{n - 1}$ par un élément de $\Ker(d^{-n}_{A_{n - 1}^\bullet})$,
nous pouvons supposer que $w_n$ s'envoie sur zéro dans
$H^{-n}(A_n^\bullet)$. Cela signifie que nous pouvons trouver $y_n$ comme
voulu. Observons que cette procédure ne modifie pas
$y_0, \ldots, y_{n - 2}$. En poursuivant ainsi indéfiniment, nous trouvons
un élément $y = (y_i)$ de $T^{-1}$ tel que $d_{T^\bullet}(y) = x$.
Ceci montre que $H^0(M^\bullet) \to H^0(T^\bullet)$ est surjective.

\medskip\noindent
Supposons que $m_0 \in \Ker(d^0_{M^\bullet})$ s'envoie sur zéro dans
$H^0(T^\bullet)$. Disons que son image est la différentielle appliquée
à $y = (y_i) \in T^{-1}$.
Alors $y_0 \in A_0^{-1}$ s'envoie sur zéro dans
$\Coker(d^{-2}_{A_1^\bullet})$.
Par hypothèse, cela signifie que
$y_0 \bmod \Im(d^{-2}_{A_0^\bullet})$
est l'image d'un certain $z \in M^{-1}$. Il s'ensuit que
$m_0 = d^{-1}_{M^\bullet}(z)$. Cela prouve l'injectivité et achève la
démonstration.
\end{proof}

\begin{lemma}
\label{lemma-resolution-gives-qis}
Soit $M^\bullet$ un complexe de groupes abéliens. Soit
$$
\ldots \to A_2^\bullet \to A_1^\bullet \to A_0^\bullet \to M^\bullet \to 0
$$
un complexe exact de complexes de groupes abéliens. Posons
$A^{p, q} = A_{-p}^q$ pour obtenir un bicomplexe. Soit
$\text{Tot}_\pi(A^{\bullet, \bullet})$
le complexe simple par produits associé au bicomplexe (voir la démonstration).
Alors l'application $\text{Tot}_\pi(A^{\bullet, \bullet}) \to M^\bullet$
induite par $A_0^\bullet \to M^\bullet$ est un quasi-isomorphisme.
\end{lemma}

\begin{proof}
En abrégeant $T^\bullet = \text{Tot}_\pi(A^{\bullet, \bullet})$,
nous posons
$$
T^n = \prod\nolimits_{p + q = n} A^{p, q} =
\prod\nolimits_{p + q = n} A_{-p}^q
\quad\text{avec}\quad
\text{d}_{T^\bullet}^n =
\prod\nolimits_{n = p + q} (f_{-p}^q + (-1)^pd_{A_{-p}^\bullet}^q)
$$
où $f_p^\bullet : A_p^\bullet \to A_{p - 1}^\bullet$
sont les applications de complexes du lemme.
Nous allons montrer que $T^\bullet$ est acyclique lorsque
$M^\bullet$ est le complexe nul. Cela suffira grâce à l'astuce suivante.
Posons $B_{n + 1}^\bullet = A_n^\bullet$ pour $n \geq 0$
et $B_0^\bullet = M^\bullet$. Nous avons alors une suite exacte
$$
\ldots \to B_2^\bullet \to B_1^\bullet \to B_0^\bullet \to 0 \to 0
$$
comme dans le lemme. Soit $S^\bullet = \text{Tot}_\pi(B^{\bullet, \bullet})$.
Il existe alors une suite exacte courte évidente de complexes
$$
0 \to M^\bullet \to S^\bullet \to T^\bullet[1] \to 0
$$
et nous concluons par la suite exacte longue de cohomologie. Certains
détails sont omis.

\medskip\noindent
Supposons $M^\bullet = 0$. Nous allons montrer $H^0(T^\bullet) = 0$. Le même
argument vaut dans les autres degrés. Soit
$x =(x_n) \in \Ker(d_{T^\bullet})$ d'image $\xi \in H^0(T^\bullet)$,
avec $x_n \in A^{-n, n} = A_n^n$.
Puisque $M^0 = 0$, nous trouvons que $x_0 = f_1^0(y_0)$ pour un certain
$y_0 \in A_1^0$.
Alors $x_1 + d^0_{A_1^\bullet}(y_0) = f_2^1(y_1)$ pour un certain
$y_1 \in A_2^1$, parce que cet élément est envoyé sur zéro par $f_1^1$,
puisque $x$ est un cocycle. En poursuivant et en raisonnant
par récurrence, nous trouvons
$y = (y_i) \in T^{-1}$ tel que $d_{T^\bullet}(y) = x$, comme voulu.
\end{proof}

\section{Injectifs}
\label{section-injectives}

\begin{definition}
\label{definition-injective}
Soit $\mathcal{A}$ une catégorie abélienne.
Un objet $J \in \Ob(\mathcal{A})$ est dit
{\it injectif} si, pour toute injection
$A \hookrightarrow B$ et tout morphisme
$A \to J$, il existe un morphisme $B \to J$ rendant
commutatif le diagramme suivant
$$
\xymatrix{
A \ar[r] \ar[d] & B \ar@{-->}[ld] \\
J &
}
$$
\end{definition}

\noindent
Voici la caractérisation incontournable des objets injectifs.

\begin{lemma}
\label{lemma-characterize-injectives}
Soit $\mathcal{A}$ une catégorie abélienne.
Soit $I$ un objet de $\mathcal{A}$.
Les assertions suivantes sont équivalentes :
\begin{enumerate}
\item l'objet $I$ est injectif ;
\item le foncteur $B \mapsto \Hom_\mathcal{A}(B, I)$
est exact ;
\item toute suite exacte courte
$$
0 \to I \to A \to B \to 0
$$
dans $\mathcal{A}$ est scindée ;
\item on a $\Ext_\mathcal{A}(B, I) = 0$ pour
tout $B \in \Ob(\mathcal{A})$.
\end{enumerate}
\end{lemma}

\begin{proof}
Omis.
\end{proof}

\begin{lemma}
\label{lemma-product-injectives}
Soit $\mathcal{A}$ une catégorie abélienne.
Supposons que $I_\omega$, $\omega \in \Omega$, soit une famille d'objets
injectifs de $\mathcal{A}$. Si $\prod_{\omega \in \Omega} I_\omega$
existe, alors il est injectif.
\end{lemma}

\begin{proof}
Omis.
\end{proof}

\begin{definition}
\label{definition-enough-injectives}
Soit $\mathcal{A}$ une catégorie abélienne.
Nous disons que $\mathcal{A}$ possède {\it suffisamment d'injectifs}
si tout objet $A$ admet un morphisme injectif
$A \to J$ vers un objet injectif $J$.
\end{definition}

\begin{definition}
\label{definition-functorial-injective-embedding}
Soit $\mathcal{A}$ une catégorie abélienne.
Nous disons que $\mathcal{A}$ possède des
{\it plongements injectifs fonctoriels}
s'il existe un foncteur
$$
J : \mathcal{A} \longrightarrow \text{Arrows}(\mathcal{A})
$$
tel que
\begin{enumerate}
\item $s \circ J = \text{id}_\mathcal{A}$ ;
\item pour tout objet $A \in \Ob(\mathcal{A})$,
le morphisme $J(A)$ soit injectif ; et
\item pour tout objet $A \in \Ob(\mathcal{A})$,
l'objet $t(J(A))$ soit un objet injectif de $\mathcal{A}$.
\end{enumerate}
Nous noterons un tel foncteur
$A \mapsto (A \to J(A))$.
\end{definition}

\section{Projectifs}
\label{section-projectives}

\begin{definition}
\label{definition-projective}
Soit $\mathcal{A}$ une catégorie abélienne.
Un objet $P \in \Ob(\mathcal{A})$ est dit
{\it projectif} si, pour toute surjection
$A \rightarrow B$ et tout morphisme
$P \to B$, il existe un morphisme $P \to A$ rendant
commutatif le diagramme suivant
$$
\xymatrix{
A \ar[r] & B \\
P \ar@{-->}[u] \ar[ru] &
}
$$
\end{definition}

\noindent
Voici la caractérisation incontournable des objets projectifs.

\begin{lemma}
\label{lemma-characterize-projectives}
Soit $\mathcal{A}$ une catégorie abélienne.
Soit $P$ un objet de $\mathcal{A}$.
Les assertions suivantes sont équivalentes :
\begin{enumerate}
\item l'objet $P$ est projectif ;
\item le foncteur $B \mapsto \Hom_\mathcal{A}(P, B)$
est exact ;
\item toute suite exacte courte
$$
0 \to A \to B \to P \to 0
$$
dans $\mathcal{A}$ est scindée ;
\item on a $\Ext_\mathcal{A}(P, A) = 0$ pour
tout $A \in \Ob(\mathcal{A})$.
\end{enumerate}
\end{lemma}

\begin{proof}
Omis.
\end{proof}

\begin{lemma}
\label{lemma-coproduct-projectives}
Soit $\mathcal{A}$ une catégorie abélienne.
Supposons que $P_\omega$, $\omega \in \Omega$, soit une famille d'objets
projectifs de $\mathcal{A}$. Si $\coprod_{\omega \in \Omega} P_\omega$
existe, alors il est projectif.
\end{lemma}

\begin{proof}
Omis.
\end{proof}

\begin{definition}
\label{definition-enough-projectives}
Soit $\mathcal{A}$ une catégorie abélienne.
Nous disons que $\mathcal{A}$ possède {\it suffisamment de projectifs}
si tout objet $A$ admet un morphisme surjectif
$P \to A$ depuis un objet projectif $P$.
\end{definition}

\begin{definition}
\label{definition-functorial-projective-surjections}
Soit $\mathcal{A}$ une catégorie abélienne.
Nous disons que $\mathcal{A}$ possède des
{\it surjections projectives fonctorielles}
s'il existe un foncteur
$$
P : \mathcal{A} \longrightarrow \text{Arrows}(\mathcal{A})
$$
tel que
\begin{enumerate}
\item $t \circ P = \text{id}_\mathcal{A}$ ;
\item pour tout objet $A \in \Ob(\mathcal{A})$,
le morphisme $P(A)$ soit surjectif ; et
\item pour tout objet $A \in \Ob(\mathcal{A})$,
l'objet $s(P(A))$ soit un objet projectif de $\mathcal{A}$.
\end{enumerate}
Nous noterons un tel foncteur
$A \mapsto (P(A) \to A)$.
\end{definition}

\section{Injectifs et foncteurs adjoints}
\label{section-adjoint}

\noindent
Voici quelques lemmes sur les foncteurs adjoints et leur relation avec
les injectifs. Voir aussi le Lemme \ref{lemma-adjoint-get-abelian}.

\begin{lemma}
\label{lemma-adjoint-preserve-injectives}
\begin{slogan}
Un foncteur qui admet un adjoint à gauche exact préserve les injectifs
\end{slogan}
Soient $\mathcal{A}$ et $\mathcal{B}$ des catégories abéliennes.
Soient $u : \mathcal{A} \to \mathcal{B}$ et
$v : \mathcal{B} \to \mathcal{A}$ des foncteurs additifs tels que
$u$ soit adjoint à droite de $v$. Considérons les conditions suivantes :
\begin{enumerate}
\item[(a)] $v$ transforme les morphismes injectifs en morphismes injectifs ;
\item[(b)] $v$ est exact ; et
\item[(c)] $u$ transforme les injectifs en injectifs.
\end{enumerate}
Alors (a) $\Leftrightarrow$ (b) $\Rightarrow$ (c). Si $\mathcal{A}$
possède suffisamment d'injectifs, les trois conditions sont équivalentes.
\end{lemma}

\begin{proof}
Observons que $v$ est exact à droite en tant qu'adjoint à gauche
(Catégories, Lemme
\ref{categories-lemma-exact-adjoint}). Conjointement au
Lemme \ref{lemma-exact-functor}, cela explique pourquoi
(a) $\Leftrightarrow$ (b).

\medskip\noindent
Supposons (a). Soit $I$ un objet injectif de $\mathcal{A}$.
Soit $\varphi : N \to M$ un morphisme injectif dans $\mathcal{B}$ et soit
$\alpha : N \to uI$ un morphisme.
Par adjonction, nous obtenons un morphisme $\alpha : vN \to I$, et
$v\varphi : vN \to vM$ est injectif par hypothèse.
Comme $I$ est un objet injectif, nous obtenons donc un morphisme
$\beta : vM \to I$ prolongeant $\alpha$. Une nouvelle application de
l'adjonction lui fait correspondre un morphisme $\beta : M \to uI$
prolongeant $\alpha$. Ainsi (c) est vraie.

\medskip\noindent
Supposons que $\mathcal{A}$ possède suffisamment d'injectifs et que (c)
soit vérifiée. Soit $f : B \to B'$ un monomorphisme dans $\mathcal{B}$,
et soit $A = \Ker(v(f))$. Choisissons un monomorphisme $g : A \to I$ avec
$I$ injectif (il existe par hypothèse). Alors $g$ se prolonge en
$g' : v(B) \to I$, d'où, par adjonction, un morphisme $B \to u(I)$.
Puisque $u(I)$ est injectif, ce morphisme se prolonge en
$h : B' \to u(I)$, d'où, par adjonction, un morphisme
$k : v(B') \to I$ prolongeant $g'$. Mais alors $k$ « prolonge » $g$,
ce qui force $A = 0$ puisque $g$ était un monomorphisme. Ainsi (a) est vraie.
\end{proof}

\begin{remark}
\label{remark-need-left-exactness}
Soit $R \to S$ un morphisme d'anneaux. Soit
$u : \text{Mod}_S \to \text{Mod}_R$ défini par $u(N) = N_R$, et soit
$v : \text{Mod}_R \to \text{Mod}_S$ défini par
$v(M) = M \otimes_R S$. Alors $u$ est adjoint à droite de $v$, le foncteur
$u$ est exact et $v$ est exact à droite. Mais les conditions (a), (b), (c) du
Lemme \ref{lemma-adjoint-preserve-injectives} ne sont pas vérifiées en général.
Par exemple, si $R = \mathbf{Z}$ et $S = \mathbf{Z}/p\mathbf{Z}$, alors
le $S$-module injectif $\mathbf{Z}/p\mathbf{Z}$ n'est pas
un $\mathbf{Z}$-module injectif. En fait, le lemme montre que
tous les $S$-modules injectifs sont injectifs comme $R$-modules si et seulement
si $R \to S$ est un morphisme plat d'anneaux.
\end{remark}

\begin{lemma}
\label{lemma-adjoint-enough-injectives}
Soient $\mathcal{A}$ et $\mathcal{B}$ des catégories abéliennes.
Soient $u : \mathcal{A} \to \mathcal{B}$ et
$v : \mathcal{B} \to \mathcal{A}$ des foncteurs additifs.
Supposons que
\begin{enumerate}
\item $u$ soit adjoint à droite de $v$ ;
\item $v$ transforme les morphismes injectifs en morphismes injectifs ;
\item $\mathcal{A}$ possède suffisamment d'injectifs ; et
\item $vB = 0$ implique $B = 0$ pour tout $B \in \Ob(\mathcal{B})$.
\end{enumerate}
Alors $\mathcal{B}$ possède suffisamment d'injectifs.
\end{lemma}

\begin{proof}
Choisissons $B \in \Ob(\mathcal{B})$.
Choisissons une injection $vB \to I$, où $I$
est un objet injectif de $\mathcal{A}$.
D'après le Lemme \ref{lemma-adjoint-preserve-injectives}
et les hypothèses, l'application correspondante
$B \to uI$ est l'injection de $B$ dans un objet injectif.
\end{proof}

\begin{remark}
\label{remark-faithfulness-needed}
Soient $\mathcal{A}$, $\mathcal{B}$,
$u : \mathcal{A} \to \mathcal{B}$ et
$v : \mathcal{B} \to \mathcal{A}$ comme dans le
Lemme \ref{lemma-adjoint-enough-injectives}.
En présence des conditions (1) et (2), la condition (4) équivaut à la
fidélité de $v$. En outre, la condition (4) est nécessaire.
Un exemple s'obtient en considérant le cas où les foncteurs $u$ et $v$
sont tous deux le foncteur nul.
\end{remark}

\begin{lemma}
\label{lemma-adjoint-functorial-injectives}
Soient $\mathcal{A}$ et $\mathcal{B}$ des catégories abéliennes.
Soient $u : \mathcal{A} \to \mathcal{B}$ et
$v : \mathcal{B} \to \mathcal{A}$ des foncteurs additifs.
Supposons que
\begin{enumerate}
\item $u$ soit adjoint à droite de $v$ ;
\item $v$ transforme les morphismes injectifs en morphismes injectifs ;
\item $\mathcal{A}$ possède suffisamment d'injectifs ;
\item $vB = 0$ implique $B = 0$ pour tout $B \in \Ob(\mathcal{B})$ ; et
\item $\mathcal{A}$ possède des plongements injectifs fonctoriels.
\end{enumerate}
Alors $\mathcal{B}$ possède des plongements injectifs fonctoriels.
\end{lemma}

\begin{proof}
Soit $A \mapsto (A \to J(A))$ un plongement injectif fonctoriel
sur $\mathcal{A}$. Alors
$B \mapsto (B \to uJ(vB))$ est un plongement injectif fonctoriel
sur $\mathcal{B}$. Comparer avec la démonstration du
Lemme \ref{lemma-adjoint-enough-injectives}.
\end{proof}

\begin{lemma}
\label{lemma-partially-defined-adjoint}
Soient $\mathcal{A}$ et $\mathcal{B}$ des catégories abéliennes.
Soit $u : \mathcal{A} \to \mathcal{B}$ un foncteur.
S'il existe un sous-ensemble $\mathcal{P} \subset \Ob(\mathcal{B})$
tel que
\begin{enumerate}
\item tout objet de $\mathcal{B}$ soit un quotient d'un élément
de $\mathcal{P}$ ; et
\item pour tout $P \in \mathcal{P}$, il existe un objet
$Q$ de $\mathcal{A}$ tel que
$\Hom_\mathcal{A}(Q, A) = \Hom_\mathcal{B}(P, u(A))$ fonctoriellement
en $A$,
\end{enumerate}
alors il existe un adjoint à gauche $v$ de $u$.
\end{lemma}

\begin{proof}
Par le lemme de Yoneda (Catégories, Lemme \ref{categories-lemma-yoneda}),
l'objet $Q$ de $\mathcal{A}$ correspondant à $P$ est défini à isomorphisme
unique près par la formule
$\Hom_\mathcal{A}(Q, A) = \Hom_\mathcal{B}(P, u(A))$. Écrivons
$Q = v(P)$. Notons $i_P : P \to u(v(P))$ l'application correspondant à
$\text{id}_{v(P)}$ dans $\Hom_\mathcal{A}(v(P), v(P))$. La fonctorialité
de (2) implique que la bijection est donnée par
$$
\Hom_\mathcal{A}(v(P), A) \to \Hom_\mathcal{B}(P, u(A)),\quad
\varphi \mapsto u(\varphi) \circ i_P
$$
Pour toute paire d'éléments $P_1, P_2 \in \mathcal{P}$, il existe une
application canonique
$$
\Hom_\mathcal{B}(P_2, P_1)
\to
\Hom_\mathcal{A}(v(P_2), v(P_1)),\quad
\varphi \mapsto v(\varphi)
$$
caractérisée par la règle
$u(v(\varphi)) \circ i_{P_2} = i_{P_1} \circ \varphi$ dans
$\Hom_\mathcal{B}(P_2, u(v(P_1)))$.
Remarquons que $\varphi \mapsto v(\varphi)$ est
compatible à la composition ; cela se voit directement à partir de la
caractérisation. Ainsi $P \mapsto v(P)$ est un foncteur de la sous-catégorie
pleine de $\mathcal{B}$ dont les objets sont les éléments de $\mathcal{P}$.

\medskip\noindent
Étant donné un objet arbitraire $B$ de $\mathcal{B}$, choisissons une suite
exacte
$$
P_2 \to P_1 \to B \to 0
$$
ce qui est possible par l'hypothèse (1). Définissons $v(B)$ comme l'objet de
$\mathcal{A}$ qui s'insère dans la suite exacte
$$
v(P_2) \to v(P_1) \to v(B) \to 0
$$
Alors
\begin{align*}
\Hom_\mathcal{A}(v(B), A)
& =
\Ker(\Hom_\mathcal{A}(v(P_1), A) \to \Hom_\mathcal{A}(v(P_2), A)) \\
& =
\Ker(\Hom_\mathcal{B}(P_1, u(A)) \to \Hom_\mathcal{B}(P_2, u(A))) \\
& =
\Hom_\mathcal{B}(B, u(A))
\end{align*}
Nous voyons donc que nous pouvons prendre
$\mathcal{P} = \Ob(\mathcal{B})$, c'est-à-dire que $v$ est défini partout.
\end{proof}

\section{Systèmes essentiellement constants}
\label{section-essentially-constant}

\noindent
Dans cette section, nous étudions les systèmes essentiellement constants
à valeurs dans des catégories additives.

\begin{lemma}
\label{lemma-essentially-constant-into-karoubian}
Soit $\mathcal{I}$ une catégorie, soit $\mathcal{A}$ une catégorie
préadditive karoubienne et soit $M : \mathcal{I} \to \mathcal{A}$ un diagramme.
\begin{enumerate}
\item Supposons $\mathcal{I}$ filtrante. Les assertions suivantes sont
équivalentes :
\begin{enumerate}
\item $M$ est essentiellement constant ;
\item $X = \colim M$ existe et il existe une sous-catégorie filtrante cofinale
$\mathcal{I}' \subset \mathcal{I}$ ainsi que, pour
$i' \in \Ob(\mathcal{I}')$, une décomposition en somme directe
$M_{i'} = X_{i'} \oplus Z_{i'}$ telle que $X_{i'}$ s'envoie
isomorphiquement sur $X$ et que $Z_{i'}$ s'envoie sur zéro dans $M_{i''}$
pour un certain $i' \to i''$ dans $\mathcal{I}'$.
\end{enumerate}
\item Supposons $\mathcal{I}$ cofiltrante. Les assertions suivantes sont
équivalentes :
\begin{enumerate}
\item $M$ est essentiellement constant ;
\item $X = \lim M$ existe et il existe une sous-catégorie cofiltrante initiale
$\mathcal{I}' \subset \mathcal{I}$ ainsi que, pour
$i' \in \Ob(\mathcal{I}')$, une décomposition en somme directe
$M_{i'} = X_{i'} \oplus Z_{i'}$
telle que $X$ s'envoie isomorphiquement sur $X_{i'}$ et que
$M_{i''} \to Z_{i'}$ soit nul pour un certain $i'' \to i'$ dans
$\mathcal{I}'$.
\end{enumerate}
\end{enumerate}
\end{lemma}

\begin{proof}
Supposons (1)(a), c'est-à-dire que $\mathcal{I}$ soit filtrante et que $M$
soit essentiellement constant. Soit $X = \colim M_i$. Choisissons $i$ et
$X \to M_i$ comme dans
Catégories, Définition
\ref{categories-definition-essentially-constant-diagram}.
Soit $\mathcal{I}'$ la sous-catégorie pleine formée des objets qui sont le but
d'un morphisme de source $i$.
Supposons $i' \in \Ob(\mathcal{I}')$ et choisissons un morphisme $i \to i'$.
Alors la composée $X \to M_i \to M_{i'}$ avec $M_{i'} \to X$ est l'identité
sur $X$. Comme $\mathcal{A}$ est karoubienne, nous obtenons une décomposition
en somme directe $M_{i'} = X_{i'} \oplus Z_{i'}$, où
$Z_{i'} = \Ker(M_{i'} \to X)$ et $X_{i'}$ s'envoie isomorphiquement sur $X$.
Choisissons $i \to k$ et $i' \to k$ tels que
$M_{i'} \to X \to M_i \to M_k$
soit égal à $M_{i'} \to M_k$, comme dans Catégories, Définition
\ref{categories-definition-essentially-constant-diagram}.
Nous voyons alors que $M_{i'} \to M_k$ annule $Z_{i'}$.
Ainsi (1)(b) est vérifiée.

\medskip\noindent
Supposons (1)(b), c'est-à-dire que $\mathcal{I}$ soit filtrante et que nous
ayons $\mathcal{I}' \subset \mathcal{I}$ ainsi que, pour
$i' \in \Ob(\mathcal{I}')$, une décomposition en somme directe
$M_{i'} = X_{i'} \oplus Z_{i'}$
comme dans l'énoncé du lemme. Pour voir que $M$ est essentiellement constant,
nous pouvons remplacer $\mathcal{I}$ par $\mathcal{I}'$, voir
Catégories, Lemme \ref{categories-lemma-cofinal-essentially-constant}.
Choisissons un $i \in \Ob(\mathcal{I})$ quelconque
et notons $X \to M_i$ la composée de l'inverse de l'isomorphisme
$X_i \to X$ et de l'inclusion $X_i \to M_i$. Si $j$ est
un second objet, choisissons $j \to k$ tel que $Z_j \to M_k$ soit
nul. Comme $\mathcal{I}$ est filtrante, nous pouvons aussi supposer qu'il
existe un morphisme $i \to k$ (quitte à agrandir $k$). Alors
$M_j \to X \to M_i \to M_k$ et $M_j \to M_k$ annulent tous deux $Z_j$.
Ainsi, après composition à droite par un morphisme $M_k \to M_l$ qui annule
le facteur direct $Z_k$, nous trouvons que
$M_j \to X \to M_i \to M_l$ et
$M_j \to M_l$ sont égaux, c'est-à-dire que $M$ est essentiellement constant.

\medskip\noindent
La démonstration de (2) est duale.
\end{proof}

\begin{lemma}
\label{lemma-direct-sum-from-product-colimit}
Soit $\mathcal{I}$ une catégorie. Soit $\mathcal{A}$ une catégorie additive
et karoubienne. Soient $F : \mathcal{I} \to \mathcal{A}$ et
$G : \mathcal{I} \to \mathcal{A}$ des foncteurs. Les assertions suivantes
sont équivalentes :
\begin{enumerate}
\item $\colim_\mathcal{I}(F \oplus G)$ existe ; et
\item $\colim_\mathcal{I} F$ et $\colim_\mathcal{I} G$ existent.
\end{enumerate}
Dans ce cas, $\colim_\mathcal{I}(F \oplus G) =
\colim_\mathcal{I} F \oplus \colim_\mathcal{I} G$.
\end{lemma}

\begin{proof}
Supposons (1). Posons $W = \colim_\mathcal{I}(F \oplus G)$.
Remarquons que la projection sur $F$ définit une transformation naturelle
$F \oplus G \to F \oplus G$ qui est idempotente. Nous obtenons donc
un endomorphisme idempotent $W \to W$ par
Catégories, Lemme \ref{categories-lemma-functorial-colimit}.
Comme $\mathcal{A}$ est karoubienne, nous obtenons une décomposition
correspondante en somme directe $W = X \oplus Y$, voir le
Lemme \ref{lemma-karoubian}.
Un argument immédiat (omis) montre que
$X = \colim_\mathcal{I} F$ et $Y = \colim_\mathcal{I} G$.
Ainsi (2) est vérifiée. Nous omettons la démonstration que (2) implique (1).
\end{proof}

\begin{lemma}
\label{lemma-direct-sum-from-product-essentially-constant}
Soit $\mathcal{I}$ une catégorie filtrante. Soit $\mathcal{A}$
une catégorie additive et karoubienne. Soient
$F : \mathcal{I} \to \mathcal{A}$ et
$G : \mathcal{I} \to \mathcal{A}$ des foncteurs. Les assertions suivantes
sont équivalentes :
\begin{enumerate}
\item $F \oplus G : \mathcal{I} \to \mathcal{A}$
est essentiellement constant ; et
\item $F$ et $G$ sont essentiellement constants.
\end{enumerate}
\end{lemma}

\begin{proof}
Supposons (1). En particulier,
$W = \colim_\mathcal{I}(F \oplus G)$ existe et, par le
Lemme \ref{lemma-direct-sum-from-product-colimit},
nous avons $W = X \oplus Y$ avec $X = \colim_\mathcal{I} F$ et
$Y = \colim_\mathcal{I} G$. Un argument immédiat (omis), utilisant par
exemple la caractérisation de
Catégories, Lemme \ref{categories-lemma-characterize-essentially-constant-ind},
montre que $F$ est essentiellement constant de valeur $X$ et que $G$ est
essentiellement constant de valeur $Y$. Ainsi (2) est vérifiée.
La démonstration que (2) implique (1) est omise.
\end{proof}

\section{Systèmes projectifs}
\label{section-inverse-systems}

\noindent
Soit $\mathcal{C}$ une catégorie.
Dans Catégories, section \ref{categories-section-posets-limits},
nous avons défini la notion de système projectif indexé par un ensemble
préordonné (à valeurs dans la catégorie $\mathcal{C}$).
Si l'ensemble préordonné est $\mathbf{N} = \{1, 2, 3, \ldots\}$
muni de l'ordre usuel, un tel système projectif indexé par $\mathbf{N}$
est souvent appelé simplement un {\it système projectif}. Il consiste tout
simplement en une paire $(M_i, f_{ii'})$, où chaque $M_i$,
$i \in \mathbf{N}$, est un objet de $\mathcal{C}$ et, pour chaque
$i > i'$, $i, i' \in \mathbf{N}$, on a un morphisme
$f_{ii'} : M_i \to M_{i'}$ tel que, de plus,
$f_{i'i''} \circ f_{ii'} = f_{ii''}$ chaque fois que cela a un sens.
Il est clair qu'il suffit en fait de donner les morphismes
$M_2 \to M_1$, $M_3 \to M_2$, et ainsi de suite. Un système projectif
est donc souvent représenté de la manière suivante
$$
M_1 \xleftarrow{\varphi_2} M_2 \xleftarrow{\varphi_3} M_3 \leftarrow \ldots
$$
En outre, nous omettons souvent de la notation les applications de transition
$\varphi_i$ et disons simplement « soit $(M_i)$ un système projectif ».

\medskip\noindent
La collection de tous les systèmes projectifs à valeurs dans
$\mathcal{C}$ forme une catégorie, avec la notion évidente de morphisme.

\begin{lemma}
\label{lemma-inverse-systems-abelian}
Soit $\mathcal{C}$ une catégorie.
\begin{enumerate}
\item Si $\mathcal{C}$ est une catégorie additive, alors la catégorie
des systèmes projectifs à valeurs dans $\mathcal{C}$ est additive.
\item Si $\mathcal{C}$ est une catégorie abélienne, alors la catégorie
des systèmes projectifs à valeurs dans $\mathcal{C}$ est abélienne.
Une suite $(K_i) \to (L_i) \to (M_i)$ de systèmes projectifs
est exacte si et seulement si chaque $K_i \to L_i \to M_i$ est exacte.
\end{enumerate}
\end{lemma}

\begin{proof}
Omis.
\end{proof}

\noindent
La limite (voir Catégories, section \ref{categories-section-posets-limits})
d'un tel système projectif se note $\lim M_i$, ou $\lim_i M_i$.
Si $\mathcal{C}$ est la catégorie des groupes abéliens (ou des ensembles),
la limite existe toujours et peut en fait être décrite comme suit
$$
\lim_i M_i
=
\{(x_i) \in \prod M_i \mid \varphi_i(x_i) = x_{i - 1}, \ i = 2, 3, \ldots\}
$$
voir Catégories, section \ref{categories-section-limit-sets}.
Cependant, étant donnée une suite exacte courte
$$
0 \to (A_i) \to (B_i) \to (C_i) \to 0
$$
de systèmes projectifs de groupes abéliens, la suite de limites associée
n'est pas toujours exacte. Pour approfondir cette question, introduisons
la notion suivante.

\begin{definition}
\label{definition-Mittag-Leffler}
Soit $\mathcal{C}$ une catégorie abélienne.
Nous disons que le système projectif $(A_i)$
satisfait la {\it condition de Mittag-Leffler}, ou, en abrégé, qu'il est
{\it ML}, si, pour tout $i$, il existe $c = c(i) \geq i$
tel que
$$
\Im(A_k \to A_i) = \Im(A_c \to A_i)
$$
pour tout $k \geq c$.
\end{definition}

\noindent
Il se trouve que la condition de Mittag-Leffler suffit à assurer
l'exactitude du foncteur $\lim$, pourvu que l'on travaille dans
la catégorie abélienne des groupes abéliens, des modules sur un anneau, etc.
A.\ Neeman montre dans un article (voir \cite{Neeman-Counterexample})
que cette condition n'est pas assez forte dans une catégorie abélienne
qui satisfait AB4* (c'est-à-dire dans laquelle les produits sont exacts).

\begin{lemma}
\label{lemma-Mittag-Leffler}
Soit
$$
0 \to (A_i) \to (B_i) \to (C_i) \to 0
$$
une suite exacte courte de systèmes projectifs de groupes abéliens.
\begin{enumerate}
\item Dans tous les cas, la suite
$$
0 \to \lim_i A_i \to \lim_i B_i \to \lim_i C_i
$$
est exacte.
\item Si $(B_i)$ est ML, alors $(C_i)$ est aussi ML.
\item Si $(A_i)$ est ML, alors
$$
0 \to \lim_i A_i \to \lim_i B_i \to \lim_i C_i \to 0
$$
est exacte.
\end{enumerate}
\end{lemma}

\begin{proof}
Bel exercice. Voir
Algèbre, Lemme \ref{algebra-lemma-Mittag-Leffler} pour la partie (3).
\end{proof}

\begin{lemma}
\label{lemma-apply-Mittag-Leffler}
Soit
$$
(A_i) \to (B_i) \to (C_i) \to (D_i)
$$
une suite exacte de systèmes projectifs de groupes abéliens. Si le
système $(A_i)$ est ML, alors la suite
$$
\lim_i B_i \to \lim_i C_i \to \lim_i D_i
$$
est exacte.
\end{lemma}

\begin{proof}
Soient $Z_i = \Ker(C_i \to D_i)$ et $I_i = \Im(A_i \to B_i)$.
Alors $\lim Z_i = \Ker(\lim C_i \to \lim D_i)$ et
nous obtenons une suite exacte courte de systèmes
$$
0 \to (I_i) \to (B_i) \to (Z_i) \to 0
$$
En outre, par le
Lemme \ref{lemma-Mittag-Leffler},
nous voyons que $(I_i)$ satisfait (ML) ; une nouvelle application du
Lemme \ref{lemma-Mittag-Leffler}
montre donc que $\lim B_i \to \lim Z_i$ est surjective, ce qui
démontre le lemme.
\end{proof}

\noindent
La caractérisation suivante des systèmes projectifs essentiellement constants
montre en particulier qu'ils satisfont ML.

\begin{lemma}
\label{lemma-essentially-constant}
Soit $\mathcal{A}$ une catégorie abélienne.
Soit $(A_i)$ un système projectif dans $\mathcal{A}$ de limite
$A = \lim A_i$.
Alors $(A_i)$ est essentiellement constant (voir
Catégories, Définition
\ref{categories-definition-essentially-constant-diagram})
si et seulement s'il existe $i$ et, pour tout $j \geq i$, une décomposition
en somme directe $A_j = A \oplus Z_j$ telle que
(a) les applications $A_{j'} \to A_j$ soient compatibles avec les
décompositions en sommes directes et (b), pour tout $j$, il existe
$j' \geq j$ tel que $Z_{j'} \to Z_j$ soit nul.
\end{lemma}

\begin{proof}
Supposons $(A_i)$ essentiellement constant. Il existe alors $i$ et
un morphisme $A_i \to A$ tels que $A \to A_i \to A$ soit l'identité et,
pour tout $j \geq i$, il existe
$j' \geq j$ tel que $A_{j'} \to A_j$ se factorise sous la forme
$A_{j'} \to A_i \to A \to A_j$ (la dernière application provient de
$A = \lim A_i$).
Ainsi, il convient de poser $Z_j = \Ker(A_j \to A)$ pour tout $j \geq i$.
La démonstration de la réciproque est omise.
\end{proof}

\noindent
Nous améliorerons le lemme suivant dans
Compléments sur l'algèbre, Lemme \ref{more-algebra-lemma-Mittag-Leffler}.

\begin{lemma}
\label{lemma-exact-sequence-ML}
Soit
$$
0 \to (A_i) \to (B_i) \to (C_i) \to 0
$$
une suite exacte de systèmes projectifs de groupes abéliens.
Si $(C_i)$ est essentiellement constant, alors $(A_i)$ satisfait ML
si et seulement si $(B_i)$ satisfait ML.
\end{lemma}

\begin{proof}
Après renumérotation, nous pouvons supposer que $C_i = C \oplus Z_i$
de manière compatible avec les applications de transition et que, pour tout
$i$, il existe $i' \geq i$ tel que $Z_{i'} \to Z_i$ soit nul, voir le
Lemme \ref{lemma-essentially-constant}.

\medskip\noindent
Supposons d'abord $C = 0$, c'est-à-dire $C_i = Z_i$.
Dans ce cas, choisissons $1 = n_1 < n_2 < n_3 < \ldots$
tel que $Z_{n_{i + 1}} \to Z_{n_i}$ soit nul.
Alors $B_{n_{i + 1}} \to B_{n_i}$ se factorise par
$A_{n_i} \subset B_{n_i}$.
Il s'ensuit que, pour $j \geq i + 1$, on a
$$
\Im(A_{n_j} \to A_{n_i}) \subset \Im(B_{n_j} \to B_{n_i}) \subset
\Im(A_{n_{j - 1}} \to A_{n_i})
$$
comme sous-objets de $A_{n_i}$.
Ainsi les images $\Im(A_{n_j} \to A_{n_i})$ se stabilisent pour
$j \geq i + 1$ si et seulement s'il en va de même des images
$\Im(B_{n_j} \to B_{n_i})$.
L'équivalence en résulte (un petit détail est omis).

\medskip\noindent
Si $C \not = 0$, notons $B'_i \subset B_i$ l'image réciproque
de $C$ par l'application $B_i \to C \oplus Z_i$. Le paragraphe précédent
montre alors que $(B'_i)$ satisfait ML si et seulement si $(B_i)$
satisfait ML. Nous pouvons donc remplacer $(B_i)$ par $(B'_i)$.
Dans ce cas, nous avons des suites exactes
$0 \to A_i \to B_i \to C \to 0$
pour tout $i$. Il s'ensuit que
$0 \to \Im(A_j \to A_i) \to \Im(B_j \to B_i) \to C \to 0$
est une suite exacte courte pour tout $j \geq i$. Les images
$\Im(A_j \to A_i)$ se stabilisent donc pour $j \geq i$
si et seulement s'il en va de même de $\Im(B_j \to B_i)$,
comme voulu.
\end{proof}

\noindent
La version « correcte » du lemme suivant est
Compléments sur l'algèbre, Lemme
\ref{more-algebra-lemma-apply-Mittag-Leffler-again}.

\begin{lemma}
\label{lemma-apply-Mittag-Leffler-again}
Soit
$$
(A^{-2}_i \to A^{-1}_i \to A^0_i \to A^1_i)
$$
un système projectif de complexes de groupes abéliens, et notons
$A^{-2} \to A^{-1} \to A^0 \to A^1$ sa limite. Notons
$(H_i^{-1})$, $(H_i^0)$ les systèmes projectifs de cohomologie, et
$H^{-1}$, $H^0$ les cohomologies de
$A^{-2} \to A^{-1} \to A^0 \to A^1$.
Si $(A^{-2}_i)$ et $(A^{-1}_i)$ sont ML et si
$(H^{-1}_i)$ est essentiellement constant, alors
$H^0 = \lim H_i^0$.
\end{lemma}

\begin{proof}
Soient $Z^j_i = \Ker(A^j_i \to A^{j + 1}_i)$ et
$I^j_i = \Im(A^{j - 1}_i \to A^j_i)$.
Remarquons que $\lim Z^0_i = \Ker(\lim A^0_i \to \lim A^1_i)$, puisque
la prise de noyaux commute aux limites.
Les systèmes $(I^{-1}_i)$ et $(I^0_i)$ sont ML en tant que quotients des
systèmes $(A^{-2}_i)$ et $(A^{-1}_i)$, voir le
Lemme \ref{lemma-Mittag-Leffler}.
Nous avons donc une suite exacte
$$
0 \to (I^{-1}_i) \to (Z^{-1}_i) \to (H^{-1}_i) \to 0
$$
de systèmes projectifs, où $(I^{-1}_i)$ est ML et où $(H^{-1}_i)$
est essentiellement constant par hypothèse.
Ainsi $(Z^{-1}_i)$ est ML par le
Lemme \ref{lemma-exact-sequence-ML}.
La suite exacte
$$
0 \to (Z^{-1}_i) \to (A^{-1}_i) \to (I^0_i) \to 0
$$
et une application du
Lemme \ref{lemma-Mittag-Leffler}
montrent que $\lim A^{-1}_i \to \lim I^0_i$ est surjective.
Enfin, la suite exacte
$$
0 \to (I^0_i) \to (Z^0_i) \to (H^0_i) \to 0
$$
et le
Lemme \ref{lemma-Mittag-Leffler}
montrent que $\lim I^0_i \to \lim Z^0_i \to \lim H^0_i \to 0$
est exacte. En rassemblant ces faits, nous obtenons le résultat.
\end{proof}

\noindent
Nous avons parfois besoin d'une version du lemme précédent où l'on prend
des limites indexées par de grands ordinaux.

\begin{lemma}
\label{lemma-ML-over-ordinals}
Soit $\alpha$ un ordinal. Soient $K_\beta^\bullet$, $\beta < \alpha$,
un système projectif de complexes de groupes abéliens indexé par $\alpha$.
Si, pour tout $\beta < \alpha$, le complexe $K_\beta^\bullet$ est acyclique
et si, pour tout $n \in \mathbf{Z}$, l'application
$$
K^n_\beta \longrightarrow \lim_{\gamma < \beta} K^n_\gamma
$$
est surjective, alors le complexe
$\lim_{\beta < \alpha} K_\beta^\bullet$ est acyclique.
\end{lemma}

\begin{proof}
Par récurrence transfinie, nous démontrons cette assertion pour tout ordinal
$\alpha$ et tout système comme dans le lemme. En particulier, pendant la
démonstration du résultat pour $\alpha$, nous pouvons supposer que les
complexes $\lim_{\gamma < \beta} K^\bullet_\gamma$ sont acycliques.

\medskip\noindent
Soit $x \in \lim_{\beta < \alpha} K^0_\beta$ avec $\text{d}(x) = 0$.
Nous allons trouver $y \in \lim_{\beta < \alpha} K^{-1}_\beta$
tel que $\text{d}(y) = x$.
Écrivons $x = (x_\beta)$, où $x_\beta \in K_\beta^0$ est l'image de $x$
pour $\beta < \alpha$. Nous allons construire $y = (y_\beta)$
par récurrence transfinie.

\medskip\noindent
Pour $\beta = 0$, soit $y_0 \in K_0^{-1}$
un élément quelconque tel que $\text{d}(y_0) = x_0$.

\medskip\noindent
Pour un ordinal successeur $\beta = \gamma + 1$, il nous faut trouver un
élément $y_\beta$ qui s'envoie à la fois sur $y_\gamma$ par l'application de
transition $f : K^\bullet_\beta \to K^\bullet_\gamma$ et sur $x_\beta$ par
la différentielle. Comme première approximation, choisissons $y'_\beta$ avec
$\text{d}(y'_\beta) = x_\beta$. Alors la différence
$y_\gamma - f(y'_\beta)$ appartient au noyau de la différentielle et est donc
égale à $\text{d}(z_\gamma)$ pour un certain $z_\gamma \in K^{-2}_\gamma$.
Par hypothèse, l'application $f^{-2} : K^{-2}_\beta \to K^{-2}_\gamma$
est surjective. Nous écrivons donc $z_\gamma = f(z_\beta)$
et remplaçons $y'_\beta$ par
$y_\beta = y'_\beta + \text{d}(z_\beta)$, ce qui convient.

\medskip\noindent
Si $\beta$ est un ordinal limite, nous avons l'élément
$(y_\gamma)_{\gamma < \beta}$ de $\lim_{\gamma < \beta} K^{-1}_\gamma$
dont la différentielle est l'image de $x_\beta$. Nous pouvons donc raisonner
exactement comme ci-dessus, en utilisant l'application de complexes
surjective terme à terme
$f : K_\beta^\bullet \to \lim_{\gamma < \beta} K_\gamma^\bullet$
et le fait (voir le premier paragraphe de la démonstration) que nous pouvons
supposer par récurrence que
$\lim_{\gamma < \beta} K_\gamma^\bullet$ est acyclique.
\end{proof}

\section{Exactitude des produits}
\label{section-product-exact}


\begin{lemma}
\label{lemma-product-abelian-groups-exact}
Soit $I$ un ensemble. Pour $i \in I$, soit $L_i \to M_i \to N_i$ un
complexe de groupes abéliens. Soit
$H_i = \Ker(M_i \to N_i)/\Im(L_i \to M_i)$
sa cohomologie. Alors
$$
\prod L_i \to \prod M_i \to \prod N_i
$$
est un complexe de groupes abéliens d'homologie $\prod H_i$.
\end{lemma}

\begin{proof}
Omis.
\end{proof}






\begin{multicols}{2}[\section{Autres chapitres}]
\noindent
Préliminaires
\begin{enumerate}
\item \hyperref[introduction-section-phantom]{Introduction}
\item \hyperref[conventions-section-phantom]{Conventions}
\item \hyperref[sets-section-phantom]{Théorie des ensembles}
\item \hyperref[categories-section-phantom]{Catégories}
\item \hyperref[topology-section-phantom]{Topologie}
\item \hyperref[sheaves-section-phantom]{Faisceaux sur les espaces}
\item \hyperref[sites-section-phantom]{Sites et faisceaux}
\item \hyperref[stacks-section-phantom]{Champs}
\item \hyperref[fields-section-phantom]{Corps}
\item \hyperref[algebra-section-phantom]{Algèbre commutative}
\item \hyperref[brauer-section-phantom]{Groupes de Brauer}
\item \hyperref[homology-section-phantom]{Algèbre homologique}
\item \hyperref[derived-section-phantom]{Catégories dérivées}
\item \hyperref[simplicial-section-phantom]{Méthodes simpliciales}
\item \hyperref[more-algebra-section-phantom]{Compléments d'algèbre}
\item \hyperref[smoothing-section-phantom]{Lissage des morphismes d'anneaux}
\item \hyperref[modules-section-phantom]{Faisceaux de modules}
\item \hyperref[sites-modules-section-phantom]{Modules sur les sites}
\item \hyperref[injectives-section-phantom]{Objets injectifs}
\item \hyperref[cohomology-section-phantom]{Cohomologie des faisceaux}
\item \hyperref[sites-cohomology-section-phantom]{Cohomologie sur les sites}
\item \hyperref[dga-section-phantom]{Algèbre différentielle graduée}
\item \hyperref[dpa-section-phantom]{Algèbre à puissances divisées}
\item \hyperref[sdga-section-phantom]{Faisceaux différentiels gradués}
\item \hyperref[hypercovering-section-phantom]{Hyperrecouvrements}
\end{enumerate}
Schémas
\begin{enumerate}
\setcounter{enumi}{25}
\item \hyperref[schemes-section-phantom]{Schémas}
\item \hyperref[constructions-section-phantom]{Constructions de schémas}
\item \hyperref[properties-section-phantom]{Propriétés des schémas}
\item \hyperref[morphisms-section-phantom]{Morphismes de schémas}
\item \hyperref[coherent-section-phantom]{Cohomologie des schémas}
\item \hyperref[divisors-section-phantom]{Diviseurs}
\item \hyperref[limits-section-phantom]{Limites de schémas}
\item \hyperref[varieties-section-phantom]{Variétés}
\item \hyperref[topologies-section-phantom]{Topologies sur les schémas}
\item \hyperref[descent-section-phantom]{Descente}
\item \hyperref[perfect-section-phantom]{Catégories dérivées de schémas}
\item \hyperref[more-morphisms-section-phantom]{Compléments sur les morphismes}
\item \hyperref[flat-section-phantom]{Compléments sur la platitude}
\item \hyperref[groupoids-section-phantom]{Schémas en groupoïdes}
\item \hyperref[more-groupoids-section-phantom]{Compléments sur les schémas en groupoïdes}
\item \hyperref[etale-section-phantom]{Morphismes étales de schémas}
\end{enumerate}
Thèmes de théorie des schémas
\begin{enumerate}
\setcounter{enumi}{41}
\item \hyperref[chow-section-phantom]{Homologie de Chow}
\item \hyperref[intersection-section-phantom]{Théorie de l'intersection}
\item \hyperref[pic-section-phantom]{Schémas de Picard des courbes}
\item \hyperref[weil-section-phantom]{Théories cohomologiques de Weil}
\item \hyperref[adequate-section-phantom]{Modules adéquats}
\item \hyperref[dualizing-section-phantom]{Complexes dualisants}
\item \hyperref[duality-section-phantom]{Dualité pour les schémas}
\item \hyperref[discriminant-section-phantom]{Discriminants et différentes}
\item \hyperref[derham-section-phantom]{Cohomologie de de Rham}
\item \hyperref[local-cohomology-section-phantom]{Cohomologie locale}
\item \hyperref[algebraization-section-phantom]{Géométrie algébrique et formelle}
\item \hyperref[curves-section-phantom]{Courbes algébriques}
\item \hyperref[resolve-section-phantom]{Résolution des surfaces}
\item \hyperref[models-section-phantom]{Réduction semi-stable}
\item \hyperref[functors-section-phantom]{Foncteurs et morphismes}
\item \hyperref[equiv-section-phantom]{Catégories dérivées de variétés}
\item \hyperref[pione-section-phantom]{Groupes fondamentaux de schémas}
\item \hyperref[etale-cohomology-section-phantom]{Cohomologie étale}
\item \hyperref[crystalline-section-phantom]{Cohomologie cristalline}
\item \hyperref[proetale-section-phantom]{Cohomologie pro-étale}
\item \hyperref[relative-cycles-section-phantom]{Cycles relatifs}
\item \hyperref[more-etale-section-phantom]{Compléments de cohomologie étale}
\item \hyperref[trace-section-phantom]{La formule des traces}
\end{enumerate}
Espaces algébriques
\begin{enumerate}
\setcounter{enumi}{64}
\item \hyperref[spaces-section-phantom]{Espaces algébriques}
\item \hyperref[spaces-properties-section-phantom]{Propriétés des espaces algébriques}
\item \hyperref[spaces-morphisms-section-phantom]{Morphismes d'espaces algébriques}
\item \hyperref[decent-spaces-section-phantom]{Espaces algébriques décents}
\item \hyperref[spaces-cohomology-section-phantom]{Cohomologie des espaces algébriques}
\item \hyperref[spaces-limits-section-phantom]{Limites d'espaces algébriques}
\item \hyperref[spaces-divisors-section-phantom]{Diviseurs sur les espaces algébriques}
\item \hyperref[spaces-over-fields-section-phantom]{Espaces algébriques sur des corps}
\item \hyperref[spaces-topologies-section-phantom]{Topologies sur les espaces algébriques}
\item \hyperref[spaces-descent-section-phantom]{Descente et espaces algébriques}
\item \hyperref[spaces-perfect-section-phantom]{Catégories dérivées d'espaces}
\item \hyperref[spaces-more-morphisms-section-phantom]{Compléments sur les morphismes d'espaces}
\item \hyperref[spaces-flat-section-phantom]{Platitude sur les espaces algébriques}
\item \hyperref[spaces-groupoids-section-phantom]{Groupoïdes dans les espaces algébriques}
\item \hyperref[spaces-more-groupoids-section-phantom]{Compléments sur les groupoïdes dans les espaces}
\item \hyperref[bootstrap-section-phantom]{Amorçage}
\item \hyperref[spaces-pushouts-section-phantom]{Sommes amalgamées d'espaces algébriques}
\end{enumerate}
Thèmes de géométrie
\begin{enumerate}
\setcounter{enumi}{81}
\item \hyperref[spaces-chow-section-phantom]{Groupes de Chow des espaces}
\item \hyperref[groupoids-quotients-section-phantom]{Quotients de groupoïdes}
\item \hyperref[spaces-more-cohomology-section-phantom]{Compléments sur la cohomologie des espaces}
\item \hyperref[spaces-simplicial-section-phantom]{Espaces simpliciaux}
\item \hyperref[spaces-duality-section-phantom]{Dualité pour les espaces}
\item \hyperref[formal-spaces-section-phantom]{Espaces algébriques formels}
\item \hyperref[restricted-section-phantom]{Algébrisation des espaces formels}
\item \hyperref[spaces-resolve-section-phantom]{Résolution des surfaces revisitée}
\end{enumerate}
Théorie des déformations
\begin{enumerate}
\setcounter{enumi}{89}
\item \hyperref[formal-defos-section-phantom]{Théorie formelle des déformations}
\item \hyperref[defos-section-phantom]{Théorie des déformations}
\item \hyperref[cotangent-section-phantom]{Le complexe cotangent}
\item \hyperref[examples-defos-section-phantom]{Problèmes de déformation}
\end{enumerate}
Champs algébriques
\begin{enumerate}
\setcounter{enumi}{93}
\item \hyperref[algebraic-section-phantom]{Champs algébriques}
\item \hyperref[examples-stacks-section-phantom]{Exemples de champs}
\item \hyperref[stacks-sheaves-section-phantom]{Faisceaux sur les champs algébriques}
\item \hyperref[criteria-section-phantom]{Critères de représentabilité}
\item \hyperref[artin-section-phantom]{Axiomes d'Artin}
\item \hyperref[quot-section-phantom]{Espaces de Quot et de Hilbert}
\item \hyperref[stacks-properties-section-phantom]{Propriétés des champs algébriques}
\item \hyperref[stacks-morphisms-section-phantom]{Morphismes de champs algébriques}
\item \hyperref[stacks-limits-section-phantom]{Limites de champs algébriques}
\item \hyperref[stacks-cohomology-section-phantom]{Cohomologie des champs algébriques}
\item \hyperref[stacks-perfect-section-phantom]{Catégories dérivées de champs}
\item \hyperref[stacks-introduction-section-phantom]{Introduction aux champs algébriques}
\item \hyperref[stacks-more-morphisms-section-phantom]{Compléments sur les morphismes de champs}
\item \hyperref[stacks-geometry-section-phantom]{La géométrie des champs}
\end{enumerate}
Thèmes de théorie des espaces de modules
\begin{enumerate}
\setcounter{enumi}{107}
\item \hyperref[moduli-section-phantom]{Champs de modules}
\item \hyperref[moduli-curves-section-phantom]{Espaces de modules de courbes}
\end{enumerate}
Divers
\begin{enumerate}
\setcounter{enumi}{109}
\item \hyperref[examples-section-phantom]{Exemples}
\item \hyperref[exercises-section-phantom]{Exercices}
\item \hyperref[guide-section-phantom]{Guide bibliographique}
\item \hyperref[desirables-section-phantom]{Desiderata}
\item \hyperref[coding-section-phantom]{Style de programmation}
\item \hyperref[obsolete-section-phantom]{Obsolète}
\item \hyperref[fdl-section-phantom]{Licence de documentation libre GNU}
\item \hyperref[index-section-phantom]{Index généré automatiquement}
\end{enumerate}
\end{multicols}


\bibliography{my}
\bibliographystyle{amsalpha}

\end{document}
